\chapter{Formal Reconstruction of the Spectral Doctrine}
\label{chap:formal-reconstruction-spectral-doctrine}

\section{Orientation: from differentiation to reconstruction}
\label{sec:formal-reconstruction-orientation}

Chapter~\Ref{chap:spectral-algebraic-geometry} constructed the spectral doctrine by
combining connective spectral contexts, stable coefficient fibres, substitution, Zariski
descent, and comprehension. Chapter~\Ref{chap:differentiation-spectral-doctrine} then
differentiated that doctrine: cotangent complexes represent first-order variation,
structured square-zero extensions realize it geometrically, finite composites generate
an infinitesimal test geometry, and endpoint adjunctions produce differential cohesion.
In parallel, each stable coefficient fibre carries the cofree cocommutative-coalgebra
exponential and its differential-linear structure.

The unresolved problem is converse. The preceding chapter extracts first-order and
infinitesimal data from a nonlinear context, but does not say when those data determine
the context from which they arose. This chapter asks when differential data cease to be
mere invariants and become reconstruction data.

The principal fixed-centre answer is comonadic. For a connective centre \(B\) over a
connective base \(R\), a complete almost finitely presented augmentation \(A\to B\) is
recovered from the anchored shifted-cotangent map
\[
    L_{B/R}[-1]\longrightarrow L_{B/A}[-1]
\]
together with the coalgebra coherence generated by structured square-zero formation. No
extra reconstruction structure is imposed: the relevant comonad is produced by the
cotangent--square-zero adjunction constructed below.

The following adjunction records that destination. Its source and target notation will
be defined only when those categories become active:
\[
\begin{tikzcd}[column sep=huge]
    \CAlg(\Sp)^{\mathrm{cn},\mathrm{afp},\wedge}_{R//B}
        \arrow[r,shift left=1.2ex,"\operatorname{cot}_{B/R}"]
        &
    \operatorname{AnchCot}^{\mathrm{aperf}}_{B/R,\geq0}
        \arrow[l,shift left=1.2ex,"\operatorname{sqz}_{B/R}"]
\end{tikzcd}
\]
For now the display is only a marker. The unrestricted adjunction will be constructed
first and then restricted to the complete theorem domain.

Two approximation coordinates lead to this theorem. \textbf{Postnikov height} resolves a
connective spectral algebra into canonical square-zero stages and recovers it by inverse
limit. \textbf{Formal order at a centre} equips an augmentation \(A\to B\) with a
canonical decreasing filtration whose completion is functorial and whose associated
graded is freely generated by the shifted relative cotangent complex. The first
coordinate measures homotopical height; the second measures order of vanishing at the
chosen centre.

The dependency diagram records the full route of the chapter. At this stage only its
upper spine is active: the two approximation mechanisms converge on complete formal
augmentations, where comonadic reconstruction is proved. The lower labels name later
exports of that theorem and need not yet be retained as definitions.
\[
\begin{tikzcd}[column sep=tiny,row sep=huge]
    & \text{differentiated spectral doctrine}
        \arrow[dl,"\text{Postnikov height}"']
        \arrow[dr,"\text{formal order}"]
    &
    \\
    \text{finite infinitesimal stages}
        \arrow[dr,"\varprojlim"']
    &&
    \text{complete adic filtration}
        \arrow[dl,"{\operatorname{Gr}\simeq\operatorname{Sym}(L[-1])}"]
    \\
    & \text{complete formal augmentation}
        \arrow[d,"\text{comonadic classification}"]
    &
    \\
    & \text{anchored shifted-cotangent coalgebra data}
        \arrow[dl,"\text{Artinian restriction}"']
        \arrow[dr,"\text{global locality}"]
    &
    \\
    \text{Schlessinger formal moduli}
        \arrow[dr,"\text{centred leaf}"']
    &&
    \text{reconstruction-local geometry}
        \arrow[dl,"\text{centred leaf}"]
    \\
    & \text{recovered tangent functor}
    &
\end{tikzcd}
\]

The argument has three phases. First we construct and control the two approximation
mechanisms, culminating in the complete fixed-centre domain and the comonadic
cotangent--square-zero theorem. Next we extract bounded Artinian tests and then remove
the centre by imposing Postnikov continuity and nilpotent excision globally. Finally we
return to a centre, stabilize split square-zero lifting, and compare the recovered
tangent functor with the two differential branches of
Chapter~\Ref{chap:differentiation-spectral-doctrine}.

That final comparison is deliberately postponed. The cofree differential-linear
structure remains in the stable coefficient fibre, whereas the square-zero branch
supplies the nonlinear infinitesimal geometry used for reconstruction. They will be
compared only after both reconstruction mechanisms are available, and only on the
represented finite-projective domain where the universal derivation appears on both
sides.

The conventions of Chapters~\Ref{chap:spectral-algebraic-geometry} and
\Ref{chap:differentiation-spectral-doctrine} remain in force. Almost finite
presentation, coherence, nilpotence, completeness, Artinianity, and reconstruction
locality are properties specifying theorem domains. Finite square-zero
factorizations, coherent refinements, Postnikov presentations, affine charts, and
centre presentations used inside proofs remain proof witnesses rather than retained
structure. Canonical filtrations, completion maps, cotangent maps, adjunction units
and counits, and localization transformations are the functorially constructed maps
used below.

Coordinate algebras remain connective. Filtered modules, cotangent complexes,
totalizations, mapping spectra, and stabilizations are formed in their corresponding
stable ambient categories.

\section{Postnikov height and finite infinitesimal reconstruction}
\label{sec:postnikov-reconstruction}

The first reconstruction coordinate is homotopical height. The preceding chapter defined
finite spectral infinitesimal substitutions as finite composites of square-zero
substitutions. The multiplicative Postnikov tower supplies canonical examples: each
transition is square-zero, every finite stage lies over the same discrete centre
\(H\pi_0(B)\), and the original connective algebra is recovered from the entire tower.

For \(n\geq1\), the common centre is visible in the comparison square
\[
\begin{tikzcd}[column sep=huge,row sep=large]
    \tau_{\leq n}^{\CAlg}B \arrow[r] \arrow[d]
        & \tau_{\leq n-1}^{\CAlg}B \arrow[d]
    \\
    H\pi_0(B) \arrow[r,equal]
        & H\pi_0(B).
\end{tikzcd}
\]
The upper arrow adds one homotopy layer, while the vertical maps keep every finite stage
over the same discrete centre. The point is stronger than analogy with the infinitesimal
site: the next theorem identifies each Postnikov transition with one of the structured
square-zero extensions already constructed.

\begin{theorem}[Postnikov stages are structured square-zero extensions]
\label{thm:postnikov-square-zero-stages}
Let \(B\) be a connective commutative ring spectrum. For every integer
\(n\geq1\), the truncation morphism
\[
    \tau_{\leq n}^{\CAlg}B
    \longrightarrow
    \tau_{\leq n-1}^{\CAlg}B
\]
is canonically a structured square-zero extension with fibre
\[
    H\pi_n(B)[n].
\]
The coefficient is regarded as a module over \(\tau_{\leq n-1}^{\CAlg}B\) through
\(\tau_{\leq n-1}^{\CAlg}B\to H\pi_0(B)\), and the extension is classified by a
canonical multiplicative Postnikov \(k\)-invariant
\[
    L_{\tau_{\leq n-1}^{\CAlg}B}
    \longrightarrow
    H\pi_n(B)[n+1].
\]
\end{theorem}

\begin{proof}
\leavevmode
\begin{enumerate}[label=\textup{(\arabic*)}]
    \item \textbf{Identify the truncation range.}
    The truncation map has the displayed fibre and lies in the bounded small
    extension range of the connective square-zero recognition theorem for
    commutative ring spectra.
    \item \textbf{Apply square-zero recognition.}
    That theorem identifies the truncation arrow with
    the structured square-zero extension classified by its multiplicative
    Postnikov \(k\)-invariant; see
    \cite[Corollary~7.4.1.28]{LurieHigherAlgebra}.
    \item \textbf{Record canonicity.}
    The classifying \(k\)-invariant is canonical
    because the Postnikov truncation functors and their transition maps are
    functorial.
\end{enumerate}
\end{proof}

\begin{corollary}[Relative Postnikov extensions]
\label{cor:relative-postnikov-square-zero}
Let \(A\to B\) be a morphism of connective commutative ring spectra. For
\(n\geq1\), the map
\[
    \tau_{\leq n}^{\CAlg}B
    \longrightarrow
    \tau_{\leq n-1}^{\CAlg}B
\]
is canonically a structured square-zero extension in commutative
\(A\)-algebras by \(H\pi_n(B)[n]\), classified by
\[
    L_{\tau_{\leq n-1}^{\CAlg}B/A}
    \longrightarrow
    H\pi_n(B)[n+1].
\]
\end{corollary}

\begin{proof}
\leavevmode
\begin{enumerate}[label=\textup{(\arabic*)}]
    \item \textbf{Lift the structural map.}
    The structural map
    \[
        A\longrightarrow\tau_{\leq n}^{\CAlg}B
    \]
    lifts the map \(A\to\tau_{\leq n-1}^{\CAlg}B\) through the absolute Postnikov
    extension of \Ref{thm:postnikov-square-zero-stages}.
    \item \textbf{Obtain the absolute nullhomotopy.}
    The absolute
    classifying derivation therefore acquires a canonical nullhomotopy after
    restriction along
    \[
        L_A\otimes_A\tau_{\leq n-1}^{\CAlg}B
        \longrightarrow
        L_{\tau_{\leq n-1}^{\CAlg}B}.
    \]
    \item \textbf{Apply cotangent transitivity.}
    Cotangent transitivity gives the cofiber sequence
    \[
    \begin{tikzcd}[column sep=large]
        L_A\otimes_A\tau_{\leq n-1}^{\CAlg}B
            \arrow[r]
            &
        L_{\tau_{\leq n-1}^{\CAlg}B}
            \arrow[r]
            &
        L_{\tau_{\leq n-1}^{\CAlg}B/A}.
    \end{tikzcd}
    \]
    \item \textbf{Factor through the relative cotangent complex.}
    The compatible absolute derivation and nullhomotopy determine a factorization
    through \(L_{\tau_{\leq n-1}^{\CAlg}B/A}\), through a contractible space of choices,
    \[
        L_{\tau_{\leq n-1}^{\CAlg}B/A}
        \longrightarrow
        H\pi_n(B)[n+1].
    \]
    \item \textbf{Classify the relative extension.}
    Structured square-zero classification in commutative \(A\)-algebras identifies
    the Postnikov transition with the extension classified by this relative map.
\end{enumerate}
\end{proof}

\begin{corollary}[Infinitesimal Postnikov reconstruction]
\label{cor:infinitesimal-postnikov-reconstruction}
Every connective commutative ring spectrum is reconstructed by the canonical
inverse tower
\[
    B
    \simeq
    \varprojlim\nolimits_n\tau_{\leq n}^{\CAlg}B,
    \qquad
    \tau_{\leq0}^{\CAlg}B\simeq H\pi_0(B),
\]
whose successive transition maps are structured square-zero extensions. For
each fixed \(n\), the composite
\[
    \tau_{\leq n}^{\CAlg}B\longrightarrow H\pi_0(B)
\]
is a finite composite of structured square-zero extensions; equivalently, its
corresponding affine morphism is a finite spectral infinitesimal substitution.
\end{corollary}

\begin{proof}
\leavevmode
\begin{enumerate}[label=\textup{(\arabic*)}]
    \item \textbf{Recover the inverse limit.}
    Postnikov completeness of connective spectra and creation of limits by the
    forgetful functor from commutative ring spectra give the inverse-limit
    equivalence.
    \item \textbf{Identify the finite composite.}
    The second assertion follows from
    \Ref{thm:postnikov-square-zero-stages} by finite composition.
\end{enumerate}
\end{proof}

\begin{remark}[Finite stages versus complete reconstruction]
\label{rem:postnikov-finite-complete}
A finite infinitesimal substitution records only membership in the
finite-composition closure of square-zero infinitesimal substitutions. On
coordinate algebras, Postnikov reconstruction is an inverse limit of the
finite structured square-zero stages. In this case the Postnikov functors
supply the canonical finite-stage factorization.
\end{remark}

\begin{remark}[Postnikov reconstruction in the differentiated doctrine]
\label{rem:postnikov-reconstruction-doctrinal}
For every \(n\), the corresponding affine morphism of the finite composite
\[
    \tau_{\leq n}^{\CAlg}B\longrightarrow H\pi_0(B)
\]
is an object of the finite infinitesimal geometry constructed in
Chapter~\Ref{chap:differentiation-spectral-doctrine}, while
\[
    B\simeq\varprojlim\nolimits_n\tau_{\leq n}^{\CAlg}B
\]
is obtained by passing from finite stages to their inverse limit.
For each \(n\), the finite stage fits into the compatible triangle
\[
\begin{tikzcd}[column sep=huge,row sep=large]
    \tau_{\leq n}^{\CAlg}B
        \arrow[r]
        \arrow[dr]
        &
    \tau_{\leq n-1}^{\CAlg}B
        \arrow[d]
    \\
        &
    H\pi_0(B).
\end{tikzcd}
\]
  The
infinitesimal site supplies the finite approximants and Postnikov completeness
supplies the reconstruction.
\end{remark}


Thus Postnikov height gives the first reconstruction prototype: finite infinitesimal
stages are canonical approximants, and the inverse limit over all heights recovers the
connective algebra. This controls homotopical height. Formal order at a fixed
augmentation is a different coordinate, constructed next.

\section{Formal order and the canonical adic filtration}
\label{sec:canonical-adic-filtration}


Postnikov reconstruction varies homotopical height over a fixed discrete centre.
Formal-order reconstruction instead fixes an augmentation
\[
    A\longrightarrow B
\]
and resolves it by order of vanishing along \(B\). The filtration is canonical: it is
created by the left adjoint to filtration-zero augmentation and then completed by a
symmetric-monoidal reflector.

The shifted cotangent complex supplies the linear generator of this filtration. The
following dependency square states the calculation the section will establish:
\[
\begin{tikzcd}[column sep=huge,row sep=large]
    \text{augmentation }(A\to B)
        \arrow[r,"\operatorname{adic}\ \text{and completion}"]
        \arrow[d,"\text{shifted cotangent}"']
    &
    \widehat{\operatorname{adic}}(A\to B)
        \arrow[d,"\operatorname{Gr}"]
    \\
    L_{B/A}[-1]
        \arrow[r,"\text{free graded commutative algebra}"']
    &
    \operatorname{Sym}^{\mathrm{gr}}_B
        \bigl(L_{B/A}[-1]\langle1\rangle\bigr).
\end{tikzcd}
\]
The square is organizational rather than internal to a single category. Its content is
that the formal neighbourhood is functorially filtered and that its graded layers are
freely generated by the same cotangent object that represented first-order variation in
Chapter~\Ref{chap:differentiation-spectral-doctrine}.

\subsection{Filtered spectral algebra and the adic left adjoint}

\begin{definition}[Filtered and graded spectral objects]
\label{def:filtered-graded-spectral-objects}
Let \(\mathbb N_{\leq}\) denote the nonnegative integers as an ordered
symmetric monoidal category under addition, and let
\(\mathbb N_{\mathrm{disc}}\) denote the same commutative monoid regarded as a
discrete symmetric monoidal category. Define
\[
    \operatorname{Fil}(\Sp)
    :=
    \Fun(\mathbb N_{\leq}^{\op},\Sp),
    \qquad
    \operatorname{Gr}(\Sp)
    :=
    \Fun(\mathbb N_{\mathrm{disc}},\Sp),
\]
with Day convolution. A filtered object is a tower
\[
    \cdots\longrightarrow F^2M\longrightarrow F^1M\longrightarrow F^0M.
\]
Its associated graded object is
\[
    \operatorname{Gr}^nM
    :=
    \operatorname{cofib}(F^{n+1}M\to F^nM),
    \qquad n\geq0,
\]
and these terms assemble into a symmetric monoidal functor
\[
    \operatorname{Gr}:
    \operatorname{Fil}(\Sp)
    \longrightarrow
    \operatorname{Gr}(\Sp).
\]
For a nonnegatively graded object \(N\) and \(w\geq0\), write
\[
    N\langle w\rangle
\]
for the upward weight shift, with
\[
    (N\langle w\rangle)_k:=N_{k-w},
\]
where negative weights are zero. In particular, if an ungraded spectrum or
module \(M\) is placed in weight zero, then \(M\langle w\rangle\) is
concentrated in weight \(w\).
For \(r<n\), write
\[
    F^rM/F^nM
    :=
    \operatorname{cofib}(F^nM\to F^rM),
\]
and put
\[
    F^\infty M:=\varprojlim\nolimits_nF^nM.
\]
The filtration is \textbf{complete} if
\[
    F^\infty M\simeq0.
\]
Equivalently, for every \(r\geq0\), the canonical morphism
\[
    F^rM
    \longrightarrow
    \varprojlim\nolimits_{n>r}F^rM/F^nM
\]
is an equivalence. The indexing by \(\mathbb N\) concerns filtration degree
only. Connectivity hypotheses will be imposed explicitly when needed; otherwise the
filtered terms range over arbitrary spectra.
\end{definition}

\begin{construction}[Completion of a filtered spectrum]
\label{con:filtered-spectrum-completion}
Let
\[
    \operatorname{const}:
    \Sp\longrightarrow\operatorname{Fil}(\Sp)
\]
denote the constant-filtration functor, so that every transition map of
\(\operatorname{const}(K)\) is the identity of \(K\). The filtration-zero
embedding of \Ref{conv:filtered-base-weight-zero} is a different functor: it
places an unfiltered object in filtration degree zero and vanishes in positive
filtration degrees.

For a filtered spectrum \(F^\bullet M\), the limit maps
\(F^\infty M\to F^rM\) assemble to a canonical morphism
\[
    \operatorname{const}(F^\infty M)
    \longrightarrow
    F^\bullet M.
\]
Define its \textbf{completion} by
\[
    \widehat{F^\bullet M}
    :=
    \operatorname{cofib}
    \left(
        \operatorname{const}(F^\infty M)
        \longrightarrow
        F^\bullet M
    \right).
\]
For every filtration degree \(r\),
\[
    \widehat F^{\,r}M
    \simeq
    \operatorname{cofib}(F^\infty M\to F^rM)
    \simeq
    \varprojlim\nolimits_{n>r}F^rM/F^nM.
\]
The multiplicative and module completion reflectors induced by this
localization are described in
\Ref{prop:filtered-completion-reflector}. Through
\Ref{conv:nonnegative-filtered-indexing}, this is the complete filtered
\(\mathbb E_\infty\)-algebra construction used in
\cite[Section~2]{AntieauFiltrationsCohomologyIII2025}.
\end{construction}

\begin{proposition}[Universal property and associated graded of completion]
\label{prop:filtered-completion-reflector}
Completion is left adjoint to the inclusion of complete filtered spectra. It
is idempotent and there are canonical equivalences
\[
    \operatorname{Gr}\widehat{F^\bullet M}
    \simeq
    \operatorname{Gr}F^\bullet M.
\]
The completion localization is compatible with Day convolution. Hence the
subcategory of complete filtered spectra inherits the canonical localized
symmetric monoidal structure for which completion is symmetric monoidal, and
completion induces the corresponding reflectors on filtered commutative
spectral algebras and on their module categories. On underlying filtered
spectra these reflectors are computed by
\Ref{con:filtered-spectrum-completion}, and associated graded is unchanged.
\end{proposition}

\begin{proof}
\leavevmode
\begin{enumerate}[label=\textup{(\arabic*)}]
    \item \textbf{Reflector property.}
    Let \(G\) be complete. The adjunction between constant filtrations and
    inverse limits gives
    \[
        \Map_{\operatorname{Fil}(\Sp)}
        \bigl(\operatorname{const}(K),G\bigr)
        \simeq
        \Map_\Sp\left(K,\varprojlim\nolimits_rG^r\right)
        \simeq *
    \]
    for every spectrum \(K\). Applying \(\Map(-,G)\) to the defining cofiber
    sequence of \Ref{con:filtered-spectrum-completion} yields
    \[
        \Map(\widehat F,G)\simeq\Map(F,G),
    \]
    so completion has the required reflector universal property.
    \item \textbf{Prove completeness.}
    Taking inverse limits in the defining cofiber sequence gives
    \[
        \varprojlim\nolimits_r\widehat F^{\,r}
        \simeq
        \operatorname{cofib}
        \bigl(F^\infty M\xrightarrow{\mathrm{id}}F^\infty M\bigr)
        \simeq0.
    \]
    Thus \(\widehat F\) is complete.
    \item \textbf{Compute associated graded.}
    The constant filtration has zero associated
    graded, and exactness of \(\operatorname{Gr}\) gives
    \[
        \operatorname{Gr}\widehat F
        \simeq
        \operatorname{Gr}F.
    \]
    \item \textbf{Local equivalences.}
    A morphism \(u:F\to G\) is inverted by completion precisely when
    \(\operatorname{Gr}(u)\) is an equivalence. For the converse direction, put
    \[
        H:=\operatorname{cofib}(u).
    \]
    If \(\operatorname{Gr}H\simeq0\), then every transition
    \(H^{n+1}\to H^n\) is an equivalence. Hence \(H\) is a constant filtration
    and \(\widehat H\simeq0\). Exactness of completion then makes
    \(\widehat u\) an equivalence.
    \item \textbf{Monoidal localization.}
    Associated graded is symmetric monoidal. Tensor products therefore preserve
    completion equivalences. The symmetric-monoidal localization theorem
    \cite[Definition~2.2.1.6 and Proposition~2.2.1.9]{LurieHigherAlgebra}
    equips the complete filtered spectra with the localized symmetric monoidal
    structure and induces the corresponding localizations on commutative algebra
    objects and module categories. Their underlying filtered spectra are
    computed by \Ref{con:filtered-spectrum-completion}.
\end{enumerate}
\end{proof}

\begin{convention}[Comparison with integer-indexed filtrations]
\label{conv:nonnegative-filtered-indexing}
When comparing with literature using decreasing \(\mathbb Z\)-indexed
filtrations, an \(\mathbb N\)-indexed filtration in this chapter is extended
canonically by
\[
    F^rM:=F^0M
    \qquad(r\leq0).
\]
This identifies our filtered objects with the full subcategory of decreasing
integer-indexed filtrations which are constant in nonpositive filtration degrees. Under
this identification, associated graded in nonnegative weights and completion
agree with the corresponding integer-indexed constructions. In particular, the comparison with
Antieau's filtered and complete filtered \(\mathbb E_\infty\)-algebras uses
this canonical extension.
\end{convention}

\begin{definition}[Weight-one filtered module]
\label{def:weight-one-filtered-module}
For a spectrum or module \(M\), let
\[
    \operatorname{wt}_1(M)
\]
denote the two-step descending filtered object
\[
\begin{tikzcd}[column sep=large]
    \cdots \arrow[r]
        & 0 \arrow[r]
        & M \arrow[r,"\mathrm{id}"]
        & M ,
\end{tikzcd}
\]
with \(M\) in filtration degrees one and zero. Equivalently,
\[
    F^0\operatorname{wt}_1(M)=M,
    \qquad
    F^1\operatorname{wt}_1(M)=M,
    \qquad
    F^n\operatorname{wt}_1(M)=0
    \quad(n\geq2).
\]
Its associated graded is concentrated in weight one:
\[
    \operatorname{Gr}\operatorname{wt}_1(M)\simeq M\langle1\rangle.
\]
\end{definition}

\begin{convention}[Unfiltered bases in filtered and graded algebra]
\label{conv:filtered-base-weight-zero}
Whenever an unfiltered commutative ring spectrum \(R\) is used as a base in
\(\operatorname{Fil}(\Sp)\) or \(\operatorname{Gr}(\Sp)\), it is embedded in
filtration degree zero in the filtered case and in graded weight zero in the graded
case. Thus in the filtered case
\[
    F^0R=R,
    \qquad
    F^nR=0\quad(n\geq1),
\]
and in the graded case the only nonzero graded term is weight zero. This
convention applies, in particular, to expressions such as
\(\CAlg(\operatorname{Fil}(\Sp))_{R/}\).
\end{convention}

\begin{proposition}[Pointwise Postnikov truncation in filtered and graded spectral algebra]
\label{prop:filtered-graded-operadic-postnikov}
Let
\[
    \mathcal C=\operatorname{Fil}(\Sp)
    \qquad\text{or}\qquad
    \mathcal C=\operatorname{Gr}(\Sp)
\]
with its pointwise Postnikov \(t\)-structure. Day convolution is compatible
with this \(t\)-structure: if \(M\in\mathcal C_{\geq a}\) and
\(N\in\mathcal C_{\geq b}\), then
\[
    M\otimes N\in\mathcal C_{\geq a+b}.
\]
Consequently the Postnikov truncation functors lift canonically to connective
commutative algebra objects. For every pointwise connective commutative
algebra \(C\) and every pointwise connective \(n\)-truncated commutative
algebra \(E\), there is a natural equivalence
\[
    \Map_{\CAlg(\mathcal C)}(C,E)
    \simeq
    \Map_{\CAlg(\mathcal C)}(\tau_{\leq n}^{\CAlg}C,E).
\]
Moreover every pointwise connective commutative algebra is recovered from its
operadic Postnikov tower:
\[
    C\simeq\varprojlim\nolimits_n\tau_{\leq n}^{\CAlg}C.
\]
\end{proposition}

\begin{proof}
\leavevmode
\begin{enumerate}[label=\textup{(\arabic*)}]
    \item \textbf{Verify graded connectivity.}
    In the graded case the tensor-product connectivity estimate is termwise.
    \item \textbf{Verify filtered connectivity.}
    In
    the filtered case each filtration term of the Day convolution is a colimit of
    tensor products of filtration terms, so the same estimate follows from the
    connective tensor-product theorem in spectra.
    \item \textbf{Lift Postnikov truncations.}
    The general compatibility
    result for a monoidal \(t\)-structure therefore supplies Postnikov truncations
    on connective commutative algebra objects; see \cite{LurieHigherAlgebra}.
    \item \textbf{Obtain the mapping equivalence.}
    Its universal property gives the displayed mapping equivalence.
    \item \textbf{Recover the operadic tower.}
    Finally the
    forgetful functor from commutative algebra objects creates limits, while the
    pointwise Postnikov \(t\)-structure is left complete, so the operadic tower
    has limit \(C\).
\end{enumerate}
\end{proof}

Fix a commutative ring spectrum \(R\). For a filtered commutative \(R\)-algebra
\(F^\bullet C\), filtration degree zero and the weight-zero associated graded determine
the augmentation
\[
    F^0C\longrightarrow\operatorname{Gr}^0C.
\]

\begin{proposition}[Augmentation from filtered algebras]
\label{prop:filtered-augmentation-functor}
The assignment
\[
    \operatorname{aug}:
    \CAlg(\operatorname{Fil}(\Sp))_{R/}
    \longrightarrow
    \Fun(\Delta^1,\CAlg(\Sp)_{R/})
\]
which sends a filtered algebra to
\(F^0C\to\operatorname{Gr}^0C\) is accessible and preserves limits. The filtration-zero augmentation is represented on underlying spectra by the
cofiber square
\[
\begin{tikzcd}[column sep=large,row sep=large]
    F^1C \arrow[r] \arrow[d]
        & F^0C \arrow[d]
    \\
    0 \arrow[r]
        & \operatorname{Gr}^0C .
\end{tikzcd}
\]
Hence it admits a left adjoint
\[
    \operatorname{adic}
    \dashv
    \operatorname{aug}.
\]
\end{proposition}

\begin{proof}
\leavevmode
\begin{enumerate}[label=\textup{(\arabic*)}]
    \item \textbf{Evaluate the filtration endpoints.}
    Evaluation in filtration degrees zero and one is accessible and preserves all
    limits.
    \item \textbf{Identify the weight-zero graded term.}
    In the stable category of spectra,
    \[
        \operatorname{Gr}^0C
        =
        \operatorname{cofib}(F^1C\to F^0C)
    \]
    is also a finite-limit construction.
    \item \textbf{Preserve arbitrary limits.}
    Finite limits commute with arbitrary
    limits, and limits of commutative algebra objects are created in the underlying
    symmetric monoidal category. Thus both endpoints and their canonical
    comparison preserve limits.
    \item \textbf{Apply the adjoint functor theorem.}
    The source and target are presentable, so the
    presentable adjoint functor theorem supplies the left adjoint.
\end{enumerate}
\end{proof}

\begin{lemma}[The module-level augmentation filtration]
\label{lem:module-level-augmentation-filtration}
Let
\[
    \operatorname{aug}_{\mathrm{mod}}:
    \operatorname{Fil}(\Sp)
    \longrightarrow
    \Fun(\Delta^1,\Sp)
\]
be the functor
\[
    F^\bullet M\longmapsto
    \bigl(F^0M\longrightarrow\operatorname{Gr}^0M\bigr).
\]
It admits a left adjoint. On an arrow \(V\to V_0\) this left adjoint is the
two-step filtered spectrum
\[
    \cdots\longrightarrow0\longrightarrow
    \operatorname{fib}(V\to V_0)
    \longrightarrow V,
\]
with \(V\) in filtration degree zero and the fibre in filtration degree one. In particular,
its filtration-zero term is \(V\) and its weight-zero associated graded is
\(V_0\).
\end{lemma}

\begin{proof}
\leavevmode
\begin{enumerate}[label=\textup{(\arabic*)}]
    \item \textbf{Identify morphisms of fibre sequences.}
    For a filtered spectrum \(F^\bullet N\), a morphism from the displayed
    two-step object to \(F^\bullet N\) is the same as a morphism of the two fibre
    sequences
    \[
    \begin{tikzcd}[column sep=large]
        \operatorname{fib}(V\to V_0) \arrow[r] \arrow[d]
            & V \arrow[r] \arrow[d]
            & V_0 \arrow[d]
        \\
        F^1N \arrow[r]
            & F^0N \arrow[r]
            & \operatorname{Gr}^0N.
    \end{tikzcd}
    \]
    \item \textbf{Compute the mapping space.}
    In a stable \(\infty\)-category, the space of such morphisms of fibre sequences
    is canonically the mapping space from the arrow \(V\to V_0\) to
    \(F^0N\to\operatorname{Gr}^0N\).
    \item \textbf{Conclude the adjunction.}
    This is the required adjunction.
    \item \textbf{Identify the filtration terms.}
    The final
    assertion is the defining fibre sequence.
\end{enumerate}
\end{proof}

\begin{proposition}[The adic construction recovers the input augmentation]
\label{prop:adic-recovers-augmentation}
For every augmented arrow
\[
    A\longrightarrow B
\]
of commutative \(R\)-algebras, the filtered algebra
\(\operatorname{adic}(A\to B)\) has canonical identifications
\[
    F^0\operatorname{adic}(A\to B)\simeq A,
    \qquad
    \operatorname{Gr}^0\operatorname{adic}(A\to B)\simeq B,
\]
under which the canonical augmentation
\[
    F^0\operatorname{adic}(A\to B)
    \longrightarrow
    \operatorname{Gr}^0\operatorname{adic}(A\to B)
\]
is the original morphism \(A\to B\).

On underlying spectra, its first filtration quotient is displayed by
\[
\begin{tikzcd}[column sep=large,row sep=large]
    F^1\operatorname{adic}(A\to B) \arrow[r] \arrow[d]
        & A \arrow[d]
    \\
    0 \arrow[r]
        & B .
\end{tikzcd}
\]
\end{proposition}

\begin{proof}
\leavevmode
\begin{enumerate}[label=\textup{(\arabic*)}]
    \item \textbf{Multiplicative left adjoint.}
    The algebraic left adjoint is the multiplicative extension of
    \Ref{lem:module-level-augmentation-filtration}.
    \item \textbf{Free augmented arrows.}
    Resolve the augmented arrow \(A\to B\) by the simplicial resolution associated with
    free augmented commutative \(R\)-algebra arrows. On a free arrow
    \[
        \operatorname{Sym}_R(V)
        \longrightarrow
        \operatorname{Sym}_R(V_0)
    \]
    induced by a morphism of \(R\)-modules \(V\to V_0\), the adjunction between
    free commutative algebras and underlying modules, together with
    \Ref{lem:module-level-augmentation-filtration}, identifies its adic image with
    the free filtered commutative algebra generated by the two-step filtration
    \[
        \operatorname{fib}(V\to V_0)\longrightarrow V.
    \]
    Consequently its filtration-zero algebra is \(\operatorname{Sym}_R(V)\), while
    its weight-zero associated graded algebra is \(\operatorname{Sym}_R(V_0)\), and
    the comparison between them is the original free arrow.
    \item \textbf{Passage to geometric realization.}
    The commutative-algebra monad preserves sifted colimits. Hence this
    free resolution is recovered by geometric realization, and evaluation at
    filtration zero, associated graded in weight zero, and the filtered free-algebra
    construction preserve the sifted colimits occurring in that realization.
    Passing to the realization therefore gives the two asserted identifications for
    an arbitrary augmented arrow. Their comparison is the adjunction unit and is
    identified with \(A\to B\).
\end{enumerate}
\end{proof}

\begin{definition}[Canonical adic filtration and completion]
\label{def:canonical-adic-filtration}
For an augmented arrow
\[
    A\longrightarrow B
\]
of commutative \(R\)-algebras, define
\[
    \operatorname{adic}(A\to B)
\]
to be its image under the left adjoint of
\Ref{prop:filtered-augmentation-functor}. Its \textbf{completed canonical adic
filtration} is the multiplicative completion
\[
    \widehat{\operatorname{adic}}(A\to B)
    :=
    \widehat{\bigl(\operatorname{adic}(A\to B)\bigr)}
\]
of \Ref{con:filtered-spectrum-completion}. On underlying filtered spectra,
for every \(r\geq0\),
\[
    F^r\widehat{\operatorname{adic}}(A\to B)
    \simeq
    \varprojlim\nolimits_{n>r}
    \operatorname{cofib}
    \left(
        F^n\operatorname{adic}(A\to B)
        \longrightarrow
        F^r\operatorname{adic}(A\to B)
    \right).
\]
Write
\[
    A^{\wedge}_B
    :=
    F^0\widehat{\operatorname{adic}}(A\to B).
\]
The adjunction unit of \Ref{prop:adic-recovers-augmentation} followed by the
completion unit gives a canonical morphism
\[
    A\longrightarrow A^{\wedge}_B.
\]
\end{definition}

\subsection{Cotangent linearization of formal order}

\begin{definition}[Shifted cotangent complex of an augmentation]
\label{def:shifted-cotangent-augmentation}
For an augmented arrow \(A\to B\) of commutative \(R\)-algebras, its shifted
cotangent complex is
\[
    L_{B/A}[-1].
\]

If \(C\) is a filtered commutative algebra, regard
\(\operatorname{Gr}^0C\) as concentrated in filtration degree zero. Its shifted
cotangent complex relative to the filtered algebra \(C\) is
\[
    L_{\operatorname{Gr}^0C/C}[-1]
\]
in the filtered module category.

If \(C\) is a nonnegatively graded commutative algebra, write
\[
    C^{\mathrm{wt}=0}
\]
for its weight-zero algebra, regarded as concentrated in weight zero. Its
shifted cotangent complex relative to \(C\) is
\[
    L_{C^{\mathrm{wt}=0}/C}[-1]
\]
in the graded module category.
\end{definition}

\begin{convention}[Almost-perfect filtered modules]
\label{conv:almost-perfect-filtered-modules}
If \(F^\bullet A\) is a pointwise connective filtered commutative ring
spectrum, an
\(F^\bullet A\)-module is called \textbf{almost perfect} when it is almost
perfect in the stable module category
\(\Mod_{F^\bullet A}(\operatorname{Fil}(\Sp))\), equipped with the pointwise
Postnikov \(t\)-structure. Perfectness and connectivity in filtered module
categories are interpreted in the same intrinsic way.
\end{convention}


\begin{convention}[Almost finite presentation in filtered and graded spectral algebra]
\label{conv:filtered-graded-almost-finite-presentation}
Let
\[
    \mathcal C=\operatorname{Fil}(\Sp)
    \qquad\text{or}\qquad
    \mathcal C=\operatorname{Gr}(\Sp)
\]
with its pointwise Postnikov \(t\)-structure, and let \(A\to C\) be a
morphism of pointwise connective commutative algebra objects. We say that
\(C\) is \textbf{almost finitely presented over \(A\)} if, for every integer
\(n\geq0\), the functor represented by \(C\) preserves filtered colimits on
the full subcategory of pointwise connective \(n\)-truncated commutative
\(A\)-algebras. This condition is a property of the morphism.
\end{convention}

\begin{remark}[Relation with formal integration]
\label{rem:bmn-formal-integration-comparison}
The filtered constructions in this section parallel the affine formal-integration
machinery of Brantner--Magidson--Nuiten. The shifted cotangent complex used
here is \(L_{B/A}[-1]\); the adic filtration is left adjoint to filtration-zero
augmentation; and the category of complete almost finitely presented augmentations is the domain of the
shifted-cotangent--structured-square-zero comonadicity theorem. The main comparison results used
below are
\cite[
    Remark~2.10,
    Lemma~2.12,
    Proposition~2.14,
    Lemma~2.16,
    Proposition~2.18,
    and Theorem~2.25
]{BrantnerMagidsonNuiten}.
The proofs are included here because the filtered connectivity and
almost-perfectness estimates are also used in the later reconstruction and
Artinian arguments.
\end{remark}


\begin{lemma}[Cotangent compatibility with grading, filtration zero, and adic formation]
\label{lem:cotangent-graded-adic-compatibility}
Let \(F^\bullet C\) be a filtered commutative algebra. There are canonical
natural equivalences
\[
    \operatorname{Gr}
    \bigl(
        L_{\operatorname{Gr}^0C/F^\bullet C}[-1]
    \bigr)
    \simeq
    L_{\operatorname{Gr}^0C/\operatorname{Gr}C}[-1],
\]
\[
    F^0
    \bigl(
        L_{\operatorname{Gr}^0C/F^\bullet C}[-1]
    \bigr)
    \simeq
    L_{\operatorname{Gr}^0C/F^0C}[-1],
\]
and, for every augmented arrow \(A\to B\),
\[
    L_{B/\operatorname{adic}(A\to B)}[-1]
    \simeq
    \operatorname{wt}_1\bigl(L_{B/A}[-1]\bigr).
\]
The first equivalence is taken in graded modules over
\(\operatorname{Gr}C\), and the second in
\(\operatorname{Gr}^0C\)-modules after evaluating the shifted relative
cotangent complex in filtration zero.

The two comparison functors are summarized by
\[
\begin{tikzcd}[column sep=huge,row sep=large]
    \operatorname{Gr}
    \bigl(L_{\operatorname{Gr}^0C/F^\bullet C}[-1]\bigr)
        \arrow[r,"\simeq"]
        &
    L_{\operatorname{Gr}^0C/\operatorname{Gr}C}[-1]
    \\
    F^0\bigl(L_{\operatorname{Gr}^0C/F^\bullet C}[-1]\bigr)
        \arrow[d,"\simeq"']
        &
    \\
    L_{\operatorname{Gr}^0C/F^0C}[-1]
        &
\end{tikzcd}
\]
where the horizontal and vertical arrows denote the canonical base-change
comparisons.
\end{lemma}

\begin{proof}
\leavevmode
\begin{enumerate}[label=\textup{(\arabic*)}]
    \item \textbf{Functoriality of cotangent complexes.}
    Associated graded and evaluation in filtration zero are symmetric monoidal,
    preserve colimits, and commute with free commutative algebra formation.
    Applying the universal derivation characterization of relative cotangent
    complexes therefore gives the first two equivalences.
    \item \textbf{Relative shifted-cotangent arrow.}
    For the third, put
    \[
        M:=L_{B/A}[-1].
    \]
    Consider the morphism of augmented arrows
    \[
        (A\to B)\longrightarrow(B\to B).
    \]
    Objectwise stabilization of the arrow category, together with the tangent
    equivalence of \Ref{thm:tangent-category-calg}, identifies its relative
    shifted-cotangent object after restriction to the fixed target \(B\) with the
    \(B\)-module arrow
    \[
        M\longrightarrow0.
    \]
    Indeed, a morphism from \(M\to0\) to a module arrow
    \(N_0\to N_1\) is a morphism \(M\to N_0\) together with a specified
    nullhomotopy of its composite with \(N_0\to N_1\); by cotangent transitivity
    this is exactly the relative shifted-cotangent datum for deforming \(A\to B\) while
    the target \(B\) is fixed. This is the arrow-category form of the ordinary
    relative cotangent classifier; compare
    \cite[Definition~2.9 and Remark~2.10]{BrantnerMagidsonNuiten}.
    \item \textbf{Module-level augmentation adjunction.}
    Work in filtered \(B\)-modules. The module-level augmentation functor
    \[
        \operatorname{aug}_{\mathrm{mod}}:
        \operatorname{Fil}(\Mod_B)
        \longrightarrow
        \Fun(\Delta^1,\Mod_B),
        \qquad
        F^\bullet N\longmapsto
        \bigl(F^0N\to\operatorname{Gr}^0N\bigr)
    \]
    has the left adjoint computed in
    \Ref{lem:module-level-augmentation-filtration}; in particular the arrow
    \(M\to0\) is sent to the two-step filtered module \(\operatorname{wt}_1(M)\).
    \item \textbf{Tangent functor of augmentation.}
    The algebraic augmentation functor of
    \Ref{prop:filtered-augmentation-functor} preserves split square-zero
    algebras. Indeed its two endpoint functors \(F^0\) and
    \(\operatorname{Gr}^0\) are symmetric monoidal and carry a filtered split
    square-zero algebra to the corresponding split square-zero arrow. Hence the
    functor induced by \(\operatorname{aug}\) on tangent module categories is
    exactly \(\operatorname{aug}_{\mathrm{mod}}\).
    \item \textbf{Yoneda identification.}
    Now let \(N\) be an arbitrary filtered \(B\)-module. Cotangent
    classification and the adjunction \(\operatorname{adic}\dashv
    \operatorname{aug}\) identify maps from the shifted cotangent complex of
    \(\operatorname{adic}(A\to B)\) with maps from the relative shifted-cotangent arrow
    \(M\to0\) computed above into \(\operatorname{aug}_{\mathrm{mod}}N\).
    Therefore
    \[
    \begin{aligned}
        \Map_{\operatorname{Fil}(\Mod_B)}
        \bigl(L_{B/\operatorname{adic}(A\to B)}[-1],N\bigr)
        &\simeq
        \Map_{\Fun(\Delta^1,\Mod_B)}
        \bigl(M\to0,\operatorname{aug}_{\mathrm{mod}}N\bigr)
        \\
        &\simeq
        \Map_{\operatorname{Fil}(\Mod_B)}
        \bigl(\operatorname{wt}_1(M),N\bigr).
    \end{aligned}
    \]
    The equivalence is natural in every filtered \(B\)-module \(N\), so Yoneda in
    the full filtered module category gives
    \[
        L_{B/\operatorname{adic}(A\to B)}[-1]
        \simeq
        \operatorname{wt}_1(M).
    \]
    The Yoneda argument takes place in the full filtered module category.
\end{enumerate}
\end{proof}


\begin{theorem}[Associated graded of the canonical adic filtration]
\label{thm:adic-associated-graded}
For every augmented arrow \(A\to B\) of commutative ring spectra there is a
canonical equivalence of graded commutative \(B\)-algebras
\[
    \operatorname{Gr}
    \widehat{\operatorname{adic}}(A\to B)
    \simeq
    \operatorname{Sym}^{\mathrm{gr}}_B
    \bigl(L_{B/A}[-1]\langle1\rangle\bigr).
\]
Completion preserves associated graded.
\end{theorem}

\begin{proof}
\leavevmode
\begin{enumerate}[label=\textup{(\arabic*)}]
    \item \textbf{Universal property.}
    Fix the source algebra \(A\). By
    \Ref{prop:adic-recovers-augmentation},
    \(\operatorname{adic}(A\to B)\) has filtration-zero algebra \(A\). Hence it
    is canonically a filtered commutative \(A\)-algebra. For a filtered
    commutative \(A\)-algebra \(F^\bullet C\), the adjunction
    \(\operatorname{adic}\dashv\operatorname{aug}\) gives
    \[
    \begin{aligned}
        &\Map_{\CAlg(\operatorname{Fil}(\Sp))_{A/}}
        \bigl(\operatorname{adic}(A\to B),F^\bullet C\bigr)
        \\
        &\qquad\simeq
        \Map_{\CAlg(\Sp)_{A/}}
        \bigl(B,\operatorname{Gr}^0C\bigr).
    \end{aligned}
    \]
    The source map from \(A\) is already fixed by the \(A\)-algebra structure.
    This is the universal property of the adic construction in
    \cite[Section~2]{BrantnerMagidsonNuiten}; through
    \Ref{conv:nonnegative-filtered-indexing}, it also agrees with the relative
    \(\mathbb E_\infty\)-Hodge filtration of
    \cite[Section~2]{AntieauFiltrationsCohomologyIII2025}.
    \item \textbf{Associated-graded calculation.}
    The associated-graded calculation
    \cite[Lemma~2.12]{BrantnerMagidsonNuiten} gives
    \[
        \operatorname{Gr}\operatorname{adic}(A\to B)
        \simeq
        \operatorname{Sym}^{\mathrm{gr}}_B
        \bigl(L_{B/A}[-1]\langle1\rangle\bigr).
    \]
    Equivalently, this is the relative Hodge calculation of
    \cite[Lemma~2.2]{AntieauFiltrationsCohomologyIII2025}. Finally,
    \Ref{prop:filtered-completion-reflector} identifies the associated graded
    before and after completion:
    \[
    \begin{tikzcd}[column sep=huge]
        \operatorname{Gr}\operatorname{adic}(A\to B)
            \arrow[r,"\sim"]
            &
        \operatorname{Gr}\widehat{\operatorname{adic}}(A\to B).
    \end{tikzcd}
    \]
    Combining the two equivalences proves the theorem.
\end{enumerate}
\end{proof}

\begin{remark}[Interpretation of the associated graded]
\label{rem:adic-associated-graded-interpretation}
The theorem identifies the associated graded of the canonical adic
filtration with the free graded commutative algebra on the shifted cotangent
complex. The completed augmented algebra is then obtained by the completion
reflector of \Ref{prop:filtered-completion-reflector}.
\end{remark}

The associated-graded theorem identifies the linear layers of formal order. Two
detection results are still needed before those layers can reconstruct an augmentation:
the canonical adic filtration must remain connective on the relevant domain, and an
associated-graded equivalence between complete filtered objects must already be an
equivalence before passing to graded pieces.

\subsection{Completion and associated-graded detection}


\begin{lemma}[Connectivity of the canonical adic filtration]
\label{lem:adic-filtration-connective}
Let \(A\to B\) be a morphism of connective commutative ring spectra which is
surjective on \(\pi_0\). Then every filtration term of
\(\operatorname{adic}(A\to B)\) is connective.
\end{lemma}

\begin{proof}
\leavevmode
\begin{enumerate}[label=\textup{(\arabic*)}]
    \item \textbf{Monadic resolution.}
    Connectivity is independent of the auxiliary base \(R\), so we may work in
    the absolute arrow category. Resolve the arrow \(A\to B\) by the
    monadic simplicial resolution for the objectwise free--forgetful adjunction
    \[
        \Fun(\Delta^1,\Sp)
        \rightleftarrows
        \Fun(\Delta^1,\CAlg(\Sp)).
    \]
    Its simplicial terms are free algebra arrows
    \[
        \operatorname{Sym}(V_q)
        \longrightarrow
        \operatorname{Sym}(V_{0,q}).
    \]
    The underlying arrow \(A\to B\) consists of connective spectra and is
    surjective on \(\pi_0\), equivalently its fibre is connective. The free
    commutative-algebra monad preserves this property. Indeed, if
    \(V\to V_0\) is a morphism of connective spectra surjective on \(\pi_0\),
    then
    \[
        \pi_0\operatorname{Sym}(V)
        \simeq
        \operatorname{Sym}_{\mathbb Z}(\pi_0V)
        \longrightarrow
        \operatorname{Sym}_{\mathbb Z}(\pi_0V_0)
        \simeq
        \pi_0\operatorname{Sym}(V_0)
    \]
    is surjective. Iterating the monad shows that every
    \(V_q\to V_{0,q}\) has connective source, target, and fibre.
    \item \textbf{Connectivity of free filtered terms.}
    By \Ref{lem:module-level-augmentation-filtration}, the module-level left
    adjoint sends \(V_q\to V_{0,q}\) to the pointwise connective two-step
    filtration
    \[
        \cdots\longrightarrow0\longrightarrow
        \operatorname{fib}(V_q\to V_{0,q})
        \longrightarrow V_q.
    \]
    For a free algebra arrow, the two free--forgetful adjunctions give a canonical
    equivalence
    \[
    \begin{aligned}
        &\operatorname{adic}
        \bigl(\operatorname{Sym}(V_q)\to\operatorname{Sym}(V_{0,q})\bigr)
        \\
        &\qquad\simeq
        \operatorname{Sym}^{\operatorname{Fil}}
        \bigl(\operatorname{adic}_{\mathrm{mod}}(V_q\to V_{0,q})\bigr).
    \end{aligned}
    \]
    Indeed both sides represent, against a filtered commutative algebra
    \(F^\bullet C\), maps from the arrow \(V_q\to V_{0,q}\) to
    \(F^0C\to\operatorname{Gr}^0C\) after passage to underlying spectra.
    \item \textbf{Geometric realization.}
    Day convolution of pointwise connective filtered spectra is pointwise
    connective, since each filtration term of a tensor product is a colimit of
    tensor products of connective spectra. The free filtered commutative algebra
    on the displayed two-step filtration is therefore pointwise connective.
    Finally \(\operatorname{adic}\), being a left adjoint, preserves the geometric
    realization of the free simplicial resolution. Evaluation in every
    filtration degree preserves the sifted colimit underlying this realization,
    and geometric realizations of simplicial connective spectra are connective.
    Hence
    \[
        F^n\operatorname{adic}(A\to B)\in\Sp_{\geq0}
    \]
    for every \(n\geq0\).
\end{enumerate}
\end{proof}


\begin{remark}[Relation with the degree-zero augmentation ideal]
\label{rem:adic-degree-zero-ideal-comparison}
For a connective augmentation surjective on \(\pi_0\), the canonical adic
filtration records the full spectral augmentation, including its higher
homotopy. The degree-zero ideal
\[
    I=\ker(\pi_0A\to\pi_0B)
\]
enters after completion: on the almost finite presentation domain,
\Ref{thm:adic-adams-derived-completion} identifies canonical adic completion with
derived \(I\)-completion. Thus the ideal-theoretic description concerns the completed object, while the
individual filtration terms retain the full spectral augmentation data.
\end{remark}

\begin{lemma}[Complete filtered detection]
\label{lem:complete-filtered-detection}
Let \(\mathcal C\) be a stable \(\infty\)-category admitting countable inverse
limits, and let
\[
    F^\bullet M\longrightarrow F^\bullet N
\]
be a morphism of complete descending filtered objects in \(\mathcal C\). If
\[
    \operatorname{Gr}M
    \xrightarrow{\simeq}
    \operatorname{Gr}N,
\]
then the filtered morphism is an equivalence. Equivalently,
\[
    F^rM\xrightarrow{\simeq}F^rN
\]
for every filtration degree \(r\geq0\).
\end{lemma}

\begin{proof}
\leavevmode
\begin{enumerate}[label=\textup{(\arabic*)}]
    \item \textbf{Finite filtration quotients.}
    Fix \(r\geq0\). The shifted tails
    \[
        F^{r+\bullet}M
        \qquad\text{and}\qquad
        F^{r+\bullet}N
    \]
    are complete. The associated-graded hypothesis identifies every successive
    cofiber
    \[
        \operatorname{cofib}
        \bigl(F^{k+1}M\to F^kM\bigr)
        \xrightarrow{\simeq}
        \operatorname{cofib}
        \bigl(F^{k+1}N\to F^kN\bigr).
    \]
    Induction on \(n\geq1\), using the cofiber sequences for the finite filtration
    quotients, gives
    \[
        F^rM/F^{r+n}M
        \xrightarrow{\simeq}
        F^rN/F^{r+n}N.
    \]
    \item \textbf{Passage to the inverse limit.}
    Completeness identifies
    \[
    \begin{tikzcd}[column sep=huge]
        F^rM
            \arrow[r,"\sim"]
            \arrow[d]
            &
        \displaystyle\varprojlim\nolimits_{n>0}F^rM/F^{r+n}M
            \arrow[d,"\sim"]
        \\
        F^rN
            \arrow[r,"\sim"']
            &
        \displaystyle\varprojlim\nolimits_{n>0}F^rN/F^{r+n}N .
    \end{tikzcd}
    \]
    The right vertical map is an equivalence because countable inverse limits
    preserve equivalences. Hence \(F^rM\to F^rN\) is an equivalence. Varying
    \(r\) proves the result.
\end{enumerate}
\end{proof}

\begin{lemma}[Connectivity from a complete filtration with connective associated graded]
\label{lem:complete-positive-connectivity}
Let \(F^\bullet C\) be a complete descending filtered spectrum such that every
associated-graded term is connective. Then \(F^0C\) is connective.
\end{lemma}

\begin{proof}
\leavevmode
\begin{enumerate}[label=\textup{(\arabic*)}]
    \item \textbf{Define the finite quotients.}
    Put \(C_n:=F^0C/F^{n+1}C\).
    \item \textbf{Establish connectivity of the quotient tower.}
    Induction on \(n\) shows that every \(C_n\) is
    connective and that the transition maps are surjective on \(\pi_0\).
    \item \textbf{Use completeness.}
    By
    completeness,
    \[
        F^0C\simeq\varprojlim\nolimits_nC_n.
    \]
    \item \textbf{Apply the Milnor exact sequence.}
    The Milnor exact sequence kills \(\pi_{-1}\) because the \(\lim^1\) term
    vanishes for the surjective tower, and all lower negative homotopy groups
    vanish termwise.
\end{enumerate}
\end{proof}


We can now retain a compact interface from this section: every augmentation has a
canonical filtered formal neighbourhood, its associated graded is cotangent-generated,
and complete filtered objects are detected on their graded layers. What remains is to
show that the finiteness and connectivity needed for comonadic reconstruction survive
through this formal-order calculus.

\section{Finiteness and convergence for fixed-centre reconstruction}
\label{sec:formal-reconstruction-finiteness}

The next task is therefore convergence rather than construction. The fixed-centre
theorem will use augmentations \(A\to B\) that are connective, finite enough for
cotangent theory, and complete for the canonical formal-order filtration. This section
establishes the estimates that make those conditions stable under the constructions used
later.

Four mechanisms provide the control. Almost-perfectness governs linear coefficients,
almost finite presentation governs algebra maps, finite-cell approximation transports
those finiteness properties through filtered constructions, and relative Hurewicz
comparison detects the first unresolved nonlinear homotopy in cotangent terms. Their
dependency is summarized by
\[
\begin{tikzcd}[column sep=tiny,row sep=large]
    & \text{almost-perfectness} 
        \arrow[dr,"\text{finite-cell approximation}"]
        \arrow[dl,"\pi_{\min}\ \text{finite}"']
    & 
    \\
    \text{first unresolved homotopy} 
        \arrow[dr,"\text{relative Hurewicz}"']
    &   
    & \text{filtered approximants}
        \arrow[dl,"\operatorname{Gr}\ \text{detection}"]
    \\
    & \text{cotangent control} 
        \arrow[d,"\text{completion}"]
    &  
    \\
    & \text{complete cotangent reconstruction data} &
\end{tikzcd}
\]


The lower-right object is the only output that must remain active: the completed adic
filtration and its shifted cotangent complex stay finite, connective, and complete
enough for Barr--Beck. The subsections establish the four links in the order in which
they are used.

\subsection{Finiteness notions and first unresolved homotopy}

\begin{definition}[Almost-perfect module]
\label{def:almost-perfect-module}
Let \(A\) be connective. An \(A\)-module \(M\) is \textbf{almost perfect} if
it is bounded below and, for every integer \(n\), its truncation
\(\tau_{\leq n}M\) is compact in the \(\infty\)-category of
\(n\)-truncated \(A\)-modules. Write
\[
    \operatorname{APerf}(A)\subseteq\Mod_A
\]
for the full stable subcategory of almost-perfect modules.
\end{definition}

\begin{lemma}[Restriction transitivity for almost-perfect modules]
\label{lem:almost-perfect-restriction-transitivity}
Let \(A\to B\) be a morphism of connective commutative ring spectra. If
\(B\) is almost perfect as an \(A\)-module, then restriction of scalars
carries almost-perfect \(B\)-modules to almost-perfect \(A\)-modules:
\[
    \operatorname{Res}^B_A
    \bigl(\operatorname{APerf}(B)\bigr)
    \subseteq
    \operatorname{APerf}(A).
\]
\end{lemma}

\begin{proof}
\leavevmode
\begin{enumerate}[label=\textup{(\arabic*)}]
    \item \textbf{Reduce to the connective case.}
    After a suspension we may take the \(B\)-module \(M\) to be connective.
    \item \textbf{Resolve by finite free modules.}
    By the connective almost-perfect resolution theorem, \(M\) is the geometric
    realization of a simplicial \(B\)-module whose terms are finite free
    \(B\)-modules; see
    \cite[Proposition~7.2.4.11(5)]{LurieHigherAlgebra}.
    \item \textbf{Restrict the simplicial terms.}
    After restriction to \(A\),
    each simplicial term is a finite direct sum of the almost-perfect
    \(A\)-module \(B\).
    \item \textbf{Use closure under realization.}
    The connective almost-perfect \(A\)-modules are closed under these geometric
    realizations by \cite[Proposition~7.2.4.11(4)]{LurieHigherAlgebra}. Hence the realization \(M\), regarded as an
    \(A\)-module, is almost perfect.
    \item \textbf{Undo the suspension.}
    Undoing the suspension gives the claim.
\end{enumerate}
\end{proof}

\begin{lemma}[First unresolved homotopy of an almost-perfect module]
\label{lem:first-unresolved-homotopy-almost-perfect}
Let \(A\) be connective and let \(M\in\operatorname{APerf}(A)\). If
\(M\) is \(n\)-connective, then its first possible nonzero homotopy module
\[
    \pi_n(M)
\]
is finitely presented over \(\pi_0(A)\). The conclusion holds for every
connective commutative ring spectrum \(A\).
\end{lemma}

\begin{proof}
\leavevmode
\begin{enumerate}[label=\textup{(\arabic*)}]
    \item \textbf{Identify the first truncation.}
    Because \(M\) is \(n\)-connective,
    \[
        \tau_{\leq n}M
        \simeq
        H\pi_n(M)[n].
    \]
    \item \textbf{Use compactness.}
    Almost-perfectness says that this truncation is compact in the
    \(\infty\)-category of \(n\)-truncated \(A\)-modules.
    \item \textbf{Test on Eilenberg--Mac Lane diagrams.}
    Test compactness only on
    filtered diagrams of Eilenberg--Mac Lane modules \(HN_\alpha[n]\). Taking
    connected components of the resulting mapping-space comparison gives
    \[
        \operatorname{Hom}_{\pi_0(A)}
        \left(\pi_n(M),\operatorname*{colim}_\alpha N_\alpha\right)
        \simeq
        \operatorname*{colim}_\alpha
        \operatorname{Hom}_{\pi_0(A)}(\pi_n(M),N_\alpha).
    \]
    \item \textbf{Read off finite presentation.}
    Thus \(\pi_n(M)\) is finitely presented over \(\pi_0(A)\).
    \item \textbf{Record the filtered and graded variants.}
    The same argument
    applies verbatim to modules over connective graded or filtered commutative ring
    spectra with their pointwise Postnikov \(t\)-structures.
\end{enumerate}
\end{proof}

\begin{definition}[Almost finite presentation]
\label{def:almost-finite-presentation}
Let \(A\to C\) be a morphism of connective commutative ring spectra. We say
that \(C\) is \textbf{almost finitely presented over \(A\)} if, for every
integer \(n\geq0\), the functor represented by \(C\) on \(n\)-truncated
connective commutative \(A\)-algebras preserves filtered colimits.
\end{definition}

\begin{definition}[Coherent commutative ring spectrum]
\label{def:coherent-commutative-ring-spectrum}
A connective commutative ring spectrum \(B\) is \textbf{coherent} if
\(\pi_0(B)\) is a coherent ordinary ring and every \(\pi_n(B)\) is finitely
presented over \(\pi_0(B)\).
\end{definition}

\begin{definition}[Coherent module]
\label{def:coherent-module}
For a coherent connective commutative ring spectrum \(B\), write
\[
    \operatorname{Coh}(B)
    \subseteq
    \operatorname{APerf}(B)
\]
for the full subcategory of eventually coconnective almost-perfect
\(B\)-modules. A \textbf{connective coherent \(B\)-module} is an object of
\[
    \operatorname{Coh}(B)\cap(\Mod_B)_{\geq0}.
\]
\end{definition}

\begin{proposition}[Almost-perfect modules over a coherent commutative ring spectrum]
\label{prop:coherent-almost-perfect-homotopy}
Let \(B\) be a coherent connective commutative ring spectrum. A
\(B\)-module \(M\) is almost perfect if and only if it is bounded below and,
for every integer \(n\), the \(\pi_0(B)\)-module \(\pi_n(M)\) is finitely
presented. Consequently
\[
    \operatorname{Coh}(B)
\]
is exactly the full subcategory of bounded \(B\)-modules whose homotopy
modules are finitely presented over \(\pi_0(B)\).
\end{proposition}

\begin{proof}
\leavevmode
\begin{enumerate}[label=\textup{(\arabic*)}]
    \item \textbf{Invoke the coherent-module characterization.}
    This is the commutative \(\mathbb E_\infty\)-specialization of
    \cite[Proposition~7.2.4.17]{LurieHigherAlgebra}.
    \item \textbf{Identify the coherent subcategory.}
    The final statement is the
    intersection with the eventually coconnective subcategory in
    \Ref{def:coherent-module}.
\end{enumerate}
\end{proof}

\begin{lemma}[Relative spectral Hurewicz in the first unresolved degree]
\label{lem:relative-spectral-hurewicz-first-degree}
Let \(A\to C\) be a morphism of connective commutative ring spectra inducing
an isomorphism on \(\pi_0\), and suppose
\[
    Q:=\operatorname{cofib}(A\to C)
\]
is \(n\)-connective for \(n\geq1\). Then the canonical relative-cofibre
Hurewicz morphism
\[
    C\otimes_AQ
    \longrightarrow
    L_{C/A}
\]
induces an isomorphism
\[
    \pi_n(Q)
    \xrightarrow{\cong}
    \pi_nL_{C/A}.
\]

The comparison factors through extension of scalars:
\[
\begin{tikzcd}[column sep=huge]
    Q \arrow[r]
        & C\otimes_AQ \arrow[r]
        & L_{C/A}.
\end{tikzcd}
\]
\end{lemma}

\begin{proof}
\leavevmode
\begin{enumerate}[label=\textup{(\arabic*)}]
    \item \textbf{Hurewicz comparison.}
    By the relative Hurewicz theorem of
    \Ref{thm:relative-hurewicz-conormal}, the fibre of
    \[
        C\otimes_AQ
        \longrightarrow
        L_{C/A}
    \]
    is \(2n\)-connective. Since \(n\geq1\), the induced morphism on the first
    possible nonzero homotopy group is therefore an isomorphism:
    \[
        \pi_n(C\otimes_AQ)
        \xrightarrow{\cong}
        \pi_nL_{C/A}.
    \]
    \item \textbf{Identify the first homotopy group.}
    It remains to identify the source. Because \(A\), \(C\), and \(Q\) are
    connective, the connective tensor-product spectral sequence applies. The
    assumption
    \[
        \pi_0(A)\xrightarrow{\cong}\pi_0(C)
    \]
    and \(n\)-connectivity of \(Q\) imply
    \[
        \pi_n(C\otimes_AQ)
        \simeq
        \pi_0(C)\otimes_{\pi_0(A)}\pi_n(Q)
        \simeq
        \pi_n(Q).
    \]
    Combining the two canonical identifications gives
    \[
        \pi_n(Q)
        \xrightarrow{\cong}
        \pi_nL_{C/A}.
    \]
\end{enumerate}
\end{proof}

Almost-perfectness reduces the first unresolved homotopy to finite module data. To
improve connectivity without losing that finiteness, we need an algebraic operation
whose linear layer is prescribed while higher word-length layers become progressively
more connective. Relation cells provide exactly that operation.

\subsection{Relation cells and algebraic finiteness}

\begin{lemma}[Suspended free algebras are finite-cell]
\label{lem:suspended-free-algebras-finite-cell}
Let \(C\) be connective. For every \(n,r\geq0\), the commutative
\(C\)-algebra
\[
    \operatorname{Sym}_C(C[n]^{\oplus r})
\]
is finite-cell in the sense of Chapter~\Ref{chap:spectral-algebraic-geometry}.
\end{lemma}

\begin{proof}
\leavevmode
\begin{enumerate}[label=\textup{(\arabic*)}]
    \item \textbf{Express the suspension as a pushout.}
    In the stable category \(\Mod_C\),
    \[
        C[1]\simeq0\amalg_C0.
    \]
    \item \textbf{Apply the free algebra functor.}
    The free commutative-algebra functor preserves colimits, so
    \[
        \operatorname{Sym}_C(C[1])
        \simeq
        C\amalg_{\operatorname{Sym}_C(C)}C.
    \]
    \item \textbf{Recognize the suspended generator as finite-cell.}
    Thus the free algebra on a once-suspended generator lies in the finite-colimit
    closure of the degree-zero free generators used in the definition of finite
    cellularity.
    \item \textbf{Iterate the construction.}
    Iterating suspension and taking finite sums proves the general
    case.
\end{enumerate}
\end{proof}


\begin{lemma}[Finite-group homotopy orbits and spectral symmetric powers]
\label{lem:finite-group-orbits-almost-perfect}
Let \(A\) be connective, or let \(F^\bullet A\) be pointwise connective in
\(\operatorname{Fil}(\Sp)\) or \(\operatorname{Gr}(\Sp)\).
\begin{enumerate}[label=(\roman*)]
    \item If \(P\) is a connective finite cell module carrying an action of a
    finite group \(G\), then \(P_{hG}\) is connective and almost perfect.
    More generally, if \(P\) is \(q\)-connective, then so is \(P_{hG}\).
    \item If \(M\) is connective and almost perfect, then for every
    \(r\geq0\),
    \[
        \operatorname{Sym}^r_A(M)
        \simeq
        (M^{\otimes_A r})_{h\Sigma_r}
    \]
    is connective and almost perfect. The same assertion holds in the
    filtered and graded module categories with their pointwise
    \(t\)-structures.
\end{enumerate}
\end{lemma}

\begin{proof}
\leavevmode
\begin{enumerate}[label=\textup{(\arabic*)}]
    \item \textbf{Finite-group homotopy orbits.}
    For \textup{(i)}, compute homotopy orbits by the simplicial bar construction.
    In every simplicial degree its underlying module is a finite direct sum of
    copies of \(P\), because \(G\) is finite. Every finite skeletal realization
    is therefore a finite cell module. The skeletal filtration has successive
    cofibres given by increasingly suspended latching cofibres; consequently the
    map from the \(s\)-skeleton to \(P_{hG}\) becomes arbitrarily connective as
    \(s\to\infty\). The finite-cell approximation criterion for almost-perfect
    modules, and its filtered or graded analogue, proves almost-perfectness.
    Connectivity is preserved because homotopy orbits are colimits of connective
    objects.
    \item \textbf{Spectral symmetric powers.}
    For \textup{(ii)}, fix a Postnikov degree \(d\). Choose a connective finite
    cell approximation
    \[
        P\longrightarrow M
    \]
    whose cofibre is \(d\)-connective; after increasing \(d\) we may take the
    approximating module itself connective. Replacing the tensor factors one at a
    time shows that
    \[
        P^{\otimes_A r}\longrightarrow M^{\otimes_A r}
    \]
    has \(d\)-connective cofibre, since every successive cofibre contains the
    \(d\)-connective module \(\operatorname{cofib}(P\to M)\) tensored with
    connective factors. Finite-group homotopy orbits preserve this connectivity.
    Thus
    \[
        \operatorname{Sym}^r_A(P)
        \longrightarrow
        \operatorname{Sym}^r_A(M)
    \]
    is \(d\)-connective. By \textup{(i)}, the source is almost perfect, hence
    admits a finite-cell approximation through degree \(d\). Composing the two
    approximations and varying \(d\) proves that
    \(\operatorname{Sym}^r_A(M)\) is almost perfect. The proof is pointwise in
    the filtered and graded cases.
\end{enumerate}
\end{proof}



\begin{lemma}[Positive relation-cell attachment]
\label{lem:positive-relation-cell}
Let \(C\) be connective, let \(n\geq0\), let
\[
    P\simeq C[n]^{\oplus r}
\]
be finite free and connective, and let \(\chi:P\to C\) be \(C\)-linear. Put
\[
    S:=\operatorname{Sym}_C(P),
\]
let \(\epsilon_\chi,\epsilon_0:S\rightrightarrows C\) be the maps classified
by \(\chi\) and zero, and define
\[
    \operatorname{RelCell}_C(\chi)
    :=
    C_{\epsilon_\chi}\otimes_SC_{\epsilon_0}.
\]
Then:
\begin{enumerate}[label=(\roman*)]
    \item \(\operatorname{RelCell}_C(\chi)\) is a finite-cell commutative \(C\)-algebra;
    \item the two-sided bar construction has a natural exhaustive
    multiplicative filtration with
    \[
        \operatorname{Gr}^k\operatorname{RelCell}_C(\chi)
        \simeq
        \operatorname{Sym}^k_C(P[1]);
    \]
    \item \(\operatorname{RelCell}_C(\chi)\) is almost perfect as a \(C\)-module.
\end{enumerate}
Now suppose that \(C\to D\) induces an isomorphism on \(\pi_0\), that
\[
    Q:=\operatorname{cofib}(C\to D)
\]
is \((n+1)\)-connective, and that
\[
    u:P\longrightarrow Q[-1]
\]
is surjective on \(\pi_n\). Let \(\partial:Q[-1]\to C\) be the boundary
of the cofibre sequence and define
\[
    \chi:=\partial\circ u.
\]
Then the canonical nullhomotopy of \(P\to C\to D\) induces a map
\[
    \operatorname{RelCell}_C(\chi)\longrightarrow D
\]
whose cofibre is \((n+2)\)-connective.
\end{lemma}

\begin{proof}
\leavevmode
\begin{enumerate}[label=\textup{(\arabic*)}]
    \item \textbf{Finite-cell algebra.}
    By \Ref{lem:suspended-free-algebras-finite-cell},
    \(\operatorname{Sym}_C(P)\) is finite-cell. The pushout defining
    \(\operatorname{RelCell}_C(\chi)\) is a finite colimit of finite-cell commutative \(C\)-algebras,
    which proves \textup{(i)}.
    \item \textbf{Bar filtration.}
    Compute the pushout by the two-sided bar construction
    \[
        \operatorname{Bar}_\bullet
        (C_{\epsilon_\chi},S,C_{\epsilon_0}).
    \]
    Filter its normalized realization by total symmetric degree in the occurrences
    of the generator \(P\). Internal multiplication faces preserve total degree,
    the endpoint differential determined by \(\chi\) lowers total degree, and the
    zero endpoint kills positive degree. Hence \(\chi\) disappears on associated
    graded. The associated graded is therefore the zero--zero derived
    self-intersection
    \[
    \begin{aligned}
        C\otimes_{\operatorname{Sym}_C(P)}C
        &\simeq
        \operatorname{Sym}_C(0)
          \amalg_{\operatorname{Sym}_C(P)}
        \operatorname{Sym}_C(0)\\
        &\simeq
        \operatorname{Sym}_C(0\amalg_P0)\\
        &\simeq
        \operatorname{Sym}_C(P[1]).
    \end{aligned}
    \]
    Thus the weight-\(k\) layer is
    \(\operatorname{Sym}^k_C(P[1])\), proving \textup{(ii)}.
    \item \textbf{Almost-perfectness.}
    Since \(P[1]\) is \((n+1)\)-connective,
    \(\operatorname{Sym}^k_C(P[1])\) is \(k(n+1)\)-connective. It is almost
    perfect by \Ref{lem:finite-group-orbits-almost-perfect}. Consequently every
    fixed Postnikov range sees only finitely many word-length layers. Finite
    extensions of almost-perfect modules are almost perfect, and the omitted tail
    can be made arbitrarily connective. This gives finite-cell approximations in
    every Postnikov range and proves \textup{(iii)}.
    \item \textbf{Connectivity improvement.}
    For the final statement, the identity \(\chi=\partial u\) supplies a canonical
    nullhomotopy of \(P\to C\to D\), hence a morphism
    \(\operatorname{RelCell}_C(\chi)\to D\). The
    induced morphism from the relative cofibre of
    \(C\to\operatorname{RelCell}_C(\chi)\) to \(Q\) has
    linear layer
    \[
        P[1]\xrightarrow{u[1]}Q,
    \]
    which is surjective on \(\pi_{n+1}\). Every nonlinear layer contains at least
    two copies of \(P[1]\), and is therefore at least
    \(2(n+1)\)-connective. Since \(2(n+1)\geq n+2\), the cofiber of the map from
    the relative cofibre of \(C\to \operatorname{RelCell}_C(\chi)\) to \(Q\) is
    \((n+2)\)-connective. Octahedrality for
    \[
        C\longrightarrow \operatorname{RelCell}_C(\chi)\longrightarrow D
    \]
    identifies this cofibre with \(\operatorname{cofib}(\operatorname{RelCell}_C(\chi)\to D)\), proving
    the assertion.
\end{enumerate}
\end{proof}


\begin{theorem}[Spectral criterion for almost finite presentation of a surjective augmentation]
\label{thm:spectral-almost-finite-criterion}
Let \(A\to B\) be a morphism of connective commutative ring spectra inducing
a surjection on \(\pi_0\). The following conditions are equivalent:
\begin{enumerate}[label=(\roman*)]
    \item \(B\) is almost finitely presented as a connective commutative
    \(A\)-algebra;
    \item \(L_{B/A}\) is almost perfect and \(\pi_0(B)\) is finitely
    presented as a \(\pi_0(A)\)-algebra;
    \item \(B\) is almost perfect as an \(A\)-module.
\end{enumerate}
The criterion is valid for arbitrary connective commutative ring spectra satisfying
the displayed hypotheses.
\end{theorem}

\begin{proof}
\leavevmode
\begin{enumerate}[label=\textup{(\arabic*)}]
    \item \textbf{Cotangent criterion.}
    The equivalence of \textup{(i)} and \textup{(ii)} is the connective
    \(\mathbb E_\infty\)-cotangent criterion for almost finite presentation;
    see \cite[Theorem~2.12(2)]{LurieDAGIV}. The same criterion is recorded in
    \cite[Proposition~A.19(1)--(2)]{BrantnerMagidsonNuiten}.
    \item \textbf{From module finiteness to algebra finiteness.}
    The implication \textup{(iii)}\(\Rightarrow\)\textup{(i)} follows from
    \cite[Proposition~A.19(3)]{BrantnerMagidsonNuiten}: almost-perfectness of
    \(B\) as an \(A\)-module implies almost finite presentation of the algebra
    map \(A\to B\).
    \item \textbf{From algebra finiteness to module finiteness.}
    For the reverse implication, assume \textup{(i)}. The standing hypothesis
    gives a surjection
    \[
        \pi_0(A)\twoheadrightarrow\pi_0(B).
    \]
    Then \cite[Proposition~A.19(4)]{BrantnerMagidsonNuiten} implies
    \[
        B\in\operatorname{APerf}(A).
    \]
    Thus \textup{(i)}, \textup{(ii)}, and \textup{(iii)} are equivalent.
\end{enumerate}
\end{proof}

\begin{corollary}[Cotangent criterion for almost finite presentation]
\label{cor:cotangent-almost-finite-general}
Let \(A\to C\) be a morphism of connective commutative ring spectra. Suppose
\(\pi_0(C)\) is finitely presented as a \(\pi_0(A)\)-algebra and
\(L_{C/A}\) is almost perfect as a \(C\)-module. Then \(C\) is almost
finitely presented over \(A\).
\end{corollary}

\begin{proof}
\leavevmode
\begin{enumerate}[label=\textup{(\arabic*)}]
    \item \textbf{Apply the cotangent criterion.}
    Apply the connective \(\mathbb E_\infty\)-cotangent criterion for almost
    finite presentation,
    \cite[Theorem~2.12(2)]{LurieDAGIV}.
    \item \textbf{Verify its hypotheses.}
    Its hypotheses are precisely finite
    presentation of \(\pi_0(C)\) over \(\pi_0(A)\) together with
    almost-perfectness of \(L_{C/A}\).
    \item \textbf{Record the equivalent formulation.}
    The equivalent formulation in
    \cite[Proposition~A.19(1)--(2)]{BrantnerMagidsonNuiten} gives the same
    conclusion.
\end{enumerate}
\end{proof}

\begin{lemma}[Finite presentation of the degree-zero augmentation kernel]
\label{lem:augmentation-kernel-finitely-presented}
Let \(A\to B\) be a morphism of connective commutative ring spectra which is
surjective on \(\pi_0\). Suppose that \(B\) is coherent and almost perfect as
an \(A\)-module. Put
\[
    I:=\ker\bigl(\pi_0(A)\to\pi_0(B)\bigr).
\]
Then \(I\) is finitely presented as a \(\pi_0(A)\)-module.
\end{lemma}

\begin{proof}
\leavevmode
\begin{enumerate}[label=\textup{(\arabic*)}]
    \item \textbf{Almost-perfect fibre.}
    Put \(K:=\operatorname{fib}(A\to B)\). Since the map is surjective on
    \(\pi_0\), the long exact homotopy sequence makes \(K\) connective. The
    almost-perfect subcategory is stable under finite limits and colimits; since
    \(A\) is perfect over itself and \(B\) is almost perfect over \(A\), the fibre
    \(K\) is almost perfect. For a connective almost-perfect module,
    \(\pi_0\) is finitely presented over \(\pi_0(A)\). Thus \(\pi_0(K)\) is
    finitely presented.
    \item \textbf{Finite restriction of the first higher homotopy.}
    Apply \Ref{lem:first-unresolved-homotopy-almost-perfect} to the connective
    almost-perfect \(A\)-module \(B\). It gives
    \[
        \pi_0(B)\in\operatorname{Mod}_{\pi_0(A)}^{\mathrm{fp}}.
    \]
    Coherence of \(B\) makes \(\pi_1(B)\) finitely presented over
    \(\pi_0(B)\). Choose a finite presentation of \(\pi_1(B)\) by finite free
    \(\pi_0(B)\)-modules. After restriction of scalars, the two finite free
    terms are finite direct sums of the finitely presented
    \(\pi_0(A)\)-module \(\pi_0(B)\). Hence \(\pi_1(B)\), and in particular
    its image in any \(\pi_0(A)\)-module, is finitely generated over
    \(\pi_0(A)\).
    \item \textbf{Degree-zero exact sequence.}
    The fibre sequence gives an exact sequence
    \[
        \pi_1(B)
        \longrightarrow
        \pi_0(K)
        \longrightarrow
        \pi_0(A)
        \longrightarrow
        \pi_0(B)
        \longrightarrow0.
    \]
    Thus there is a canonical exact sequence
    \[
        \operatorname{im}\bigl(\pi_1(B)\to\pi_0(K)\bigr)
        \longrightarrow
        \pi_0(K)
        \longrightarrow
        I
        \longrightarrow0.
    \]
    The middle module is finitely presented over \(\pi_0(A)\) by the first
    step of the proof, while the image on the left is finitely generated by
    the preceding restriction-of-scalars argument. A quotient of a finitely
    presented module by a finitely generated submodule is finitely presented.
    Therefore \(I\) is finitely presented over \(\pi_0(A)\).
\end{enumerate}
\end{proof}

\begin{lemma}[The positive graded symmetric augmentation is almost finitely presented]
\label{lem:positive-graded-symmetric-aft}
Let \(B\) be connective and let \(M\) be a connective almost-perfect
\(B\)-module. Then the augmentation
\[
    \operatorname{Sym}^{\mathrm{gr}}_B(M\langle1\rangle)
    \longrightarrow
    B
\]
is almost finitely presented in connective nonnegatively graded commutative
\(B\)-algebras.
\end{lemma}

\begin{proof}
\leavevmode
\begin{enumerate}[label=\textup{(\arabic*)}]
    \item \textbf{Fix the graded algebra and test object.}
    Put \(C:=\operatorname{Sym}^{\mathrm{gr}}_B(M\langle1\rangle)\). Let \(D\) be an
    \(n\)-truncated connective graded commutative \(C\)-algebra.
    \item \textbf{Apply the free graded-algebra adjunction.}
    By the free
    graded-algebra adjunction, a morphism of augmented graded \(B\)-algebras from
    \(B\) to \(D\) over \(C\) is equivalently a nullhomotopy, compatible with the
    augmentation, of the image of the universal weight-one map
    \[
        M\longrightarrow D_1.
    \]
    \item \textbf{Express the mapping space as a finite limit.}
    Thus its mapping space is a finite-limit expression built from
    \(\Map_{\Mod_B}(M,D_1)\).
    \item \textbf{Use almost-perfectness and weight evaluation.}
    Evaluation in weight one preserves filtered colimits, and
    almost-perfectness of \(M\) makes mapping out of \(M\) commute with filtered
    colimits on every fixed truncated connective module category.
    \item \textbf{Commute filtered colimits with finite limits.}
    Filtered
    colimits of spaces commute with finite limits, proving the claim.
\end{enumerate}
\end{proof}

\begin{lemma}[The positive symmetric augmentation module is almost perfect]
\label{lem:positive-graded-symmetric-augmentation-module}
Let \(B\) be connective and let \(M\) be a connective almost-perfect
\(B\)-module. Put
\[
    C:=\operatorname{Sym}^{\mathrm{gr}}_B(M\langle1\rangle).
\]
Then the weight-zero augmentation module \(B\), concentrated in graded
weight zero, is almost perfect as a graded \(C\)-module.
\end{lemma}

\begin{proof}
\leavevmode
\begin{enumerate}[label=\textup{(\arabic*)}]
    \item \textbf{Resolve the augmentation module.}
    Resolve the augmentation module by the two-sided bar construction for
    \(C\to B\).
    \item \textbf{Filter by word length.}
    Filter its normalization by word length in the positive
    weight-one generator.
    \item \textbf{Identify the homogeneous terms.}
    Every homogeneous term is assembled from finite tensor
    powers of \(M[1]\) and finite-group homotopy orbits.
    \item \textbf{Approximate each fixed word length.}
    For a fixed word length
    and a fixed Postnikov range, a sufficiently large finite skeleton of the
    corresponding finite-group free resolution gives a perfect approximation.
    \item \textbf{Use increasing connectivity.}
    Because \(M\) is connective and occurs in positive weight, the connectivity of
    word length \(k\) tends to \(+\infty\) with \(k\).
    \item \textbf{Bound the relevant word lengths.}
    Hence every fixed
    Postnikov range sees only finitely many word lengths.
    \item \textbf{Assemble the perfect approximations.}
    Finite extensions of the
    resulting perfect approximations give a perfect graded \(C\)-module
    approximation to \(B\) through every Postnikov degree.
\end{enumerate}
\end{proof}

The preceding relation-cell calculus is unfiltered. To use it for the canonical adic
filtration, we now transport the same approximation argument to filtered modules.
Compact generators indexed by filtration degree make almost-perfectness accessible to
finite filtered cells, while completion lets associated graded detect the result.

\subsection{Filtered cell approximation and complete detection}

\begin{construction}[Compact filtered weight generators]
\label{con:free-filtered-weight-generators}
For \(w\geq0\), let \(\mathbb S\langle w\rangle\in\operatorname{Fil}(\Sp)\)
be the filtered sphere
\[
    F^n\mathbb S\langle w\rangle
    =
    \begin{cases}
        \mathbb S,&0\leq n\leq w,\\
        0,&n>w,
    \end{cases}
\]
with identity transition maps in the nonzero range. Its associated graded is
the sphere concentrated in weight \(w\). For a filtered commutative ring
spectrum \(F^\bullet A\), put
\[
    F^\bullet A\langle w\rangle
    :=
    F^\bullet A\otimes\mathbb S\langle w\rangle.
\]
Then for every filtered \(A\)-module \(N\) there is a natural equivalence
\[
    \map_{F^\bullet A}
    \bigl(F^\bullet A\langle w\rangle,N\bigr)
    \simeq
    F^wN,
\]
and
\[
    \operatorname{Gr}
    \bigl(F^\bullet A\langle w\rangle\bigr)
    \simeq
    \operatorname{Gr}(A)\langle w\rangle.
\]
The objects \(F^\bullet A\langle w\rangle\), together with their suspensions,
form a family of compact generators of the filtered module category. A
\textbf{finite filtered cell module} is an object in the smallest full subcategory
containing these suspended weight generators and closed under finite colimits.
The perfect filtered modules form the thick subcategory generated by the same
weight generators. Thus every finite filtered cell module is perfect, and every
perfect filtered module is a retract of a finite filtered cell module.
\end{construction}

\begin{proof}
\leavevmode
\begin{enumerate}[label=\textup{(\arabic*)}]
    \item \textbf{Represent filtration evaluation.}
    The filtered sphere \(\mathbb S\langle w\rangle\) represents evaluation in
    filtration degree \(w\).
    \item \textbf{Apply free-module adjunction.}
    Free-module adjunction gives the mapping formula.
    \item \textbf{Prove compactness of the weight generators.}
    Evaluation preserves filtered colimits, so the weight generators are compact.
    \item \textbf{Prove generation.}
    If all mapping spectra from all suspensions of all weight generators vanish,
    then every filtration term of \(N\) vanishes; hence they generate.
    \item \textbf{Compute associated graded.}
    Associated
    graded is symmetric monoidal and carries \(\mathbb S\langle w\rangle\) to
    the graded sphere in weight \(w\), proving the associated-graded formula.
    \item \textbf{Identify finite filtered cells as perfect.}
    Finite colimits of compact objects are compact, so every finite filtered
    cell module is compact and therefore perfect.
    \item \textbf{Identify perfect modules as retracts.}
    Compact objects in a compactly
    generated stable category form the thick subcategory generated by any set of
    compact generators; equivalently, every perfect filtered module is a retract
    of a finite filtered cell module.
\end{enumerate}
\end{proof}

\begin{lemma}[Finite-cell approximation criterion for filtered almost-perfect modules]
\label{lem:filtered-aperf-cell-criterion}
Let \(F^\bullet A\) be pointwise connective and let \(N\) be a bounded-below
filtered \(A\)-module. Then \(N\) is almost perfect if and only if, for every
integer \(d\), there exists a finite filtered cell module \(P_d\) and a map
\[
    P_d\longrightarrow N
\]
whose cofibre is \(d\)-connective in the pointwise Postnikov \(t\)-structure.
\end{lemma}

\begin{proof}
\leavevmode
\begin{enumerate}[label=\textup{(\arabic*)}]
    \item \textbf{Inductive approximation.}
    After a suspension, assume that \(N\) is connective. Suppose first that
    \(N\) is almost perfect. Set
    \[
        C_0:=N,
        \qquad
        P_0:=0.
    \]
    Assume inductively that
    \[
    \begin{tikzcd}[column sep=large]
        P_d \arrow[r] & N \arrow[r] & C_d
    \end{tikzcd}
    \]
    is a cofiber sequence with \(P_d\) a finite filtered cell and \(C_d\)
    pointwise \(d\)-connective. Since finite filtered cells are perfect and
    almost-perfect modules are closed under finite cofibers, \(C_d\) is almost
    perfect. Therefore
    \[
        \tau_{\leq d}C_d
        \simeq
        H\pi_d(C_d)[d]
    \]
    is compact in the category of pointwise \(d\)-truncated filtered
    \(A\)-modules, so \(\pi_d(C_d)\) is finitely presented in the heart.
    \item \textbf{Kill the first unresolved homotopy.}
    For each \(w\geq0\), the free weight module \(A\langle w\rangle\) represents
    evaluation in filtration degree \(w\). Its heart object
    \(\pi_0(A\langle w\rangle)\) is therefore projective, and these heart objects
    form a family of generators. Finite presentation of \(\pi_d(C_d)\) gives
    weights \(w_1,\ldots,w_r\) and an epimorphism
    \[
        \bigoplus_{i=1}^r
        \pi_0\bigl(A\langle w_i\rangle\bigr)
        \twoheadrightarrow
        \pi_d(C_d).
    \]
    Choose representatives of these classes to obtain
    \[
        Q_d:=
        \bigoplus_{i=1}^rA\langle w_i\rangle[d]
        \longrightarrow
        C_d
    \]
    surjective on pointwise \(\pi_d\). Let \(C_{d+1}\) be its cofiber and put
    \[
        P_{d+1}:=\operatorname{fib}(N\to C_{d+1}).
    \]
    Then \(C_{d+1}\) is pointwise \((d+1)\)-connective, and the \(3\times3\)
    lemma supplies a cofiber sequence
    \[
    \begin{tikzcd}[column sep=large]
        P_d \arrow[r] & P_{d+1} \arrow[r] & Q_d .
    \end{tikzcd}
    \]
    Thus \(P_{d+1}\) is again a finite filtered cell. Induction yields finite
    filtered-cell approximations with arbitrarily connective cofiber.
    \item \textbf{Converse implication.}
    Conversely, suppose that such approximations exist. Fix \(n\) and choose
    \(P\to N\) with \((n+1)\)-connective cofiber. Then
    \[
        \tau_{\leq n}P\xrightarrow{\simeq}\tau_{\leq n}N.
    \]
    A finite filtered cell is compact, and mapping into a pointwise
    \(n\)-truncated module depends only on its \(n\)-truncation. Hence
    \(\tau_{\leq n}N\) is compact in the pointwise \(n\)-truncated module
    category. Varying \(n\) proves that \(N\) is almost perfect. Desuspending
    gives the bounded-below case.
\end{enumerate}
\end{proof}


\begin{lemma}[Restriction transitivity for filtered almost-perfect modules]
\label{lem:filtered-almost-perfect-restriction-transitivity}
Let
\[
    F^\bullet A\longrightarrow F^\bullet B
\]
be a morphism of pointwise connective filtered commutative ring spectra. If
\(F^\bullet B\) is almost perfect as a filtered \(F^\bullet A\)-module, then
restriction of scalars carries almost-perfect filtered \(B\)-modules to
almost-perfect filtered \(A\)-modules.
\end{lemma}

\begin{proof}
\leavevmode
\begin{enumerate}[label=\textup{(\arabic*)}]
    \item \textbf{Filtered \(B\)-cell approximation.}
    After a suspension it is enough to treat a connective almost-perfect filtered
    \(B\)-module \(M\). For every integer \(d\),
    \Ref{lem:filtered-aperf-cell-criterion} supplies a finite filtered
    \(B\)-cell
    \[
        P_d\longrightarrow M
    \]
    with \(d\)-connective cofibre.
    \item \textbf{Restriction to \(A\).}
    Every filtered weight generator \(B\langle w\rangle\), after restriction to
    \(A\), is almost perfect. Indeed \(B\) is almost perfect over \(A\), and
    tensoring with the finite filtered sphere \(\mathbb S\langle w\rangle\)
    preserves finite-cell approximations and their pointwise connectivity. The
    almost-perfect subcategory is stable under finite limits, colimits, shifts,
    and retracts, so the restricted module underlying every finite filtered
    \(B\)-cell \(P_d\) is almost perfect over \(A\).
    \item \textbf{Finite \(A\)-cell approximation.}
    Apply \Ref{lem:filtered-aperf-cell-criterion} to the restricted module
    \(P_d\). Choose a finite filtered \(A\)-cell
    \[
        Q_d\longrightarrow P_d
    \]
    with \(d\)-connective cofibre. The composite
    \[
        Q_d\longrightarrow P_d\longrightarrow M
    \]
    again has \(d\)-connective cofibre by octahedrality. Since this can be done
    for every \(d\), the filtered finite-cell criterion makes \(M\) almost
    perfect over \(A\).
\end{enumerate}
\end{proof}


\begin{theorem}[Complete filtered almost-perfect detection]
\label{thm:filtered-almost-perfect-detection}
Let \(F^\bullet A\) be a complete pointwise connective filtered commutative
ring spectrum, and let \(F^\bullet N\) be a pointwise connective filtered
\(A\)-module. Then the following are equivalent:
\begin{enumerate}[label=(\roman*)]
    \item \(F^\bullet N\) is almost perfect in
    \(\Mod_{F^\bullet A}(\operatorname{Fil}(\Sp))\);
    \item \(F^\bullet N\) is complete and
    \(\operatorname{Gr}N\) is almost perfect over
    \(\operatorname{Gr}A\).
\end{enumerate}
Under these conditions, \(F^0N\) is almost perfect as an
\(F^0A\)-module.
\end{theorem}

\begin{proof}
\leavevmode
\begin{enumerate}[label=\textup{(\arabic*)}]
    \item \textbf{Associated graded.}
    Assume first that \(N\) is almost perfect. By
    \Ref{lem:filtered-aperf-cell-criterion}, for every \(d\) choose a finite
    filtered cell module
    \[
        P_d\longrightarrow N
    \]
    with \(d\)-connective cofibre \(C_d\). Associated graded sends each compact
    weight generator to the corresponding compact graded weight generator and is
    exact. Hence it sends finite filtered cells to finite graded cells, and the
    same approximations show that \(\operatorname{Gr}N\) is almost perfect.
    \item \textbf{Completeness.}
    Every weight generator \(A\langle w\rangle\) is complete when \(A\) is
    complete: its descending tail is a finite filtration shift of the tail of
    \(A\). Completeness is stable under finite fibres, cofibres, sums, and
    retracts, so every \(P_d\) is complete. Therefore
    \[
        \varprojlim\nolimits_nF^nN
        \simeq
        \varprojlim\nolimits_nF^nC_d.
    \]
    Each \(F^nC_d\) is \(d\)-connective. The Milnor exact sequence makes the
    inverse limit \((d-1)\)-connective. Since \(d\) is arbitrary, the inverse
    limit vanishes and \(N\) is complete.
    \item \textbf{Filtration zero.}
    Evaluation in filtration zero is exact and sends every weight generator
    \(A\langle w\rangle\) to the free rank-one \(F^0A\)-module. Hence it sends
    finite filtered cell modules to perfect \(F^0A\)-modules and preserves the
    connectivity of the cofibre of each approximation. Thus \(F^0N\) is almost
    perfect over \(F^0A\).
    \item \textbf{Inductive converse.}
    Conversely suppose that \(N\) is complete and \(\operatorname{Gr}N\) is
    almost perfect. Set \(C_0=N\). Suppose inductively that \(C_d\) is complete,
    pointwise \(d\)-connective, and has almost-perfect associated graded. Then
    \(\operatorname{Gr}C_d\) is \(d\)-connective, and
    \Ref{lem:first-unresolved-homotopy-almost-perfect} in the graded module
    category makes
    \[
        \pi_d\operatorname{Gr}C_d
    \]
    a finitely presented graded module. Choose finitely many homogeneous
    generators in weights \(w_1,\ldots,w_r\).
    \item \textbf{Lift graded generators.}
    For each \(w_i\), the cofibre sequence
    \[
        F^{w_i+1}C_d
        \longrightarrow
        F^{w_i}C_d
        \longrightarrow
        \operatorname{Gr}^{w_i}C_d
    \]
    and pointwise \(d\)-connectivity give a surjection on \(\pi_d\). Choose lifts
    of the selected generators. They assemble to a morphism
    \[
        \bigoplus_{i=1}^r
        F^\bullet A\langle w_i\rangle[d]
        \longrightarrow
        C_d.
    \]
    Let \(C_{d+1}\) be its cofibre. Associated graded is exact, so the chosen
    classes kill the first possible homotopy group and
    \(\operatorname{Gr}C_{d+1}\) is \((d+1)\)-connective. It remains almost
    perfect. The source is a finite filtered cell and is complete; hence
    \(C_{d+1}\) is complete.
    \item \textbf{Recover pointwise connectivity.}
    Apply \Ref{lem:complete-positive-connectivity} to every shifted filtration tail
    of \(C_{d+1}\). Since each associated-graded term is
    \((d+1)\)-connective, every filtration term of \(C_{d+1}\) is
    \((d+1)\)-connective. Iterating and splicing the finite cell attachments by
    octahedrality gives, through every Postnikov degree, a finite filtered-cell
    approximation to \(N\). By \Ref{lem:filtered-aperf-cell-criterion}, \(N\) is
    almost perfect.
\end{enumerate}
\end{proof}

\begin{corollary}[Almost-perfectness from a complete filtration]
\label{lem:almost-perfect-complete-positive-filtration}
Let \(F^\bullet A\) be complete and pointwise connective, and let
\(F^\bullet N\) be pointwise connective and complete. If \(\operatorname{Gr}N\) is almost perfect over
\(\operatorname{Gr}A\), then \(F^0N\) is almost perfect over \(F^0A\).
\end{corollary}

\begin{proof}
\leavevmode
\begin{enumerate}[label=\textup{(\arabic*)}]
    \item \textbf{Apply complete filtered detection.}
    Apply \Ref{thm:filtered-almost-perfect-detection}.
\end{enumerate}
\end{proof}

\begin{corollary}[Almost-perfect filtered modules are complete]
\label{lem:almost-perfect-filtered-complete}
Let \(F^\bullet A\) be a complete pointwise connective filtered commutative
ring spectrum. Every pointwise connective almost-perfect filtered
\(A\)-module is complete.
\end{corollary}

\begin{proof}
\leavevmode
\begin{enumerate}[label=\textup{(\arabic*)}]
    \item \textbf{Apply the forward detection implication.}
    Apply the forward implication of
    \Ref{thm:filtered-almost-perfect-detection}.
\end{enumerate}
\end{proof}

Complete filtered detection now controls objects once their graded layers are known. The
last missing estimate is cotangent control: the nonlinear Hurewicz comparison must
persist in filtered spectral algebra. The following results combine filtered relation
cells, Postnikov approximation, and cotangent transitivity to supply it.

\subsection{Filtered Hurewicz control and cotangent convergence}

\begin{lemma}[Filtered relative Hurewicz in the first unresolved degree]
\label{lem:filtered-relative-hurewicz-first-degree}
Let \(F^\bullet A\to F^\bullet C\) be a morphism of pointwise connective
filtered commutative ring spectra. Suppose its relative cofibre \(Q\) is
pointwise \(n\)-connective for some \(n\geq1\), and that the map is an
isomorphism on pointwise \(\pi_0\). Then \(L_{C/A}\) is pointwise
\(n\)-connective and the canonical map
\[
    F^\bullet C\otimes_{F^\bullet A}Q
    \longrightarrow
    L_{C/A}
\]
induces an isomorphism on pointwise \(\pi_n\). More generally, if
\(F^\bullet A\to F^\bullet C\) is pointwise surjective on \(\pi_0\), then
\(L_{C/A}\) is pointwise \(1\)-connective.
\end{lemma}

\begin{proof}
\leavevmode
\begin{enumerate}[label=\textup{(\arabic*)}]
    \item \textbf{Verify the Hurewicz hypotheses.}
    The stable presentable symmetric monoidal category
    \[
        \operatorname{Fil}(\Sp)
        =
        \Fun(\mathbb N_{\leq}^{\op},\Sp)
    \]
    with Day convolution and the pointwise Postnikov \(t\)-structure satisfies the
    hypotheses of the general Hurewicz theorem of
    \cite[Construction~7.4.3.10 and Theorem~7.4.3.12]{LurieHigherAlgebra}:
    Day convolution preserves colimits separately, its unit is pointwise
    connective, and tensor products of pointwise connective filtered spectra are
    pointwise connective.
    \item \textbf{Apply the filtered Hurewicz theorem.}
    Applying that theorem inside
    \(\operatorname{Fil}(\Sp)\) gives a canonical Hurewicz morphism
    \[
        F^\bullet C\otimes_{F^\bullet A}Q
        \longrightarrow
        L_{C/A}
    \]
    whose fibre is pointwise \(2n\)-connective when \(Q\) is pointwise
    \(n\)-connective.
    \item \textbf{Read off the first homotopy group.}
    This implies pointwise \(n\)-connectivity of
    \(L_{C/A}\) and the asserted isomorphism on the first possible homotopy
    object.
    \item \textbf{Handle the surjective degree-zero case.}
    If the original morphism is pointwise surjective on \(\pi_0\), its
    relative cofibre is pointwise \(1\)-connective, and the same theorem gives
    pointwise \(1\)-connectivity of the relative cotangent complex.
\end{enumerate}
\end{proof}


\begin{lemma}[Filtered positive relation-cell attachment]
\label{lem:filtered-positive-relation-cell}
Let \(F^\bullet C\) be a pointwise connective filtered commutative ring
spectrum, let \(n\geq0\), let \(P\) be a pointwise \(n\)-connective finite
filtered cell \(C\)-module, and let \(\chi:P\to C\) be \(C\)-linear. Put
\[
    S:=\operatorname{Sym}_C(P),
\]
let \(\epsilon_\chi,\epsilon_0:S\rightrightarrows C\) be the maps classified
by \(\chi\) and zero, and define
\[
    \operatorname{RelCell}_C(\chi)
    :=
    C_{\epsilon_\chi}\otimes_SC_{\epsilon_0}.
\]
Then \(\operatorname{RelCell}_C(\chi)\) is compact as a commutative filtered \(C\)-algebra, is
almost perfect as a filtered \(C\)-module, and its relative cotangent complex
is perfect. More precisely,
\[
    L_{\operatorname{RelCell}_C(\chi)/C}
    \simeq
    \operatorname{RelCell}_C(\chi)\otimes_C P[1].
\]
The two-sided bar construction has an exhaustive multiplicative filtration
whose word-length \(k\) associated-graded term is
\[
    \operatorname{Sym}^k_C(P[1]).
\]

Now let \(C\to D\) be a morphism of pointwise connective filtered
commutative ring spectra, put
\[
    Q:=\operatorname{cofib}(C\to D),
\]
and suppose \(Q\) is pointwise \((n+1)\)-connective. If
\[
    u:P\longrightarrow Q[-1]
\]
is surjective on the first possible homotopy object and
\(\partial:Q[-1]\to C\) is the boundary morphism, define
\[
    \chi:=\partial\circ u.
\]
Then the canonical nullhomotopy of \(P\to C\to D\) constructs a morphism
\[
    \operatorname{RelCell}_C(\chi)\longrightarrow D
\]
whose relative cofibre is pointwise \((n+2)\)-connective.
\end{lemma}

\begin{proof}
\leavevmode
\begin{enumerate}[label=\textup{(\arabic*)}]
    \item \textbf{Compactness.}
    Because \(P\) is a finite filtered cell, it is compact and perfect as a
    filtered \(C\)-module. The forgetful functor from commutative filtered
    \(C\)-algebras preserves filtered colimits, so the free-algebra adjunction
    makes \(\operatorname{Sym}_C(P)\) compact. Mapping out of the pushout
    \(\operatorname{RelCell}_C(\chi)\) is a finite pullback of mapping spaces, and filtered colimits of
    spaces commute with finite limits; hence \(\operatorname{RelCell}_C(\chi)\) is compact as a
    commutative filtered \(C\)-algebra.
    \item \textbf{Relative cotangent complex.}
    Cotangent base change for the defining pushout and transitivity for
    \[
        C\longrightarrow\operatorname{Sym}_C(P)
        \xrightarrow{\epsilon_0}C,
    \]
    whose composite is the identity, give
    \[
        L_{\operatorname{RelCell}_C(\chi)/C}
        \simeq
        \operatorname{RelCell}_C(\chi)\otimes_C P[1].
    \]
    This module is perfect.
    \item \textbf{Bar filtration.}
    For the underlying filtered module, compute the pushout by the two-sided bar
    construction and filter its normalized realization by total symmetric degree
    in the occurrences of \(P\). Internal multiplication preserves this degree,
    the endpoint differential determined by \(\chi\) lowers it, and the zero
    endpoint kills positive degree. Thus the associated graded is the zero--zero
    self-intersection
    \[
    \begin{aligned}
        C\otimes_{\operatorname{Sym}_C(P)}C
        &\simeq
        \operatorname{Sym}_C(0)
        \amalg_{\operatorname{Sym}_C(P)}
        \operatorname{Sym}_C(0)\\
        &\simeq
        \operatorname{Sym}_C(P[1]).
    \end{aligned}
    \]
    The word-length \(k\) layer is therefore
    \(\operatorname{Sym}^k_C(P[1])\). It is pointwise
    \(k(n+1)\)-connective and almost perfect by
    \Ref{lem:finite-group-orbits-almost-perfect}. Hence every fixed Postnikov
    range sees only finitely many word lengths; finite extensions of their
    finite-cell approximations prove that \(\operatorname{RelCell}_C(\chi)\) is almost perfect as a
    filtered \(C\)-module.
    \item \textbf{Connectivity improvement.}
    For the final statement, \(\chi=\partial u\) supplies a canonical
    nullhomotopy of \(P\to C\to D\), and hence a map
    \(\operatorname{RelCell}_C(\chi)\to D\). The
    linear word-length layer of the induced map on relative cofibres is
    \[
        P[1]\longrightarrow Q,
    \]
    which is surjective on the first possible homotopy object. Every nonlinear
    layer contains at least two copies of \(P[1]\) and is therefore pointwise at
    least \(2(n+1)\)-connective. Since \(2(n+1)\geq n+2\), octahedrality for
    \[
        C\longrightarrow \operatorname{RelCell}_C(\chi)\longrightarrow D
    \]
    shows that the remaining relative cofibre is pointwise
    \((n+2)\)-connective.
\end{enumerate}
\end{proof}



\begin{lemma}[Composition of compact spectral algebra morphisms]
\label{lem:compact-algebra-morphism-composition}
Let \(\mathcal C\) be \(\Sp\), \(\operatorname{Fil}(\Sp)\), or
\(\operatorname{Gr}(\Sp)\), and let
\[
    A\longrightarrow C\longrightarrow D
\]
be morphisms of commutative algebra objects. If \(C\) is compact as a
commutative \(A\)-algebra and \(D\) is compact as a commutative
\(C\)-algebra, then \(D\) is compact as a commutative \(A\)-algebra.
\end{lemma}

\begin{proof}
\leavevmode
\begin{enumerate}[label=\textup{(\arabic*)}]
    \item \textbf{Choose a filtered target diagram.}
    Let \(\{E_\alpha\}\) be a filtered diagram of commutative \(A\)-algebras and
    let
    \[
        \varphi:D\longrightarrow\operatorname*{colim}_\alpha E_\alpha
    \]
    be an \(A\)-algebra map.
    \item \textbf{Restrict the algebra map to \(C\).}
    Its restriction along \(C\to D\) is a point of
    \[
        \Map_{\CAlg(\mathcal C)_{A/}}
        \left(C,\operatorname*{colim}_\alpha E_\alpha\right).
    \]
    \item \textbf{Represent the \(C\)-structure at a finite stage.}
    Compactness of \(C\) implies that this \(C\)-algebra structure is represented
    at some stage \(E_{\alpha_0}\).
    \item \textbf{Pass to a cofinal diagram of \(C\)-algebras.}
    After restricting to the cofinal filtered
    subcategory of stages under \(\alpha_0\), all objects carry the induced
    compatible \(C\)-algebra structure.
    \item \textbf{Use compactness of \(D\).}
    Compactness of \(D\) in
    \(\CAlg(\mathcal C)_{C/}\) then implies that \(\varphi\), together with every higher homotopy
    between such maps, is represented at a finite stage.
    \item \textbf{Identify the mapping-space comparison.}
    Therefore the canonical
    map
    \[
        \operatorname*{colim}_\alpha
        \Map_{\CAlg(\mathcal C)_{A/}}(D,E_\alpha)
        \longrightarrow
        \Map_{\CAlg(\mathcal C)_{A/}}
        \left(D,\operatorname*{colim}_\alpha E_\alpha\right)
    \]
    is an equivalence.
    \item \textbf{Conclude compactness.}
    Thus \(D\) is compact as a commutative \(A\)-algebra.
\end{enumerate}
\end{proof}

\begin{lemma}[Filtered surjective almost finite presentation and cotangent criterion]
\label{lem:filtered-surjective-almost-finite}
Let
\[
    F^\bullet A\longrightarrow F^\bullet B
\]
be a morphism of pointwise connective filtered commutative ring spectra which
is surjective on \(\pi_0\) in every filtration degree. Suppose that
\(F^\bullet B\) is almost perfect as a filtered \(F^\bullet A\)-module. Then
the morphism is almost finitely presented in the sense of
\Ref{conv:filtered-graded-almost-finite-presentation}, and its relative
cotangent complex is an almost-perfect filtered \(F^\bullet B\)-module. The
shifted relative cotangent complex is pointwise connective.
\end{lemma}

\begin{proof}
\leavevmode
\begin{enumerate}[label=\textup{(\arabic*)}]
    \item \textbf{Initial relative cofiber.}
    Write \(A\) and \(B\) for the two filtered algebras and put
    \[
        Q_0:=\operatorname{cofib}(A\to B).
    \]
    Pointwise surjectivity on \(\pi_0\) makes \(Q_0\) pointwise
    \(1\)-connective. Since \(A\) is perfect over itself and \(B\) is almost
    perfect over \(A\), the relative cofibre \(Q_0\) is almost perfect as a
    filtered \(A\)-module.
    \item \textbf{Inductive relation-cell construction.}
    We construct compact filtered \(A\)-algebras
    \[
        C_d\longrightarrow B,
        \qquad d\geq0,
    \]
    starting with \(C_0=A\), such that
    \[
        Q_d:=\operatorname{cofib}(C_d\to B)
    \]
    is pointwise \((d+1)\)-connective and both \(C_d\) and \(Q_d\) are almost
    perfect as filtered \(A\)-modules. Suppose \(C_d\) has been constructed.
    Then \(Q_d[-1]\) is pointwise \(d\)-connective and almost perfect over
    \(A\). By \Ref{lem:first-unresolved-homotopy-almost-perfect} in the
    filtered module category, its first possible homotopy object is finitely
    presented in the pointwise heart. The images of the weight generators
    \(A\langle w\rangle\) generate that heart, so there are weights
    \(w_1,\ldots,w_r\) and a morphism
    \[
        G_d:=\bigoplus_{j=1}^rA\langle w_j\rangle[d]
        \longrightarrow
        Q_d[-1]
    \]
    surjective on pointwise \(\pi_d\).
    \item \textbf{Compact algebra approximants.}
    Base change the finite filtered cell to \(C_d\):
    \[
        P_d^{\mathrm{rel}}
        :=
        C_d\otimes_AG_d.
    \]
    The induced map
    \[
        u_d:P_d^{\mathrm{rel}}\longrightarrow Q_d[-1]
    \]
    remains surjective on the first possible homotopy object. Let
    \(\partial_d:Q_d[-1]\to C_d\) be the boundary and put
    \[
        \chi_d:=\partial_du_d,
        \qquad
        C_{d+1}:=\operatorname{RelCell}_{C_d}(\chi_d).
    \]
    By \Ref{lem:filtered-positive-relation-cell}, the canonical morphism
    \(C_{d+1}\to B\) has relative cofibre \(Q_{d+1}\) pointwise
    \((d+2)\)-connective, and \(C_{d+1}\) is almost perfect as a filtered
    \(C_d\)-module. Since \(C_d\) is almost perfect over \(A\), the filtered
    restriction-transitivity lemma
    \Ref{lem:filtered-almost-perfect-restriction-transitivity} makes
    \(C_{d+1}\) almost perfect over \(A\). The cofibre sequence
    \[
        C_{d+1}\longrightarrow B\longrightarrow Q_{d+1}
    \]
    then makes \(Q_{d+1}\) almost perfect over \(A\).
    \item \textbf{Almost finite presentation.}
    Each relative relation-cell morphism \(C_d\to C_{d+1}\) is compact as a
    commutative filtered algebra by
    \Ref{lem:filtered-positive-relation-cell}. Compact algebra morphisms are
    stable under composition by
    \Ref{lem:compact-algebra-morphism-composition}. Hence every \(C_d\) is
    compact as a filtered \(A\)-algebra. The relative cotangent complex of each
    relation-cell stage is perfect by
    \Ref{lem:filtered-positive-relation-cell}, and cotangent transitivity
    propagates perfectness to
    \[
        L_{C_d/A}
    \]
    for every \(d\).
    \item \textbf{Cotangent transitivity.}
    Let \(E\) be a pointwise connective \(m\)-truncated filtered commutative
    \(A\)-algebra and choose \(d\geq m\). Since \(Q_d\) is pointwise
    \((d+1)\)-connective, \(C_d\to B\) induces an equivalence on operadic
    Postnikov truncations through degree \(m\). By
    \Ref{prop:filtered-graded-operadic-postnikov},
    \[
    \begin{aligned}
        \Map_{\CAlg(\operatorname{Fil}(\Sp))_{A/}}(B,E)
        &\simeq
        \Map_{\CAlg(\operatorname{Fil}(\Sp))_{A/}}(\tau_{\leq m}^{\CAlg}B,E)\\
        &\simeq
        \Map_{\CAlg(\operatorname{Fil}(\Sp))_{A/}}(\tau_{\leq m}^{\CAlg}C_d,E)\\
        &\simeq
        \Map_{\CAlg(\operatorname{Fil}(\Sp))_{A/}}(C_d,E).
    \end{aligned}
    \]
    The right-hand mapping functor preserves filtered colimits because \(C_d\) is
    compact. Since \(m\) was arbitrary, \(A\to B\) is almost finitely presented
    in the sense of
    \Ref{conv:filtered-graded-almost-finite-presentation}.
    \item \textbf{Perfect cotangent approximants.}
    It remains to prove that the relative cotangent complex is almost perfect.
    For every \(d\geq1\), cotangent transitivity for
    \[
        A\longrightarrow C_d\longrightarrow B
    \]
    gives a cofibre sequence
    \[
        B\otimes_{C_d}L_{C_d/A}
        \longrightarrow
        L_{B/A}
        \longrightarrow
        L_{B/C_d}.
    \]
    Put
    \[
        P_d^{\mathrm{cot}}
        :=
        B\otimes_{C_d}L_{C_d/A}.
    \]
    Since \(L_{C_d/A}\) is perfect, \(P_d^{\mathrm{cot}}\) is a perfect
    filtered \(B\)-module.
    \item \textbf{Hurewicz estimate.}
    For \(d\geq1\), the morphism \(C_d\to B\) is pointwise an isomorphism on
    \(\pi_0\), and its relative cofibre \(Q_d\) is pointwise
    \((d+1)\)-connective. Hence
    \Ref{lem:filtered-relative-hurewicz-first-degree} makes
    \[
        L_{B/C_d}
    \]
    pointwise \((d+1)\)-connective. Consequently, for every fixed integer
    \(m\geq0\) and every \(d\geq\max\{m,1\}\), the preceding transitivity
    sequence induces an equivalence
    \[
        \tau_{\leq m}P_d^{\mathrm{cot}}
        \xrightarrow{\simeq}
        \tau_{\leq m}L_{B/A}.
    \]
    \item \textbf{Truncation compactness.}
    The module \(P_d^{\mathrm{cot}}\) is perfect and therefore compact in the
    filtered \(B\)-module category. Its \(m\)-truncation is consequently compact
    in the full subcategory of pointwise \(m\)-truncated filtered \(B\)-modules:
    for every filtered diagram \(\{E_\alpha\}\) in that subcategory,
    \[
    \begin{aligned}
        \Map
        \left(
            \tau_{\leq m}P_d^{\mathrm{cot}},
            \operatorname*{colim}_\alpha E_\alpha
        \right)
        &\simeq
        \Map
        \left(
            P_d^{\mathrm{cot}},
            \operatorname*{colim}_\alpha E_\alpha
        \right)
        \\
        &\simeq
        \operatorname*{colim}_\alpha
        \Map(P_d^{\mathrm{cot}},E_\alpha)
        \\
        &\simeq
        \operatorname*{colim}_\alpha
        \Map(\tau_{\leq m}P_d^{\mathrm{cot}},E_\alpha).
    \end{aligned}
    \]
    Therefore
    \[
        \tau_{\leq m}L_{B/A}
    \]
    is compact in the pointwise \(m\)-truncated filtered module category for every
    \(m\).
    \item \textbf{Connectivity of the shifted cotangent complex.}
    Finally, the pointwise-surjective case of
    \Ref{lem:filtered-relative-hurewicz-first-degree} makes
    \(L_{B/A}\) pointwise \(1\)-connective. Thus \(L_{B/A}\) is bounded below,
    and the preceding truncation-compactness statement proves directly that
    \(L_{B/A}\) is almost perfect. Shifting by \([-1]\) shows that the relative
    shifted cotangent complex is pointwise connective.
\end{enumerate}
\end{proof}


\begin{lemma}[Cotangent completeness for the canonical adic filtration]
\label{lem:cotangent-completion-adic}
Let \(A\to B\) be connective and surjective on \(\pi_0\), and suppose
\[
    M:=L_{B/A}[-1]
\]
is connective and almost perfect. Then the shifted cotangent complex of the completed
canonical adic filtration is a connective almost-perfect complete filtered
\(B\)-module, and completion induces an equivalence
\[
    L_{B/\operatorname{adic}(A\to B)}[-1]
    \xrightarrow{\simeq}
    L_{B/\widehat{\operatorname{adic}}(A\to B)}[-1].
\]
\end{lemma}

\begin{proof}
\leavevmode
\begin{enumerate}[label=\textup{(\arabic*)}]
    \item \textbf{Completed filtered algebra.}
    Put
    \[
        F^\bullet C:=\widehat{\operatorname{adic}}(A\to B).
    \]
    By \Ref{thm:adic-associated-graded},
    \[
        \operatorname{Gr}C
        \simeq
        \operatorname{Sym}^{\mathrm{gr}}_B(M\langle1\rangle).
    \]
    The associated graded is connective. Since \(F^\bullet C\) is complete,
    \Ref{lem:complete-positive-connectivity}, applied to every shifted tail, shows
    that \(F^\bullet C\) is pointwise connective. Its augmentation to the filtered
    algebra \(B\) concentrated in filtration degree zero is therefore pointwise surjective on
    \(\pi_0\).
    \item \textbf{Almost-perfect augmentation module.}
    Regard \(B\) as a filtered \(C\)-module concentrated in filtration degree zero. It is
    complete, and its associated graded is the augmentation module over
    \(\operatorname{Sym}^{\mathrm{gr}}_B(M\langle1\rangle)\). By
    \Ref{lem:positive-graded-symmetric-augmentation-module}, that graded module is
    almost perfect. The filtered detection theorem
    \Ref{thm:filtered-almost-perfect-detection} therefore gives
    \[
        B\in\operatorname{APerf}(F^\bullet C).
    \]
    Apply \Ref{lem:filtered-surjective-almost-finite} to the filtered augmentation
    \(F^\bullet C\to B\). Its shifted cotangent complex
    \[
        L_{B/F^\bullet C}[-1]
    \]
    is connective and almost perfect. Because the base filtered algebra is
    complete, \Ref{lem:almost-perfect-filtered-complete} makes this shifted cotangent complex
    complete.
    \item \textbf{Define the cotangent comparison.}
    Put
    \[
        C_0^\bullet:=\operatorname{adic}(A\to B),
        \qquad
        C_1^\bullet:=\widehat{\operatorname{adic}}(A\to B),
    \]
    and let \(q:C_0^\bullet\to C_1^\bullet\) be the completion morphism. The
    induced base-change map gives the canonical cotangent comparison
    \[
        \kappa:
        L_{B/C_0^\bullet}[-1]
        \longrightarrow
        L_{B/C_1^\bullet}[-1].
    \]
    \item \textbf{Apply naturality of graded cotangent compatibility.}
    Naturality of the first equivalence in
    \Ref{lem:cotangent-graded-adic-compatibility} gives the commutative diagram
    \[
    \begin{tikzcd}[column sep=huge,row sep=large]
        \operatorname{Gr}L_{B/C_0^\bullet}[-1]
            \arrow[r,"\operatorname{Gr}(\kappa)"]
            \arrow[d,"\simeq"']
        &
        \operatorname{Gr}L_{B/C_1^\bullet}[-1]
            \arrow[d,"\simeq"]
        \\
        L_{B/\operatorname{Gr}C_0}[-1]
            \arrow[r]
        &
        L_{B/\operatorname{Gr}C_1}[-1].
    \end{tikzcd}
    \]
    \item \textbf{Identify the associated-graded comparison.}
    Completion preserves associated graded, so
    \(\operatorname{Gr}(q)\) is an equivalence. Cotangent functoriality
    therefore makes the lower horizontal morphism, and hence
    \(\operatorname{Gr}(\kappa)\), an equivalence.
    \item \textbf{Identify the source before completion.}
    Before completion,
    \Ref{lem:cotangent-graded-adic-compatibility} also identifies
    \[
        L_{B/C_0^\bullet}[-1]
        \simeq
        \operatorname{wt}_1(M),
    \]
    which is complete.
    \item \textbf{Use completeness of the target.}
    The target is complete by the preceding step.
    \item \textbf{Apply complete filtered detection.}
    Thus
    \(\kappa\) is an associated-graded equivalence between complete filtered
    modules, and \Ref{lem:complete-filtered-detection} makes it an equivalence.
\end{enumerate}
\end{proof}


This completes the technical package needed later: finite-cell approximation controls
algebra and module finiteness, associated graded detects complete filtered objects, and
filtered Hurewicz comparison keeps the shifted cotangent classifier connective and
almost perfect after completion. We can now define the nonlinear fixed-centre domain on
which formal-order reconstruction converges.

\section{The complete fixed-centre reconstruction domain}
\label{sec:complete-formal-augmentations}


The preceding estimates isolate the desired nonlinear domain. A complete almost finitely
presented augmentation is a fixed point of canonical adic completion; on the connective,
surjective, almost-finitely-presented sector, that fixed-point condition will be
identified with derived completeness for the degree-zero augmentation ideal.

Two facts are needed before the comonadic theorem. First, canonical adic completion must
be identified with its Amitsur and ideal-theoretic descriptions. Second, structured
square-zero augmentations generated by connective almost-perfect cotangent data must
already lie in the complete domain. Together these facts allow the cotangent and
square-zero functors to restrict to one reconstruction category.

\subsection{Canonical completeness and the complete almost finitely presented sector}

Let \(A\to B\) be a connective augmented arrow surjective on \(\pi_0\), and write
\[
    I:=\ker\bigl(\pi_0(A)\to\pi_0(B)\bigr).
\]
Derived \(I\)-completeness is used in the sense of
\cite[Sections~7.3.3--7.3.4]{LurieSAG}. It gives the intrinsic fixed-point
description of the canonical adic completion.

\begin{definition}[Complete almost finitely presented augmentation]
\label{def:complete-almost-finite-augmentation}
Fix a morphism \(R\to B\) of connective commutative ring spectra. Define
\[
    \CAlg(\Sp)^{\mathrm{cn},\mathrm{afp},\wedge}_{R//B}
    \subseteq
    \CAlg(\Sp)^{\mathrm{cn}}_{R//B}
\]
to be the full subcategory spanned by augmented arrows \(A\to B\) such that:
\begin{enumerate}[label=(\roman*)]
    \item \(\pi_0(A)\to\pi_0(B)\) is surjective;
    \item \(B\) is almost perfect as an \(A\)-module;
    \item the canonical adic completion morphism of
    \Ref{def:canonical-adic-filtration} is an equivalence
    \[
        A\xrightarrow{\simeq}A_B^\wedge.
    \]
\end{enumerate}
Objects of this full subcategory are \textbf{complete almost finitely presented
augmentations} over \(R\) with centre \(B\). Completeness is a fixed-point
property of a functorially constructed completion.

The fixed-point condition is compatible with the augmentation:
\[
\begin{tikzcd}[column sep=huge,row sep=large]
    A \arrow[r] \arrow[dr]
        & A_B^\wedge \arrow[d]
    \\
        & B .
\end{tikzcd}
\]
\end{definition}

\begin{lemma}[Nilpotent modules are derived complete]
\label{lem:nilpotent-modules-derived-complete}
Let \(A\) be a connective commutative ring spectrum and let
\(I\subseteq\pi_0(A)\) be a finitely generated ideal. If an \(A\)-module
\(N\) is annihilated by a power \(I^m\), then \(N\) is derived
\(I\)-complete. The full subcategory of derived \(I\)-complete \(A\)-modules
is stable and closed under limits.
\end{lemma}

\begin{proof}
\leavevmode
\begin{enumerate}[label=\textup{(\arabic*)}]
    \item \textbf{Use nilpotence on homotopy modules.}
    Every homotopy module \(\pi_k(N)\) is annihilated by \(I^m\).
    \item \textbf{Identify ordinary adic completeness.}
    Its ordinary
    \(I\)-adic tower is therefore eventually constant, so it is \(I\)-adically
    complete and hence \(I\)-complete.
    \item \textbf{Apply the spectral completion criterion.}
    The homotopy-group criterion for spectral
    derived completion,
    \cite[Theorem~7.3.4.1]{LurieSAG}, implies that \(N\) is derived
    \(I\)-complete.
    \item \textbf{Record stability under limits.}
    Stability and closure under limits are the formal properties
    of the complete subcategory in the completion theory of
    \cite[Section~7.3.3]{LurieSAG}.
\end{enumerate}
\end{proof}

\begin{lemma}[A complete filtration gives derived ideal completeness]
\label{lem:filtered-completeness-augmentation-kernel}
Let \(F^\bullet C\) be a complete pointwise connective filtered commutative
ring spectrum, put
\[
    A:=F^0C,
    \qquad
    B:=\operatorname{Gr}^0C,
    \qquad
    I:=\ker\bigl(\pi_0(A)\to\pi_0(B)\bigr),
\]
and suppose \(I\) is finitely generated. If \(F^\bullet M\) is a complete
filtered \(C\)-module, then \(F^0M\) is derived \(I\)-complete. In
particular, \(A\) is derived \(I\)-complete.
\end{lemma}

\begin{proof}
\leavevmode
\begin{enumerate}[label=\textup{(\arabic*)}]
    \item \textbf{Express filtration zero as an inverse limit.}
    Completeness gives
    \[
        F^0M\simeq
        \varprojlim\nolimits_n\operatorname{cofib}(F^nM\to F^0M).
    \]
    \item \textbf{Identify the augmentation ideal.}
    Since \(F^1C\to F^0C\to B\) is a cofibre sequence of connective spectra, the
    ideal \(I\) is the image of \(\pi_0F^1C\to\pi_0F^0C\).
    \item \textbf{Use multiplicativity of the filtration.}
    Multiplicativity of
    the filtration sends a product of \(n\) lifts from \(F^1C\) into \(F^nC\),
    and the action of \(F^nC\) on \(F^0M\) lands in \(F^nM\). Hence \(I^n\)
    acts trivially on
    \[
        \operatorname{cofib}(F^nM\to F^0M).
    \]
    \item \textbf{Complete the finite quotients.}
    Each finite filtration quotient is therefore derived \(I\)-complete by
    \Ref{lem:nilpotent-modules-derived-complete}.
    \item \textbf{Pass completeness through the inverse limit.}
    The same lemma shows that the
    derived \(I\)-complete subcategory is closed under limits.
    \item \textbf{Conclude derived completeness.}
    Taking the
    displayed inverse limit proves the claim.
\end{enumerate}
\end{proof}


\begin{theorem}[Hodge--Adams interpretation of canonical adic completion]
\label{thm:adic-adams-completion}
Let \(A\to B\) be a morphism of connective commutative ring spectra which is
surjective on \(\pi_0\). Then there is a canonical equivalence
\[
    A_B^\wedge
    \simeq
    \operatorname*{Tot}_{[n]\in\Delta}
    B^{\otimes_A(n+1)}.
\]
More generally, for every bounded-below \(A\)-module \(N\), filtration zero
of the completion of
\[
    N\otimes_A\operatorname{adic}(A\to B)
\]
is canonically equivalent to the Amitsur totalization
\[
    \operatorname*{Tot}_{[n]\in\Delta}
    \bigl(N\otimes_A B^{\otimes_A(n+1)}\bigr).
\]
The comparison uses only connectivity and surjectivity on \(\pi_0\).
\end{theorem}

\begin{proof}
\leavevmode
\begin{enumerate}[label=\textup{(\arabic*)}]
    \item \textbf{Identify with the Hodge filtration in the literature.}
    The universal property used in
    \Ref{thm:adic-associated-graded} identifies
    \(\operatorname{adic}(A\to B)\) with the relative
    \(\mathbb E_\infty\)-Hodge filtration attached to the augmentation
    \(A\to B\). Through
    \Ref{conv:nonnegative-filtered-indexing}, the completion is therefore the
    complete relative Hodge filtration considered in
    \cite[Construction~2.1 and Lemma~2.2]{AntieauFiltrationsCohomologyIII2025}.
    \item \textbf{Apply the Hodge--Adams comparison.}
    For a connective morphism \(A\to B\) that is surjective on \(\pi_0\),
    \cite[Proposition~2.17]{AntieauFiltrationsCohomologyIII2025} identifies the
    completed \(N\)-valued relative Hodge filtration with a complete
    filtration on the \(B\)-Adams completion
    \[
        N_B^\wedge
        :=
        \operatorname*{Tot}_{[n]\in\Delta}
        \bigl(N\otimes_A B^{\otimes_A(n+1)}\bigr)
    \]
    for every bounded-below \(A\)-module \(N\). Its filtration-zero term is
    therefore the displayed Amitsur totalization. Setting \(N=A\) gives the
    canonical equivalence for \(A_B^\wedge\).
\end{enumerate}
\end{proof}

\begin{theorem}[Ideal-theoretic interpretation of canonical adic completion]
\label{thm:adic-adams-derived-completion}
Let \(A\to B\) be a morphism of connective commutative ring spectra which is
surjective on \(\pi_0\) and almost finitely presented. Put
\[
    I:=\ker\bigl(\pi_0(A)\to\pi_0(B)\bigr).
\]
Then \(I\) is finitely generated and, for every bounded-below \(A\)-module
\(N\), there is a canonical equivalence
\[
    N_I^\wedge
    \simeq
    \operatorname*{Tot}_{[n]\in\Delta}
    \bigl(N\otimes_A B^{\otimes_A(n+1)}\bigr).
\]
In particular,
\[
    A_B^\wedge
    \simeq
    A_I^\wedge,
\]
where \(A_I^\wedge\) denotes derived \(I\)-completion in spectral algebraic
geometry. Consequently
\[
    A\longrightarrow A_B^\wedge
\]
is an equivalence if and only if \(A\) is derived \(I\)-complete.
\end{theorem}

\begin{proof}
\leavevmode
\begin{enumerate}[label=\textup{(\arabic*)}]
    \item \textbf{Finite augmentation ideal.}
    Almost finite presentation and surjectivity on \(\pi_0\) give
    \[
        \pi_0(B)\simeq\pi_0(A)/I
    \]
    with \(I\) finitely generated. Put
    \[
        X:=\Spec(A),
        \qquad
        Z:=\Spec(B).
    \]
    Then \(Z\to X\) is an almost finitely presented closed immersion with
    underlying closed support \(V(I)\).
    \item \textbf{Čech descent on the formal completion.}
    Let \(\widehat X_Z\) be the formal completion of \(X\) along this support.
    Formal-completion descent for an almost finitely presented closed immersion
    identifies bounded-below quasi-coherent modules on \(\widehat X_Z\) with the
    totalization over the Čech nerve of \(Z\to X\); see
    \cite[Corollary~3.1.6]{HalpernLeistnerPreygel2023}. The \(n\)-simplices of
    that Čech nerve are
    \[
        Z^{\times_X(n+1)}
        \simeq
        \Spec\bigl(B^{\otimes_A(n+1)}\bigr),
    \]
    so the restriction of the quasi-coherent module associated with \(N\) has
    global sections
    \[
        \operatorname*{Tot}_{[n]\in\Delta}
        \bigl(N\otimes_A B^{\otimes_A(n+1)}\bigr).
    \]
    The descent geometry can be displayed as
    \[
    \begin{tikzcd}[column sep=large,row sep=large]
        \cdots
            \arrow[r,shift left=1.1ex]
            \arrow[r,shift right=1.1ex]
            &
        \Spec(B\otimes_A B)
            \arrow[r,shift left=1.1ex]
            \arrow[r,shift right=1.1ex]
            &
        \Spec(B)
            \arrow[r]
            &
        \widehat X_Z .
    \end{tikzcd}
    \]
    \item \textbf{Derived completion description.}
    Since \(I\) is finitely generated, the same formal completion is
    \[
        \widehat X_Z
        \simeq
        \operatorname{Spf}(A,I).
    \]
    The affine formal-completion theorem
    \cite[Corollary~8.2.4.15]{LurieSAG} identifies quasi-coherent modules on
    \(\operatorname{Spf}(A,I)\) with derived \(I\)-complete \(A\)-modules. Under
    this equivalence, the restriction of \(N\) corresponds to \(N_I^\wedge\).
    Thus the two descriptions of the same formal restriction yield
    \[
        N_I^\wedge
        \simeq
        \operatorname*{Tot}_{[n]\in\Delta}
        \bigl(N\otimes_A B^{\otimes_A(n+1)}\bigr).
    \]
    \item \textbf{Compare the two descriptions.}
    Taking \(N=A\) and using
    \Ref{thm:adic-adams-completion} gives
    \[
        A_B^\wedge\simeq A_I^\wedge.
    \]
    The fixed points of derived \(I\)-completion are precisely the derived
    \(I\)-complete \(A\)-modules, which proves the final assertion. On the
    complete almost finitely presented theorem domain, this is the same
    completion criterion recorded in
    \cite[Proposition~2.18]{BrantnerMagidsonNuiten}.
\end{enumerate}
\end{proof}


\begin{proposition}[Intrinsic characterization of complete almost finitely presented augmentations]
\label{prop:complete-aft-intrinsic-characterization}
Let \(A\to B\) be connective and surjective on \(\pi_0\), and put
\(I=\ker(\pi_0(A)\to\pi_0(B))\). The following are equivalent:
\begin{enumerate}[label=(\roman*)]
    \item \(A\to B\) belongs to
    \(\CAlg(\Sp)^{\mathrm{cn},\mathrm{afp},\wedge}_{R//B}\);
    \item \(B\) is almost perfect over \(A\) and \(A\) is derived
    \(I\)-complete;
    \item \(A\to B\) is almost finitely presented and \(A\) is derived
    \(I\)-complete.
\end{enumerate}
Under these conditions \(I\) is finitely generated and
\[
    \operatorname{Gr}\widehat{\operatorname{adic}}(A\to B)
    \simeq
    \operatorname{Sym}^{\mathrm{gr}}_B
    \bigl(L_{B/A}[-1]\langle1\rangle\bigr).
\]
\end{proposition}

\begin{proof}
\leavevmode
\begin{enumerate}[label=\textup{(\arabic*)}]
    \item \textbf{Compare the finiteness conditions.}
    The equivalence of \textup{(ii)} and \textup{(iii)} is
    \Ref{thm:spectral-almost-finite-criterion}.
    \item \textbf{Deduce finite generation of the ideal.}
    Under either condition the
    augmentation is almost finitely presented, so \(I\) is finitely generated.
    \item \textbf{Compare the two completions.}
    By \Ref{thm:adic-adams-derived-completion},
    \[
        A_B^\wedge\simeq A_I^\wedge.
    \]
    \item \textbf{Identify the fixed-point condition.}
    Thus the fixed-point condition in \textup{(i)} is exactly derived
    \(I\)-completeness, proving the equivalence of all three conditions.
    \item \textbf{Record the associated graded.}
    The
    associated-graded formula is \Ref{thm:adic-associated-graded}.
\end{enumerate}
\end{proof}

\subsection{Structured square-zero finiteness and coherent refinement}

\begin{lemma}[Structured square-zero finiteness and coherent refinement]
\label{lem:square-zero-almost-finite}
Let \(R\to B\) be a morphism of connective commutative ring spectra, let
\(M\) be a connective almost-perfect \(B\)-module, and let
\[
    B_\alpha\longrightarrow B
\]
be the structured square-zero extension classified by
\[
    \alpha:L_{B/R}[-1]\longrightarrow M.
\]
Then \(B\) is almost perfect as a \(B_\alpha\)-module and
\(B_\alpha\to B\) is almost finitely presented. If \(B\) is coherent, then
\[
    I_\alpha:=\ker\bigl(\pi_0(B_\alpha)\to\pi_0(B)\bigr)
\]
is finitely presented as a \(\pi_0(B_\alpha)\)-module and \(B_\alpha\) is a
coherent connective commutative ring spectrum.
\end{lemma}

\begin{proof}
\leavevmode
\begin{enumerate}[label=\textup{(\arabic*)}]
    \item \textbf{Almost finite presentation.}
    The spectral square-zero finiteness theorem
    \cite[Corollary~5.2.2.5]{LurieSAG} says that a structured square-zero extension by a
    connective almost-perfect coefficient is almost finitely presented. Since the
    map is surjective on \(\pi_0\),
    \Ref{thm:spectral-almost-finite-criterion} then gives
    \[
        B\in\operatorname{APerf}(B_\alpha).
    \]
    \item \textbf{Degree-zero kernel.}
    Assume now that \(B\) is coherent. Put
    \[
        A_0:=\pi_0(B_\alpha),
        \qquad
        B_0:=\pi_0(B).
    \]
    The long exact sequence of the square-zero fibre sequence gives
    \[
        \pi_1(B)\longrightarrow\pi_0(M)
        \longrightarrow I_\alpha\longrightarrow0.
    \]
    By \Ref{prop:coherent-almost-perfect-homotopy},
    \(\pi_1(B)\) and \(\pi_0(M)\) are finitely presented over the coherent
    ring \(B_0\). The image of
    \(\pi_1(B)\to\pi_0(M)\) is therefore finitely generated, hence finitely
    presented over \(B_0\); its cokernel \(I_\alpha\) is consequently finitely
    presented over \(B_0\).

    The degree-zero kernel of a square-zero extension is square-zero, so
    \(I_\alpha^2=0\) and
    \[
        B_0\simeq A_0/I_\alpha.
    \]
    A finite \(B_0\)-generating set of \(I_\alpha\) is also a finite
    \(A_0\)-generating set. Applying
    \Ref{lem:finite-presentation-quotient-comparison} to
    \(A_0\twoheadrightarrow B_0\) now makes \(I_\alpha\) finitely presented
    over \(A_0\).
    \item \textbf{Spectral coherence.}
    Apply the square-zero case of
    \Ref{lem:coherence-nilpotent-extension} to
    \(A_0\twoheadrightarrow B_0\). It makes \(A_0\) coherent. Moreover every
    homotopy module of \(B\) and \(M\) is finitely presented over \(B_0\);
    since \(A_0\twoheadrightarrow B_0\) has finitely generated kernel, these
    modules are finitely presented over \(A_0\) by
    \Ref{lem:finite-presentation-quotient-comparison}. The long exact homotopy
    sequence of
    \[
        M\longrightarrow B_\alpha\longrightarrow B
    \]
    then expresses each \(\pi_n(B_\alpha)\) through kernels, cokernels, and
    extensions of finitely presented \(A_0\)-modules. Coherence of \(A_0\)
    keeps these modules finitely presented, proving spectral coherence.
\end{enumerate}
\end{proof}

\begin{corollary}[Structured square-zero objects lie in the complete almost finitely presented sector]
\label{cor:square-zero-complete-aft}
Fix a morphism \(R\to B\) of connective commutative ring spectra. Let \(M\) be connective and almost
perfect, let
\[
    \alpha:L_{B/R}[-1]\longrightarrow M,
\]
and let \(B_\alpha\to B\) be the corresponding structured square-zero
extension over \(R\). Then
\[
    B_\alpha\to B
    \in
    \CAlg(\Sp)^{\mathrm{cn},\mathrm{afp},\wedge}_{R//B}.
\]
\end{corollary}

\begin{proof}
\leavevmode
\begin{enumerate}[label=\textup{(\arabic*)}]
    \item \textbf{Establish finiteness of the augmentation.}
    By \Ref{lem:square-zero-almost-finite}, \(B\) is almost perfect over
    \(B_\alpha\), and the augmentation is almost finitely presented.
    \item \textbf{Deduce finite generation of the ideal.}
    Hence its
    degree-zero augmentation ideal \(I_\alpha\) is finitely generated.
    \item \textbf{Use square-zero nilpotence.}
    The
    square-zero construction gives \(I_\alpha^2=0\); hence \(I_\alpha^2\)
    annihilates the underlying \(B_\alpha\)-module \(B_\alpha\).
    \item \textbf{Apply derived completeness.}
    Therefore
    \(B_\alpha\) is derived \(I_\alpha\)-complete by
    \Ref{lem:nilpotent-modules-derived-complete}.
    \item \textbf{Invoke the intrinsic characterization.}
    Apply
    \Ref{prop:complete-aft-intrinsic-characterization}.
\end{enumerate}
\end{proof}


We now have the theorem domain and know that structured square-zero formation preserves
it. Completeness comes from the canonical adic filtration, finiteness from the cotangent
criterion, and the ideal-theoretic comparison makes the fixed-point condition intrinsic.
These are exactly the inputs required for the fixed-centre cotangent--square-zero
theorem.

\section{Comonadic reconstruction from anchored shifted-cotangent data}
\label{sec:cotangent-square-zero-comonadic}


Fix a connective centre \(B\) over \(R\). The principal theorem compares two
constructions: an augmentation \(A\to B\) produces the anchored shifted cotangent object
\[
    L_{B/R}[-1]\longrightarrow L_{B/A}[-1],
\]
whereas an anchored map \(L_{B/R}[-1]\to M\) produces a structured square-zero
augmentation of \(B\). These assignments are adjoint. The proof of comonadicity has
three stages: restrict the adjunction to the complete theorem domain, prove
conservativity by associated-graded detection, and verify the split totalizations
required by Barr--Beck through filtered reconstruction.

The adjunction unit already exhibits the resulting coalgebra structure:
\[
\begin{tikzcd}[column sep=huge,row sep=large]
    A
        \arrow[r,"\eta_A"]
    &
    \operatorname{sqz}_{B/R}\operatorname{cot}_{B/R}(A)
    \\
    \operatorname{cot}_{B/R}(A)
        \arrow[r,"{\operatorname{cot}_{B/R}(\eta_A)}"']
    &
    \mathbb C^{\mathrm{sqz}}_{B/R}
        \operatorname{cot}_{B/R}(A).
\end{tikzcd}
\]
The lower map is therefore generated by the adjunction rather than imposed
independently. Comonadicity will show that this anchored cotangent coalgebra determines
the complete augmentation. Coherence of \(B\) is unnecessary here; it enters only when
the complete theory is restricted to bounded coherent Artinian tests.

\subsection{Anchored shifted-cotangent data}

\begin{definition}[Anchored shifted-cotangent category]
\label{def:anchored-shifted-cotangent-category}
Fix a morphism of connective commutative ring spectra
\[
    R\longrightarrow B.
\]
Define
\[
    \operatorname{AnchCot}^{\mathrm{sp}}_{B/R}
    :=
    (\Mod_B)_{L_{B/R}[-1]/}.
\]
Thus an object is a pair \((M,\alpha)\) with
\[
    \alpha:L_{B/R}[-1]\longrightarrow M.
\]
Let
\[
    \operatorname{AnchCot}^{\mathrm{aperf}}_{B/R,\geq0}
\]
be the full subcategory in which \(M\) is connective and almost perfect.
\end{definition}

\begin{definition}[Shifted-cotangent and structured square-zero functors]
\label{def:cotangent-square-zero-functors}
For an augmented \(R\)-algebra \(A\to B\), define
\[
    \operatorname{cot}_{B/R}(A)
    :=
    \left(
        L_{B/R}[-1]
        \longrightarrow
        L_{B/A}[-1]
    \right),
\]
where the anchor is induced by cotangent transitivity.

For \((M,\alpha)\in\operatorname{AnchCot}^{\mathrm{sp}}_{B/R}\), let
\[
    B\xrightarrow{d_\alpha}B\oplus M[1]
\]
be the section classified by the degree-one derivation
\(L_{B/R}\to M[1]\). Define
\[
    B_\alpha:=B\times_{B\oplus M[1]}B
\]

so that
\[
\begin{tikzcd}[column sep=large,row sep=large]
    B_\alpha \arrow[r] \arrow[d]
        & B \arrow[d,"d_\alpha"]
    \\
    B \arrow[r,"d_0"']
        & B\oplus M[1]
\end{tikzcd}
\]
is Cartesian, where \(d_0\) is the zero derivation. Set
\[
    \operatorname{sqz}_{B/R}(M,\alpha)
    :=
    (B_\alpha\to B).
\]
This structured square-zero functor is defined in unrestricted commutative ring spectra;
it restricts to connective arrows on connective coefficients by
\Ref{prop:structured-square-zero-fibre}.
\end{definition}

\begin{proposition}[Shifted-cotangent--structured-square-zero adjunction]
\label{prop:cotangent-square-zero-adjunction}
There is a canonical adjunction
\[
    \operatorname{cot}_{B/R}
    \dashv
    \operatorname{sqz}_{B/R}.
\]
More explicitly,
\[
\Map_{\operatorname{AnchCot}^{\mathrm{sp}}_{B/R}}
\bigl(
    \operatorname{cot}_{B/R}(A),(M,\alpha)
\bigr)
\simeq
\Map_{\CAlg(\Sp)_{R//B}}
(A,B_\alpha).
\]

A point of the right-hand mapping space is equivalently a commutative lifting
diagram
\[
\begin{tikzcd}[column sep=large,row sep=large]
    A \arrow[r,dashed] \arrow[d]
        & B_\alpha \arrow[d]
    \\
    B \arrow[r,equal]
        & B .
\end{tikzcd}
\]
\end{proposition}

\begin{proof}
\leavevmode
\begin{enumerate}[label=\textup{(\arabic*)}]
    \item \textbf{Interpret the right-hand mapping space.}
    A point of the right-hand mapping space is a lift of the augmentation
    \(A\to B\) through the structured square-zero extension classified by
    \(L_{B/R}\to M[1]\).
    \item \textbf{Apply the square-zero lifting classifier.}
    By
    \Ref{prop:formal-lifting-fixed-square-zero}, such lifts are equivalent to maps
    \[
        L_{B/A}\longrightarrow M[1]
    \]
    whose restriction along \(L_{B/R}\to L_{B/A}\) is the specified classifying
    term.
    \item \textbf{Shift to anchored cotangent data.}
    Shifting by \([-1]\) identifies this exactly with a morphism from
    \(\operatorname{cot}_{B/R}(A)\) to \((M,\alpha)\) in the undercategory of
    \(L_{B/R}[-1]\).
\end{enumerate}
\end{proof}

\subsection{Filtered reconstruction}

\begin{lemma}[Filtered reconstruction of a complete almost finitely presented augmentation]
\label{lem:filtered-reconstruction-complete-aft}
Let \(F^\bullet C\) be a complete filtered commutative \(R\)-algebra whose
associated-graded terms are connective. Suppose its filtration-zero augmentation
is
\[
    A:=F^0C\longrightarrow B:=\operatorname{Gr}^0C
\]
and that for a connective almost-perfect \(B\)-module \(M\) there is an
equivalence
\[
    \operatorname{Gr}C
    \simeq
    \operatorname{Sym}^{\mathrm{gr}}_B(M\langle1\rangle).
\]
Then \(A\to B\) belongs to
\(\CAlg(\Sp)^{\mathrm{cn},\mathrm{afp},\wedge}_{R//B}\). Moreover the filtered augmentation
module \(B\), concentrated in filtration degree zero, is almost perfect over
\(F^\bullet C\).
\end{lemma}

\begin{proof}
\leavevmode
\begin{enumerate}[label=\textup{(\arabic*)}]
    \item \textbf{Connectivity.}
    Apply \Ref{lem:complete-positive-connectivity} to every shifted filtration tail
    of \(F^\bullet C\). Every \(F^nC\), in particular \(A=F^0C\), is
    connective. Since
    \[
        F^1C\longrightarrow F^0C\longrightarrow B
    \]
    is a cofibre sequence with connective first term, \(\pi_0(A)\to\pi_0(B)\) is
    surjective.
    \item \textbf{Almost-perfect augmentation module.}
    Regard \(B\) as the filtered \(C\)-module concentrated in filtration degree
    zero. It is complete. Its associated graded is the augmentation module over
    \(\operatorname{Sym}^{\mathrm{gr}}_B(M\langle1\rangle)\), which is almost perfect by
    \Ref{lem:positive-graded-symmetric-augmentation-module}. Hence
    \Ref{thm:filtered-almost-perfect-detection} makes \(B\) almost perfect as a
    filtered \(C\)-module and, after evaluation at filtration zero, gives
    \[
        B\in\operatorname{APerf}(A).
    \]
    The spectral criterion for almost finite presentation therefore makes the degree-zero
    augmentation ideal finitely generated.
    \item \textbf{Derived completeness.}
    Finally apply
    \Ref{lem:filtered-completeness-augmentation-kernel} to the
    complete filtered algebra \(F^\bullet C\), viewed as a module over itself.
    Its filtration-zero term \(A\) is derived complete for that augmentation
    ideal. The intrinsic characterization
    \Ref{prop:complete-aft-intrinsic-characterization} now gives
    \[
        A\to B\in\CAlg(\Sp)^{\mathrm{cn},\mathrm{afp},\wedge}_{R//B}.
    \]
\end{enumerate}
\end{proof}

\subsection{Comonadicity}

\begin{theorem}[Shifted-cotangent--structured-square-zero comonadicity]
\label{thm:cotangent-square-zero-comonadicity}
Let \(R\to B\) be a morphism of connective commutative ring spectra. The
adjunction of \Ref{prop:cotangent-square-zero-adjunction} restricts to
\[
    \operatorname{cot}_{B/R}
    :
    \CAlg(\Sp)^{\mathrm{cn},\mathrm{afp},\wedge}_{R//B}
    \rightleftarrows
    \operatorname{AnchCot}^{\mathrm{aperf}}_{B/R,\geq0}
    :
    \operatorname{sqz}_{B/R},
\]
and the left adjoint \(\operatorname{cot}_{B/R}\) is comonadic.
\end{theorem}

\begin{proof}
\leavevmode
\begin{enumerate}[label=\textup{(\arabic*)}]
    \item \textbf{Restriction to the theorem domain.}
    The right adjoint lands in the displayed source category by
    \Ref{cor:square-zero-complete-aft}. Conversely, if \(A\to B\) is complete
    almost finitely presented, \Ref{prop:complete-aft-intrinsic-characterization} makes the
    augmentation almost finitely presented. Hence \(L_{B/A}\) is almost perfect,
    and surjectivity on \(\pi_0\) makes \(L_{B/A}[-1]\) connective. Thus the
    shifted-cotangent functor lands in
    \(\operatorname{AnchCot}^{\mathrm{aperf}}_{B/R,\geq0}\).
    This is the spectral form of the theorem domain appearing in
    \cite[Theorem~2.25]{BrantnerMagidsonNuiten}.
    \item \textbf{Conservativity.}
    Suppose a morphism \(A\to A'\) over \(B\) induces an equivalence
    \[
        L_{B/A}[-1]\xrightarrow{\simeq}L_{B/A'}[-1]
    \]
    of anchored shifted cotangent complexes. Functoriality of canonical adic formation gives
    a morphism between the completed filtrations. By
    \Ref{thm:adic-associated-graded}, its associated graded is the free graded
    symmetric algebra on the displayed cotangent equivalence and is therefore an
    equivalence. Both filtered algebras are complete, so
    \Ref{lem:complete-filtered-detection} makes the filtered morphism an
    equivalence. Taking filtration zero and using that both augmented arrows are
    fixed by canonical adic completion gives \(A\simeq A'\).
    \item \textbf{Split totalization data.}
    Let \(A^\bullet\) be a cosimplicial object of
    \(\CAlg(\Sp)^{\mathrm{cn},\mathrm{afp},\wedge}_{R//B}\) which is
    \(\operatorname{cot}_{B/R}\)-split in the Barr--Beck sense: there is an
    augmentation
    \[
        (M,\alpha)
        \longrightarrow
        \operatorname{cot}_{B/R}(A^\bullet)
    \]
    in \(\operatorname{AnchCot}^{\mathrm{aperf}}_{B/R,\geq0}\) making the resulting
    augmented cosimplicial object split. In particular,
    \[
        (M,\alpha)
        \xrightarrow{\simeq}
        \operatorname*{Tot}_{[n]\in\Delta}
        \operatorname{cot}_{B/R}(A^n),
    \]
    and this totalization is absolute: it is preserved by every functor out of
    \(\operatorname{AnchCot}^{\mathrm{aperf}}_{B/R,\geq0}\). After forgetting the
    anchors, the underlying split cosimplicial \(B\)-module therefore has
    absolute totalization
    \[
        M
        \xrightarrow{\simeq}
        \operatorname*{Tot}_{[n]\in\Delta}
        L_{B/A^n}[-1].
    \]
    \item \textbf{Totalization in filtered algebras.}
    Form
    \[
        A^u:=\operatorname*{Tot}_{[n]\in\Delta}A^n
    \]
    in the unrestricted category \(\CAlg(\Sp)_{R//B}\), and form, likewise in
    unrestricted filtered commutative spectra, the complete filtered algebra
    \[
        F^\bullet C
        :=
        \operatorname*{Tot}_{[n]\in\Delta}
        \widehat{\operatorname{adic}}(A^n\to B).
    \]
    Limits preserve complete filtered objects. Evaluation in filtration zero
    commutes with limits, and every \(A^n\) is fixed by completion, so
    \[
        F^0C\simeq A^u.
    \]
    The totalization has been formed in unrestricted commutative ring spectra;
    filtered reconstruction will show that \(A^u\) is connective.
    \item \textbf{Commute associated graded with totalization.}
    Associated graded commutes with this totalization. Indeed each graded piece
    is the cofibre of two filtration evaluations; filtration evaluation preserves
    limits, and in a stable category finite cofibres are finite limits and
    therefore commute with arbitrary limits.
    \item \textbf{Compute the totalized associated graded.}
    Hence
    \[
    \begin{aligned}
        \operatorname{Gr}C
        &\simeq
        \operatorname*{Tot}_{[n]\in\Delta}
        \operatorname{Gr}
        \widehat{\operatorname{adic}}(A^n\to B)
        \\
        &\simeq
        \operatorname*{Tot}_{[n]\in\Delta}
        \operatorname{Sym}^{\mathrm{gr}}_B
        \bigl(L_{B/A^n}[-1]\langle1\rangle\bigr).
    \end{aligned}
    \]
    \item \textbf{Use splitness of the cotangent cosimplicial module.}
    Because the underlying augmented cosimplicial \(B\)-module
    \(L_{B/A^\bullet}[-1]\) is split, its totalization is absolute.
    \item \textbf{Identify the free graded algebra.}
    Therefore the
    graded free commutative-algebra functor preserves this totalization, giving
    \[
    \begin{aligned}
        \operatorname{Gr}C
        &\simeq
        \operatorname{Sym}^{\mathrm{gr}}_B
        \left(
            \left(
            \operatorname*{Tot}_{[n]\in\Delta}
            L_{B/A^n}[-1]
            \right)\langle1\rangle
        \right)
        \\
        &\simeq
        \operatorname{Sym}^{\mathrm{gr}}_B\bigl(M\langle1\rangle\bigr).
    \end{aligned}
    \]
    \item \textbf{Record finiteness of the totalized coefficient.}
    Since
    \[
        (M,\alpha)\in\operatorname{AnchCot}^{\mathrm{aperf}}_{B/R,\geq0},
    \]
    the underlying \(B\)-module \(M\) is connective and almost perfect.
    \item \textbf{Recover the connective totalization.}
    Lemma~\Ref{lem:filtered-reconstruction-complete-aft} now shows that
    \(A^u=F^0C\) is connective and that
    \[
        A^u\to B
    \]
    belongs to the complete almost finitely presented sector. Because the inclusion of
    connective commutative ring spectra is full, this unrestricted totalization
    therefore computes the required totalization in the source category. Write
    \[
        A:=A^u
    \]
    from now on.
    \item \textbf{Construct the shifted-cotangent comparison.}
    Write
    \[
        C_n^\bullet
        :=
        \widehat{\operatorname{adic}}(A^n\to B).
    \]
    The totalization cone defining \(F^\bullet C\) has canonical projection
    morphisms
    \[
        p_n:F^\bullet C\longrightarrow C_n^\bullet
    \]
    of filtered commutative \(R\)-algebras, all compatible with the displayed
    augmentations to \(B\). Cotangent functoriality in the base therefore gives
    natural morphisms
    \[
        L_{B/F^\bullet C}[-1]
        \longrightarrow
        L_{B/C_n^\bullet}[-1].
    \]
    Their cosimplicial compatibility is inherited from the totalization cone.
    By the universal property of the limit they assemble to the canonical
    comparison
    \[
        \gamma:
        L_{B/F^\bullet C}[-1]
        \longrightarrow
        \operatorname*{Tot}_{[n]\in\Delta}
        L_{B/C_n^\bullet}[-1].
    \]
    \item \textbf{Apply associated graded to the comparison.}
    Apply associated graded to \(\gamma\). Since associated graded commutes with
    the totalization, and since
    \Ref{lem:cotangent-graded-adic-compatibility} is natural in the filtered
    algebra, \(\operatorname{Gr}(\gamma)\) sits in the canonical comparison
    diagram
    \[
    \begin{tikzcd}[column sep=huge,row sep=large]
        \operatorname{Gr}L_{B/F^\bullet C}[-1]
            \arrow[r,"\operatorname{Gr}(\gamma)"]
            \arrow[d,"\simeq"']
        &
        \displaystyle
        \operatorname*{Tot}_{[n]\in\Delta}
        \operatorname{Gr}L_{B/C_n^\bullet}[-1]
            \arrow[d,"\simeq"]
        \\
        L_{B/\operatorname{Gr}C}[-1]
            \arrow[r]
        &
        \displaystyle
        \operatorname*{Tot}_{[n]\in\Delta}
        L_{B/\operatorname{Gr}C_n}[-1].
    \end{tikzcd}
    \]
    \item \textbf{Define the stagewise cotangent modules.}
    Put
    \[
        M_n:=L_{B/A^n}[-1].
    \]
    \item \textbf{Compute the associated-graded algebras.}
    By \Ref{thm:adic-associated-graded} and the calculation in the preceding
    step,
    \[
        \operatorname{Gr}C_n
        \simeq
        \operatorname{Sym}^{\mathrm{gr}}_B(M_n\langle1\rangle),
        \qquad
        \operatorname{Gr}C
        \simeq
        \operatorname{Sym}^{\mathrm{gr}}_B(M\langle1\rangle).
    \]
    \item \textbf{Compute the graded relative cotangent complexes.}
    Cotangent transitivity for the zero augmentation of a free graded
    commutative \(B\)-algebra gives canonical equivalences
    \[
        L_{B/\operatorname{Gr}C_n}[-1]
        \simeq
        M_n\langle1\rangle,
        \qquad
        L_{B/\operatorname{Gr}C}[-1]
        \simeq
        M\langle1\rangle.
    \]
    \item \textbf{Identify the lower comparison map.}
    Under these equivalences the lower horizontal morphism is exactly
    \[
        M\langle1\rangle
        \longrightarrow
        \operatorname*{Tot}_{[n]\in\Delta}M_n\langle1\rangle.
    \]
    \item \textbf{Use split totalization.}
    This is an equivalence because the augmented cosimplicial module
    \(M\to M_\bullet\) is split and hence has absolute totalization. Therefore
    \(\operatorname{Gr}(\gamma)\) is an equivalence.
    \item \textbf{Completeness of the comparison.}
    For every \(n\), the shifted cotangent module
    \(M_n=L_{B/A^n}[-1]\) is connective and almost perfect by the theorem-domain
    condition. Hence \Ref{lem:cotangent-completion-adic} makes
    \[
        L_{B/C_n^\bullet}[-1]
    \]
    a complete filtered \(B\)-module. Its totalization, the target of
    \(\gamma\), is therefore complete because complete filtered modules are
    closed under limits.

    For the source, \Ref{lem:filtered-reconstruction-complete-aft} makes the
    augmentation module \(B\) almost perfect over \(F^\bullet C\). The
    augmentation is pointwise surjective on \(\pi_0\), so
    \Ref{lem:filtered-surjective-almost-finite} makes
    \(L_{B/F^\bullet C}[-1]\) connective and almost perfect as a filtered
    \(B\)-module. The coefficient algebra \(B\), concentrated in filtration
    degree zero, is complete; hence
    \Ref{lem:almost-perfect-filtered-complete} makes the source complete.
    Since \(\gamma\) is an associated-graded equivalence between complete
    filtered modules, \Ref{lem:complete-filtered-detection} makes
    \(\gamma\) an equivalence.
    \item \textbf{Filtration-zero comparison.}
    Finally
    \Ref{lem:cotangent-graded-adic-compatibility} identifies filtration zero of
    these shifted cotangent complexes in filtered modules with the ordinary shifted cotangent complexes.
    Thus the underlying \(B\)-modules satisfy
    \[
        \operatorname{cot}_{B/R}(A)
        \xrightarrow{\simeq}
        \operatorname*{Tot}_{[n]\in\Delta}
        \operatorname{cot}_{B/R}(A^n).
    \]
    Naturality of cotangent transitivity identifies this equivalence with the
    totalization of the canonical anchor maps from \(L_{B/R}[-1]\). Hence the
    comparison is an equivalence in the undercategory
    \[
        (\Mod_B)_{L_{B/R}[-1]/}
        =
        \operatorname{AnchCot}^{\mathrm{sp}}_{B/R},
    \]
    and therefore preserves the anchor from \(L_{B/R}[-1]\). Consequently,
    \(\operatorname{cot}_{B/R}\) preserves every
    \(\operatorname{cot}_{B/R}\)-split totalization required by the
    Barr--Beck--Lurie criterion.
    \item \textbf{Barr--Beck conclusion.}
    Together with conservativity, the Barr--Beck--Lurie theorem gives
    comonadicity.
\end{enumerate}
\end{proof}

\begin{definition}[Shifted-cotangent--structured-square-zero comonad]
\label{def:cotangent-square-zero-comonad}
Under the hypotheses of
\Ref{thm:cotangent-square-zero-comonadicity}, define
\[
    \mathbb C^{\mathrm{sqz}}_{B/R}
    :=
    \operatorname{cot}_{B/R}
    \operatorname{sqz}_{B/R}
\]
on \(\operatorname{AnchCot}^{\mathrm{aperf}}_{B/R,\geq0}\), with its canonical
comonad structure induced by the adjunction.
\end{definition}

\begin{corollary}[Comonadic classification of complete formal augmentations]
\label{cor:cotangent-square-zero-coalgebra-classification}
The Barr--Beck comparison functor induces a canonical equivalence
\[
    \CAlg(\Sp)^{\mathrm{cn},\mathrm{afp},\wedge}_{R//B}
    \xrightarrow{\simeq}
    \operatorname{Coalg}_{\mathbb C^{\mathrm{sqz}}_{B/R}}
    \bigl(
        \operatorname{AnchCot}^{\mathrm{aperf}}_{B/R,\geq0}
    \bigr),
\]
where \(\operatorname{Coalg}_{\mathbb C^{\mathrm{sqz}}_{B/R}}(-)\) denotes
the \(\infty\)-category of coalgebras for the displayed comonad. Thus the Eilenberg--Moore category is the precise category of comonadic
shifted-cotangent data used for reconstruction.
\end{corollary}

\begin{proof}
\leavevmode
\begin{enumerate}[label=\textup{(\arabic*)}]
    \item \textbf{Invoke the comonadic comparison.}
    This is the comparison equivalence furnished by the comonadicity statement of
    \Ref{thm:cotangent-square-zero-comonadicity}.
    \item \textbf{Identify the comonad.}
    Its comonad is
    \(\mathbb C^{\mathrm{sqz}}_{B/R}\) by
    \Ref{def:cotangent-square-zero-comonad}.
\end{enumerate}
\end{proof}

\begin{corollary}[Canonical structured-square-zero cobar resolution]
\label{cor:canonical-square-zero-cobar-resolution}
Every
\[
    A\to B
    \in
    \CAlg(\Sp)^{\mathrm{cn},\mathrm{afp},\wedge}_{R//B}
\]
admits a canonical augmented cosimplicial resolution
\[
    A\longrightarrow A^\bullet.
\]
Its first terms are functorially of the form
\[
\begin{tikzcd}[column sep=huge]
    A \arrow[r]
        &
    \operatorname{sqz}_{B/R}\operatorname{cot}_{B/R}(A)
        \arrow[r,shift left=1.2ex]
        \arrow[r,shift right=1.2ex]
        &
    \operatorname{sqz}_{B/R}
    \mathbb C^{\mathrm{sqz}}_{B/R}
    \operatorname{cot}_{B/R}(A)
        \arrow[r,shift left=1.5ex]
        \arrow[r]
        \arrow[r,shift right=1.5ex]
        &
    \cdots ,
\end{tikzcd}
\]
The resolution satisfies:
\begin{enumerate}[label=(\roman*)]
    \item every positive term is a constructed structured square-zero
    extension of \(B\);
    \item \(\operatorname{cot}_{B/R}(A^\bullet)\) is split cosimplicial;
    \item the augmentation exhibits
    \[
        A\xrightarrow{\simeq}\operatorname*{Tot}A^\bullet.
    \]
\end{enumerate}
The resolution is the cobar resolution functorially generated by the
adjunction.
\end{corollary}

\begin{proof}
\leavevmode
\begin{enumerate}[label=\textup{(\arabic*)}]
    \item \textbf{Invoke the canonical cobar resolution.}
    This is the canonical cobar resolution associated to the comonadic adjunction
    of \Ref{thm:cotangent-square-zero-comonadicity}.
    \item \textbf{Identify its positive terms.}
    Positive terms lie in the
    image of the right adjoint \(\operatorname{sqz}_{B/R}\),
    \item \textbf{Identify the splitting codegeneracies.}
    the canonical extra
    codegeneracies split the shifted-cotangent image,
    \item \textbf{Recover the totalization.}
    and comonadicity identifies the
    augmentation with the totalization.
\end{enumerate}
\end{proof}

\begin{remark}[Comonadic interpretation]
\label{rem:formal-augmentations-comonadic-interpretation}
The comonadic adjunction reconstructs the complete almost finitely presented augmentation
category from shifted relative cotangent data. Its cobar construction supplies the
corresponding functorial cosimplicial resolution.
\end{remark}


The fixed-centre problem is now solved on the complete almost-finitely-presented sector:
the nonlinear augmentation is equivalent to its anchored shifted-cotangent object with
the canonical comonadic coherence, and the cobar construction reconstructs it by a
functorial cosimplicial square-zero resolution. The next task is to extract the bounded
finite-test fragment needed for formal-moduli semantics.

\section{Artinian finite tests inside the infinitesimal geometry}
\label{sec:spectral-artinian-extensions}


The complete fixed-centre category contains infinite formal neighbourhoods, whereas
Schlessinger semantics is built from bounded tests. Spectral Artinian extensions provide
this finite fragment by combining nilpotence on \(\pi_0\), almost finite presentation,
and eventual coconnectivity of the relative fibre.

The section proves that these intrinsic conditions are equivalent to finite structured
square-zero generation with coherent coefficients. The relation between the intrinsic
and factorized descriptions is summarized by
\[
\begin{tikzcd}[column sep=huge,row sep=large]
    \text{intrinsic Artinian conditions}
        \arrow[r,"\text{finite generation}"]
        \arrow[d,"\text{complete formal object}"']
    &
    \text{finite coherent square-zero factorization}
        \arrow[d,"\operatorname{Spec}"]
    \\
    \text{complete augmentation at }B
        \arrow[r,"\text{affine realization}"']
    &
    \text{finite infinitesimal test at }\Spec(B).
\end{tikzcd}
\]
The finite factorization is a theorem and remains a proof witness. Affine realization,
by contrast, produces a presentation-free object of the finite infinitesimal geometry
already constructed in Chapter~\Ref{chap:differentiation-spectral-doctrine}. Thus
Artinianity adds bounded coherent control without replacing intrinsic infinitesimal
geometry by a chosen presentation.

\begin{remark}[Relation with the Artinian formal-integration category]
\label{rem:bmn-artinian-comparison}
The intrinsic Artinian conditions below agree with the fixed-centre
formal-integration test category of Brantner--Magidson--Nuiten. Their
finite coherent structured square-zero description is given in
\cite[
    Proposition~4.3 and Corollaries~4.4--4.5
]{BrantnerMagidsonNuiten}. The argument below is written in the present
spectral conventions so that the resulting factorizations can be compared
directly with the finite infinitesimal substitutions of
Chapter~\Ref{chap:differentiation-spectral-doctrine}.
\end{remark}

\subsection{Intrinsic Artinian extensions}

\begin{definition}[Spectral Artinian extension]
\label{def:spectral-artinian-extension}
Fix a morphism of connective commutative ring spectra
\[
    R\longrightarrow B
\]
with \(B\) coherent. A connective augmented \(R\)-algebra
\[
    A\longrightarrow B
\]
is a \textbf{spectral Artinian extension} if:
\begin{enumerate}[label=(\roman*)]
    \item \(\pi_0(A)\to\pi_0(B)\) is surjective with nilpotent kernel;
    \item \(A\to B\) is almost finitely presented;
    \item \(\operatorname{fib}(A\to B)\) is eventually coconnective.
\end{enumerate}
Write
\[
    \operatorname{Art}^{\mathrm{sp}}_{B/R}
    \subseteq
    \CAlg(\Sp)^{\mathrm{cn}}_{R//B}
\]
for the full subcategory on these arrows.
\end{definition}

\begin{lemma}[Finite presentation across a finitely generated quotient]
\label{lem:finite-presentation-quotient-comparison}
Let \(R\twoheadrightarrow S=R/J\) be a surjection of ordinary commutative
rings with \(J\) finitely generated.
\begin{enumerate}[label=(\roman*)]
    \item Every finitely presented \(S\)-module is finitely presented after
    restriction of scalars to \(R\).
    \item If an \(R\)-module \(M\) is annihilated by \(J\) and finitely
    presented over \(R\), then it is finitely presented over \(S\).
\end{enumerate}
\end{lemma}

\begin{proof}
\leavevmode
\begin{enumerate}[label=\textup{(\arabic*)}]
    \item \textbf{Prove finite presentation after restriction.}
    For \textup{(i)}, start with a finite presentation over \(S\), lift its
    finite matrix of relations to \(R\), and adjoin finitely many relations
    expressing the action of a chosen finite generating set of \(J\) as zero on
    the finite set of module generators.
    \item \textbf{Tensor the second presentation with the quotient.}
    For \textup{(ii)}, tensor a finite
    \(R\)-presentation of \(M\) with \(S\).
    \item \textbf{Identify the resulting cokernel.}
    Since \(J\) acts trivially on
    \(M\), the resulting cokernel is still \(M\), giving a finite
    \(S\)-presentation.
\end{enumerate}
\end{proof}

\begin{lemma}[Coherence across a finite nilpotent extension]
\label{lem:coherence-nilpotent-extension}
Let \(A_0\twoheadrightarrow B_0\) be a surjection of ordinary commutative
rings with nilpotent kernel \(I\). Suppose \(B_0\) is coherent and \(I\) is
finitely presented as an \(A_0\)-module. Then:
\begin{enumerate}[label=(\roman*)]
    \item \(A_0\) is coherent;
    \item every power \(I^k\) is finitely presented as an \(A_0\)-module;
    \item for every finitely presented \(A_0\)-module \(M\), each quotient
    \[
        I^kM/I^{k+1}M
    \]
    is finitely presented over \(B_0\).
\end{enumerate}
\end{lemma}

\begin{proof}
\leavevmode
\begin{enumerate}[label=\textup{(\arabic*)}]
    \item \textbf{Square-zero base case.}
    We first treat a square-zero extension
    \[
        C\twoheadrightarrow D=C/J,
        \qquad J^2=0,
    \]
    with \(D\) coherent and \(J\) finitely presented over \(C\). Since \(J\) is
    finitely generated, \(D\) is finitely presented over \(C\), and because \(J\)
    is annihilated by itself the finite-presentation comparison of
    \Ref{lem:finite-presentation-quotient-comparison} makes \(J\) finitely
    presented over \(D\).
    \item \textbf{Choose generators of the ideal.}
    Let \(K=(x_1,\ldots,x_r)\subseteq C\), let
    \[
        \varphi:C^{\oplus r}\twoheadrightarrow K
    \]
    send the standard basis to the chosen generators, and reduce modulo \(J\).
    \item \textbf{Reduce the presentation modulo the square-zero ideal.}
    Write
    \[
        \bar\varphi:D^{\oplus r}\twoheadrightarrow\bar K
    \]
    for the resulting presentation of the image ideal \(\bar K\subseteq D\).
    \item \textbf{Choose lifts of the reduced relations.}
    Since \(D\) is coherent, the kernel
    \[
        L:=\ker(\bar\varphi)
    \]
    is finitely generated. Choose lifts
    \(\widetilde\ell_1,\ldots,\widetilde\ell_s\in C^{\oplus r}\) of a finite
    generating set of \(L\).
    \item \textbf{Prove finite generation of the intersection.}
    We first show that \(K\cap J\) is finitely generated as a \(D\)-module.
    If \(y\in K\cap J\), choose \(v\in C^{\oplus r}\) with
    \(\varphi(v)=y\). Its reduction \(\bar v\) lies in \(L\). Subtracting a
    \(C\)-linear combination of the \(\widetilde\ell_j\), we may arrange that
    the resulting vector lies in \(J^{\oplus r}\). Consequently \(y\) is a
    \(D\)-linear combination of the finitely many elements
    \(\varphi(\widetilde\ell_j)\in J\) together with the image of
    \(J^{\oplus r}\) under \(\varphi\). The latter image is finitely generated
    because \(J\) is a finitely generated \(D\)-module.
    \item \textbf{Promote the intersection to finite presentation.}
    Thus \(K\cap J\) is
    finitely generated over \(D\), hence finitely presented because \(D\) is
    coherent.
    \item \textbf{Compare finite presentations over the extension.}
    The ideal \(\bar K\) is finitely presented over \(D\) by coherence. The
    finite-presentation comparison
    \Ref{lem:finite-presentation-quotient-comparison} carries both
    \(K\cap J\) and \(\bar K\) to finitely presented \(C\)-modules.
    \item \textbf{Conclude coherence of the square-zero extension.}
    Therefore
    the exact sequence
    \[
        0\longrightarrow K\cap J
        \longrightarrow K
        \longrightarrow \bar K
        \longrightarrow0
    \]
    makes \(K\) finitely presented over \(C\). Every finitely generated ideal
    of \(C\) is therefore finitely presented, so \(C\) is coherent.
    \item \textbf{Nilpotent induction.}
    Choose \(m\) with \(I^m=0\) and put
    \[
        R_n:=A_0/I^n,
        \qquad 1\leq n\leq m.
    \]
    Then \(R_1=B_0\), \(R_m=A_0\), and for \(n\geq2\) the kernel of
    \[
        R_n\longrightarrow R_{n-1}
    \]
    is
    \[
        J_n:=I^{n-1}/I^n,
    \]
    which is square-zero. Suppose inductively that \(R_{n-1}\) is coherent.
    Because \(I\) is finitely presented over \(A_0\), base change gives a finitely
    presented \(R_{n-1}\)-module
    \[
        I\otimes_{A_0}R_{n-1}
        \simeq
        I/I^n.
    \]
    The submodule
    \[
        J_n=I^{n-1}/I^n\subseteq I/I^n
    \]
    is finitely generated: \(I\) is finitely generated, so the same is true of
    \(I^{n-1}\). Coherence of \(R_{n-1}\) therefore makes \(J_n\) finitely
    presented over \(R_{n-1}\). Since its \(R_n\)-action factors through
    \(R_{n-1}=R_n/J_n\), the finite-presentation comparison of
    \Ref{lem:finite-presentation-quotient-comparison} makes \(J_n\) finitely presented over \(R_n\). The
    square-zero step now shows that \(R_n\) is coherent. Induction gives
    coherence of every \(R_n\), in particular \(A_0\).
    \item \textbf{Establish finite presentation of the powers.}
    Every \(I^k\) is a finitely generated ideal in the coherent ring \(A_0\) and
    is therefore finitely presented.
    \item \textbf{Establish finite presentation of the associated quotients.}
    If \(M\) is finitely presented over
    \(A_0\), the submodules \(I^kM\) and \(I^{k+1}M\) are finitely generated
    submodules of the finitely presented module \(M\), hence finitely presented by
    coherence. Their quotient is finitely presented over \(A_0\) and is
    annihilated by \(I\); \Ref{lem:finite-presentation-quotient-comparison} makes it finitely presented
    over \(B_0=A_0/I\).
\end{enumerate}
\end{proof}


Nilpotence and boundedness alone do not yet give a finite structured presentation. The
next step constructs one by separating the degree-zero nilpotent tower from the
positive-degree Postnikov layers and then refining the resulting coefficients over the
fixed centre.

\subsection{Finite structured square-zero generation}

\begin{lemma}[Bounded refinement of a square-zero coefficient filtration]
\label{lem:refine-square-zero-coefficient-filtration}
Let
\[
    C\longrightarrow D
\]
be a structured square-zero extension of connective commutative ring spectra
classified by a connective \(D\)-module \(M\). Suppose \(M\) admits a finite
filtration by \(D\)-submodules whose successive cofibres are
\(M_1,\ldots,M_r\), and for every \(j\) there exists \(n_j\geq0\) such that
\[
    M_j\in(\Mod_D)_{\geq n_j}\cap(\Mod_D)_{\leq2n_j}.
\]
Then \(C\to D\) admits a finite factorization by structured
square-zero extensions whose successive coefficients are
\(M_1,\ldots,M_r\). The factorization is functorial for the chosen filtered
coefficient.
\end{lemma}

\begin{proof}
\leavevmode
\begin{enumerate}[label=\textup{(\arabic*)}]
    \item \textbf{One filtration step.}
    It is enough to treat a cofibre sequence of connective \(D\)-modules
    \[
        M'\longrightarrow M\longrightarrow M''
    \]
    with \(M'\) in the stated small-extension range. Let
    \[
        \eta:L_{D/A}\longrightarrow M[1]
    \]
    classify \(C\to D\) in some compatible relative algebra category. Compose
    with \(M[1]\to M''[1]\) and let
    \[
        C''\longrightarrow D
    \]
    be the corresponding structured square-zero extension. Functoriality of
    split square-zero formation and of the defining pullbacks produces a
    canonical morphism
    \[
        C\longrightarrow C''
    \]
    over \(D\).
    \item \textbf{Compare the two fibre sequences.}
    Compare the underlying fibre sequences
    \[
        M\longrightarrow C\longrightarrow D,
        \qquad
        M''\longrightarrow C''\longrightarrow D.
    \]
    \item \textbf{Identify the relative fibre.}
    The \(3\times3\) lemma in spectra identifies
    \[
        \operatorname{fib}(C\to C'')\simeq M'.
    \]
    \item \textbf{Analyze the augmentation-ideal action.}
    We must also identify the module structure on this fibre. The augmentation
    ideal
    \[
        \operatorname{fib}(C''\to D)\simeq M''
    \]
    acts on \(M'\) through multiplication inside the original square-zero fibre
    \(M\).
    \item \textbf{Use null multiplication.}
    Since the nonunital multiplication
    \[
        M\otimes_C M\longrightarrow M
    \]
    is coherently null, this ideal action is null.
    \item \textbf{Identify the restricted module structure.}
    Hence the \(C''\)-module
    structure on \(M'\) factors canonically through \(C''\to D\), so the fibre is
    precisely the restriction of the specified \(D\)-module \(M'\). Its
    nonunital multiplication is likewise coherently null.
    \item \textbf{Verify the small-extension range.}
    If
    \[
        M'\in(\Mod_D)_{\geq n}\cap(\Mod_D)_{\leq2n},
    \]
    then \(C\to C''\) is an \(n\)-small extension in the sense of
    \cite[Theorem~7.4.1.26 and Corollary~7.4.1.27]{LurieHigherAlgebra}.
    \item \textbf{Apply small-extension recognition.}
    The
    small-extension recognition theorem therefore identifies it with a structured
    square-zero extension by \(M'\).
    \item \textbf{Factor the relative class.}
    In a relative algebra category, the
    structural map from the base supplies the required nullhomotopy of the
    absolute class on restriction, and cotangent transitivity factors the class
    through the relative cotangent complex.
    \item \textbf{Iterate the refinement.}
    Iterating this construction along the finite coefficient filtration gives the
    asserted factorization. In all applications below, the successive coefficients
    are Eilenberg--Mac Lane modules \(HN[j]\), which lie in the required range
    with \(n=j\).
\end{enumerate}
\end{proof}


\begin{lemma}[Connective pullbacks over a degree-zero surjection]
\label{lem:connective-nilpotent-pullback}
Let
\[
    C^u:=A_0\times_{A_{01}}A_1
\]
be a pullback formed in unrestricted commutative ring spectra, where
\(A_0,A_1,A_{01}\) are connective. If one of the maps
\[
    \pi_0(A_i)\longrightarrow\pi_0(A_{01})
\]
is surjective, then \(C^u\) is connective. Hence it also computes the pullback
in connective commutative ring spectra.
\end{lemma}

\begin{proof}
\leavevmode
\begin{enumerate}[label=\textup{(\arabic*)}]
    \item \textbf{Create the underlying fibre sequence.}
    The forgetful functor from commutative ring spectra to spectra creates limits,
    so there is a fibre sequence
    \[
        C^u
        \longrightarrow
        A_0\oplus A_1
        \longrightarrow
        A_{01}.
    \]
    \item \textbf{Identify the possible negative homotopy.}
    All three displayed algebras are connective. Consequently the only possible
    negative homotopy group of \(C^u\) is
    \[
        \pi_{-1}(C^u)
        \simeq
        \operatorname{coker}
        \bigl(
            \pi_0(A_0)\oplus\pi_0(A_1)
            \longrightarrow
            \pi_0(A_{01})
        \bigr).
    \]
    \item \textbf{Use degree-zero surjectivity.}
    Surjectivity of either component kills this cokernel, and all lower negative
    homotopy groups vanish termwise.
    \item \textbf{Conclude connectivity.}
    Thus \(C^u\) is connective.
    \item \textbf{Recover the connective pullback.}
    The inclusion of
    connective commutative ring spectra is full, so the same object satisfies the
    pullback universal property in the connective subcategory.
\end{enumerate}
\end{proof}

\begin{lemma}[Discrete square-zero quotients are structured]
\label{lem:discrete-square-zero-structured}
Let \(Q\twoheadrightarrow P\) be a surjection of ordinary commutative rings
with square-zero kernel \(J\). Then the induced morphism of connective
commutative ring spectra
\[
    HQ\longrightarrow HP
\]
is a structured square-zero extension of connective commutative ring spectra with coefficient \(HJ\). The
same assertion holds in any compatible relative commutative-algebra category.
\end{lemma}

\begin{proof}
\leavevmode
\begin{enumerate}[label=\textup{(\arabic*)}]
    \item \textbf{Identify the square-zero fibre.}
    The homotopy fibre of \(HQ\to HP\) is the discrete \(HP\)-module \(HJ\), and
    the multiplication on that fibre is null because \(J^2=0\).
    \item \textbf{Recognize a small extension.}
    Thus the arrow is
    a connective \(0\)-small extension.
    \item \textbf{Apply square-zero recognition.}
    The connective square-zero recognition
    theorem for commutative ring spectra identifies such an arrow with the
    structured extension classified by its canonical derivation
    \[
        L_{HP}\longrightarrow HJ[1];
    \]
    see \cite{LurieDAGIV,LurieHigherAlgebra}.
    \item \textbf{Factor the relative class.}
    If the original square-zero quotient
    is equipped with a compatible base-algebra structure, the structural map from
    the base supplies the nullhomotopy of the absolute class after restriction;
    cotangent transitivity therefore factors the class uniquely through the
    relative cotangent complex.
    \item \textbf{Conclude the relative statement.}
    This gives the relative statement.
\end{enumerate}
\end{proof}

\begin{proposition}[Bounded nilpotent finite structured square-zero generation]
\label{prop:bounded-nilpotent-square-zero-generation}
Let \(A\to B\) be a morphism of connective commutative ring spectra such that
\(\pi_0(A)\to\pi_0(B)\) is surjective with nilpotent kernel and
\(\operatorname{fib}(A\to B)\) is eventually coconnective. Then \(A\to B\)
admits a finite factorization by structured square-zero
extensions whose coefficients are restrictions of connective eventually
coconnective \(B\)-modules. The hypotheses are exactly surjectivity on \(\pi_0\) with nilpotent kernel
and eventual coconnectivity of the relative fibre.
\end{proposition}

\begin{proof}
\leavevmode
\begin{enumerate}[label=\textup{(\arabic*)}]
    \item \textbf{Recall the reconstruction tower.}
    The general reconstruction tower by structured square-zero extensions for a
    map surjective on \(\pi_0\) with nilpotent kernel is
    \cite[Lemma~2.1.14]{LurieDAGXIV}. We retain the explicit bounded argument
    here because the theorem additionally records finiteness under eventual
    coconnectivity and refines every coefficient to the restriction of a
    \(B\)-module.
    \item \textbf{Separate degree zero.}
    Put
    \[
        I:=\ker\bigl(\pi_0(A)\to\pi_0(B)\bigr),
        \qquad
        I^m=0.
    \]
    We first separate the degree-zero nilpotent part. Form, in unrestricted
    commutative ring spectra,
    \[
        A^{(1)}
        :=
        H\pi_0(A)\times_{H\pi_0(B)}B.
    \]
    Since \(H\pi_0(A)\to H\pi_0(B)\) is surjective on \(\pi_0\),
    \Ref{lem:connective-nilpotent-pullback} shows that \(A^{(1)}\) is connective;
    hence the same pullback computes the limit in connective commutative ring
    spectra. The canonical map \(A\to A^{(1)}\) is an isomorphism on \(\pi_0\).
    \item \textbf{Define the degree-zero pullbacks.}
    For \(1\leq k\leq m\), put
    \[
        D_k:=H\bigl(\pi_0(A)/I^k\bigr),
        \qquad
        \widetilde D_k^u:=D_k\times_{D_1}B
    \]
    with the pullback formed first in unrestricted commutative ring spectra.
    \item \textbf{Prove the pullbacks are connective.}
    The map \(D_k\to D_1\) is surjective on \(\pi_0\), so
    \Ref{lem:connective-nilpotent-pullback} makes
    \(\widetilde D_k^u\) connective. It therefore computes the same pullback
    in connective commutative ring spectra; write
    \[
        \widetilde D_k:=\widetilde D_k^u.
    \]
    \item \textbf{Identify the endpoints and transition pullbacks.}
    Then
    \[
        \widetilde D_1\simeq B,
        \qquad
        \widetilde D_m\simeq A^{(1)},
    \]
    and pullback pasting gives
    \[
        \widetilde D_{k+1}
        \simeq
        D_{k+1}\times_{D_k}\widetilde D_k.
    \]
    \item \textbf{Identify the discrete square-zero stages.}
    The map \(D_{k+1}\to D_k\) is the discrete structured square-zero
    quotient with coefficient \(H(I^k/I^{k+1})\) by
    \Ref{lem:discrete-square-zero-structured}.
    \item \textbf{Recognize their algebraic pullbacks.}
    The displayed identity
    identifies
    \[
        \widetilde D_{k+1}\longrightarrow\widetilde D_k
    \]
    with its algebraic pullback along
    \(\widetilde D_k\to D_k\).
    \item \textbf{Obtain the finite structured factorization.}
    Hence
    \Ref{lem:structured-square-zero-cartesian-pullback} gives
    \[
        A^{(1)}=\widetilde D_m
        \longrightarrow\cdots\longrightarrow
        \widetilde D_1=B
    \]
    is a finite structured square-zero factorization.
    \item \textbf{Identify the coefficients over the centre.}
    Since
    \(I(I^k/I^{k+1})=0\), the coefficient action factors through
    \(\pi_0(B)=\pi_0(A)/I\), and each coefficient is the restriction of the
    connective discrete \(B\)-module \(H(I^k/I^{k+1})\).
    \item \textbf{Begin the positive-degree tower.}
    It remains to factor \(A\to A^{(1)}\). The fibre of
    \(A^{(1)}\to B\) is \(HI\), while the fibre of \(A\to B\) is eventually
    coconnective. Fibre transitivity for
    \[
        A\longrightarrow A^{(1)}\longrightarrow B
    \]
    therefore shows that \(\operatorname{fib}(A\to A^{(1)})\) is eventually
    coconnective. Put
    \[
        C_1:=\operatorname{cofib}(A\to A^{(1)}).
    \]
    The map is an isomorphism on \(\pi_0\), so \(C_1\) is \(1\)-connective; choose
    \(T\) such that it is \(T\)-coconnective.
    \item \textbf{State the induction hypothesis.}
    Suppose inductively that maps
    \[
        A\longrightarrow A^{(i)}\longrightarrow A^{(1)}
    \]
    have been constructed with
    \[
        C_i:=\operatorname{cofib}(A\to A^{(i)})
    \]
    \(i\)-connective and \(T\)-coconnective.
    \item \textbf{Construct the relative Hurewicz map.}
    The canonical relative
    cofibre--cotangent morphism is the composite
    \[
        h_i:
        C_i
        \longrightarrow
        A^{(i)}\otimes_AC_i
        \longrightarrow
        L_{A^{(i)}/A}.
    \]
    \item \textbf{Identify its first homotopy map.}
    Since \(\pi_0(A)\simeq\pi_0(A^{(i)})\), the connective tensor-product spectral
    sequence makes the first arrow an isomorphism on \(\pi_i\), and
    \Ref{lem:relative-spectral-hurewicz-first-degree} makes the second arrow an
    isomorphism on \(\pi_i\). Hence
    \[
        \pi_i(h_i):\pi_iC_i\xrightarrow{\cong}\pi_iL_{A^{(i)}/A}.
    \]
    \item \textbf{Define the coefficient modules.}
    Put
    \[
        N_i:=\pi_iC_i,
        \qquad
        M_i:=HN_i[i-1].
    \]
    \item \textbf{Restrict the coefficient action.}
    Regard \(M_i\) as an \(A^{(i)}\)-module through the canonical truncation
    \[
        A^{(i)}\longrightarrow H\pi_0(A^{(i)})
        \simeq H\pi_0(A).
    \]
    \item \textbf{Define the classifying map.}
    The relative cotangent complex is \(i\)-connective. Its first Postnikov
    projection followed by the inverse of \(\pi_i(h_i)\) defines
    \[
        \eta_i:
        L_{A^{(i)}/A}
        \longrightarrow
        H\pi_iL_{A^{(i)}/A}[i]
        \xrightarrow{\ H(\pi_i(h_i)^{-1})[i]\ }
        HN_i[i]
        =M_i[1].
    \]
    \item \textbf{Compute the composite on first homotopy.}
    The composite
    \[
        C_i\xrightarrow{h_i}L_{A^{(i)}/A}
        \xrightarrow{\eta_i}HN_i[i]
    \]
    induces the identity on \(N_i=\pi_iC_i\).
    \item \textbf{Identify the Postnikov projection.}
    Because \(C_i\) is
    \(i\)-connective and \(HN_i[i]\) is \(i\)-truncated, maps between them are
    detected on \(\pi_i\); the composite is therefore exactly the canonical
    Postnikov projection
    \[
        C_i\longrightarrow H\pi_i(C_i)[i].
    \]
    \item \textbf{Construct the next structured square-zero extension.}
    Let \(A^{(i+1)}\to A^{(i)}\) be the structured square-zero extension in
    commutative \(A\)-algebras classified by \(\eta_i\).
    \item \textbf{Record its fibre sequence.}
    Its underlying fibre
    sequence is
    \[
        M_i\longrightarrow A^{(i+1)}\longrightarrow A^{(i)}.
    \]
    \item \textbf{Apply octahedrality.}
    Octahedrality for
    \(A\to A^{(i+1)}\to A^{(i)}\) gives a cofibre sequence
    \[
        C_{i+1}\longrightarrow C_i\longrightarrow M_i[1].
    \]
    \item \textbf{Identify the boundary map.}
    By the square-zero boundary comparison
    \Ref{lem:square-zero-boundary-hurewicz}, the final map is the composite
    \[
        C_i\xrightarrow{h_i}L_{A^{(i)}/A}
        \xrightarrow{\eta_i}M_i[1],
    \]
    which was just identified with the canonical Postnikov projection.
    \item \textbf{Identify the next Postnikov tail.}
    Therefore
    \[
        C_{i+1}
        \simeq
        \operatorname{fib}
        \bigl(C_i\to H\pi_i(C_i)[i]\bigr)
        \simeq
        \tau_{\geq i+1}C_i.
    \]
    \item \textbf{Update the connectivity bounds.}
    Thus the lower connectivity increases by one and the upper coconnective bound
    remains \(T\).
    \item \textbf{Refine coefficients over the centre.}
    The \(\pi_0(A)\)-module \(N_i\) has the finite \(I\)-adic filtration
    \[
        N_i\supseteq IN_i\supseteq\cdots\supseteq I^mN_i=0.
    \]
    Its successive quotients are annihilated by \(I\), so their action factors
    through \(\pi_0(B)\). Consequently \(M_i\) has a finite filtration whose
    successive cofibres are restrictions of the connective, eventually
    coconnective \(B\)-modules
    \[
        H(I^kN_i/I^{k+1}N_i)[i-1].
    \]
    Each displayed quotient is concentrated in homotopy degree \(i-1\), so it
    lies in the small-extension range required by
    \Ref{lem:refine-square-zero-coefficient-filtration}. Applying that lemma
    refines the single stage \(A^{(i+1)}\to A^{(i)}\) into finitely many
    structured square-zero extensions with coefficients restricted from those
    \(B\)-modules; the refinement preserves the composite and the Postnikov
    calculation above.
    \item \textbf{Termination.}
    For \(i>T\), one has \(C_i\simeq0\), so \(A\simeq A^{(i)}\). Concatenating
    the finitely many refined Postnikov-killing stages with the degree-zero
    nilpotent tower from \(A^{(1)}\) to \(B\) gives the asserted finite
    structured square-zero factorization.
\end{enumerate}
\end{proof}

\begin{proposition}[Finite structured square-zero generation of spectral Artinian extensions]
\label{prop:artinian-finite-square-zero-generation}
An augmented arrow \(A\to B\) is spectral Artinian if and only if it admits a finite factorization
\[
\begin{tikzcd}[column sep=large]
    A=A_m
        \arrow[r,"q_m"]
        &
    A_{m-1}
        \arrow[r]
        &
    \cdots
        \arrow[r]
        &
    A_1
        \arrow[r,"q_1"]
        &
    A_0=B .
\end{tikzcd}
\]
in which every
\[
    q_i:A_i\longrightarrow A_{i-1}
\]
is a structured square-zero extension whose coefficient is the restriction of
a connective coherent \(B\)-module.
Moreover, whenever \(A\to B\) is spectral Artinian, \(A\) is coherent and
\[
    A\to B
    \in
    \CAlg(\Sp)^{\mathrm{cn},\mathrm{afp},\wedge}_{R//B}.
\]
\end{proposition}

\begin{proof}
\leavevmode
\begin{enumerate}[label=\textup{(\arabic*)}]
    \item \textbf{Forward induction on the structured square-zero chain.}
    Suppose first that such a finite structured square-zero factorization is given. We prove
    inductively upward from \(A_0=B\) that every source \(A_i\) is coherent, that
    \(A_{i-1}\) is almost perfect over \(A_i\), and that \(B\) is almost perfect
    over \(A_i\). The assertions hold at \(i=0\). Assume them through
    \(A_{i-1}\). The coefficient of \(A_i\to A_{i-1}\) is the restriction of a
    connective coherent \(B\)-module \(M_i\). Since \(B\) is almost perfect over
    \(A_{i-1}\) and \(M_i\) is almost perfect over \(B\),
    \Ref{lem:almost-perfect-restriction-transitivity} makes \(M_i\) almost
    perfect over \(A_{i-1}\). Apply
    \Ref{lem:square-zero-almost-finite} with centre \(A_{i-1}\). It makes
    \(A_{i-1}\) almost perfect over \(A_i\), makes the stage almost finitely
    presented, and makes \(A_i\) coherent. \Ref{lem:almost-perfect-restriction-transitivity} then makes \(B\) almost
    perfect over \(A_i\).
    \item \textbf{Verify the Artinian conditions.}
    Thus the composite \(A\to B\) is almost finitely presented. Every stage has
    square-zero degree-zero kernel, so the kernel of the composite on \(\pi_0\) is
    nilpotent. Every connective coherent \(B\)-module is eventually coconnective, and the
    relative fibre of a finite composite is obtained from the stage fibres by
    finitely many extensions; hence \(\operatorname{fib}(A\to B)\) is eventually
    coconnective. Therefore \(A\to B\) is spectral Artinian. The nilpotent
    augmentation ideal makes \(A\) derived complete, while \(B\) is almost perfect
    over \(A\); \Ref{prop:complete-aft-intrinsic-characterization} puts the arrow
    in the complete almost finitely presented sector. The induction already proved coherence
    of \(A\).
    \item \textbf{Reverse implication: initial finiteness.}
    Conversely let \(A\to B\) be spectral Artinian and put
    \[
        I:=\ker\bigl(\pi_0(A)\to\pi_0(B)\bigr),
        \qquad
        I^m=0.
    \]
    The spectral criterion for almost finite presentation gives
    \[
        B\in\operatorname{APerf}(A).
    \]
    Since \(B\) is coherent,
    \Ref{lem:augmentation-kernel-finitely-presented} makes \(I\) finitely
    presented as a \(\pi_0(A)\)-module. Notice that the long exact sequence gives
    only a canonical quotient
    \[
        \pi_0\operatorname{fib}(A\to B)\twoheadrightarrow I;
    \]
    we never identify these two modules. By
    \Ref{lem:coherence-nilpotent-extension}, the ordinary ring \(\pi_0(A)\) is
    coherent, every \(I^k\) is finitely presented, and for every finitely presented
    \(\pi_0(A)\)-module \(N\) the pieces
    \[
        I^kN/I^{k+1}N
    \]
    are finitely presented over \(\pi_0(B)\). At this point the available coherence input concerns the ordinary ring
    \(\pi_0(A)\), and the following construction uses precisely that input.
    \item \textbf{Invoke bounded nilpotent factorization.}
    We now run the explicit bounded-nilpotent factorization of
    \Ref{prop:bounded-nilpotent-square-zero-generation}, keeping track of
    almost-perfectness so that its coefficients can be refined to coherent
    \(B\)-modules.
    \item \textbf{Degree-zero coherent stages.}
    First form
    \[
        A^{(1)}
        :=
        H\pi_0(A)\times_{H\pi_0(B)}B
    \]
    and the degree-zero tower
    \[
        A^{(1)}=\widetilde D_m
        \longrightarrow\cdots\longrightarrow
        \widetilde D_1=B,
        \qquad
        \widetilde D_k
        =H(\pi_0(A)/I^k)\times_{H\pi_0(B)}B.
    \]
    Each displayed pullback may be formed in unrestricted commutative ring
    spectra: the map
    \(H(\pi_0(A)/I^k)\to H\pi_0(B)\) is surjective on \(\pi_0\), so
    \Ref{lem:connective-nilpotent-pullback} shows that the result is
    connective and hence is also the pullback in the connective algebra
    category. Its stage coefficient is
    \[
        H(I^k/I^{k+1}).
    \]
    By \Ref{lem:coherence-nilpotent-extension}, \(I^k/I^{k+1}\) is finitely
    presented over \(\pi_0(B)\); since \(B\) is coherent, the corresponding
    Eilenberg--Mac Lane module is a connective coherent \(B\)-module. Thus the
    degree-zero tower already has the required coefficient class.
    \item \textbf{Almost-perfectness of the degree-zero tower.}
    We also need its sources to be almost perfect over \(A\). The target
    \(B=\widetilde D_1\) is almost perfect over \(A\). Every stage coefficient is
    almost perfect over \(B\), hence is almost perfect over \(A\) by
    \Ref{lem:almost-perfect-restriction-transitivity}. The underlying fibre sequence
    of each structured square-zero stage then shows inductively that every \(\widetilde D_k\),
    in particular \(A^{(1)}\), is almost perfect over \(A\). Therefore
    \[
        C_1:=\operatorname{cofib}(A\to A^{(1)})
    \]
    is almost perfect over \(A\), because \(A\) is perfect over itself.
    As in \Ref{prop:bounded-nilpotent-square-zero-generation}, \(C_1\) is
    \(1\)-connective and eventually coconnective; choose \(t\) such that it is
    \(t\)-coconnective.
    \item \textbf{State the positive-degree induction hypothesis.}
    Inductively suppose that
    \[
        A\longrightarrow A^{(i)}\longrightarrow A^{(1)}
    \]
    has been constructed such that \(A^{(i)}\) and
    \[
        C_i:=\operatorname{cofib}(A\to A^{(i)})
    \]
    are almost perfect over \(A\), while \(C_i\) is \(i\)-connective and
    \(t\)-coconnective.
    \item \textbf{Apply relative Hurewicz.}
    The relative Hurewicz construction of
    \Ref{prop:bounded-nilpotent-square-zero-generation} gives an isomorphism
    \[
        \pi_iC_i
        \xrightarrow{\cong}
        \pi_iL_{A^{(i)}/A}.
    \]
    \item \textbf{Define the coefficient modules.}
    Put
    \[
        N_i:=\pi_iC_i,
        \qquad
        M_i:=HN_i[i-1].
    \]
    \item \textbf{Restrict the coefficient action.}
    Regard \(M_i\) as an \(A^{(i)}\)-module through
    \(A^{(i)}\to H\pi_0(A^{(i)})\simeq H\pi_0(A)\).
    \item \textbf{Establish finite presentation of the coefficient.}
    Since \(C_i\) is
    \(i\)-connective and almost perfect,
    \Ref{lem:first-unresolved-homotopy-almost-perfect} makes \(N_i\) finitely
    presented over \(\pi_0(A)\).
    \item \textbf{Define the classifying term.}
    The first Postnikov projection of the relative
    cotangent complex and the inverse Hurewicz identification define the
    classifying term
    \[
        \eta_i:
        L_{A^{(i)}/A}
        \longrightarrow
        HN_i[i]
        =M_i[1].
    \]
    \item \textbf{Construct the next square-zero stage.}
    Let \(A^{(i+1)}\to A^{(i)}\) be the corresponding structured square-zero
    extension over \(A\).
    \item \textbf{Identify the next Postnikov tail.}
    Exactly as in
    \Ref{prop:bounded-nilpotent-square-zero-generation}, octahedrality identifies
    \[
        C_{i+1}:=\operatorname{cofib}(A\to A^{(i+1)})
        \simeq
        \tau_{\geq i+1}C_i.
    \]
    \item \textbf{Update the connectivity bound.}
    Thus connectivity increases by one and the fixed upper bound \(t\) remains.
    \item \textbf{Defer almost-perfectness to the refinement.}
    Almost-perfectness will be established from the refined structured square-zero stage
    using the coherence information already available for the graded pieces.
    \item \textbf{Refine the coefficients.}
    The finitely presented \(\pi_0(A)\)-module \(N_i\) has the finite
    \(I\)-adic filtration
    \[
        N_i\supseteq IN_i\supseteq\cdots\supseteq I^mN_i=0.
    \]
    By \Ref{lem:coherence-nilpotent-extension}, every quotient
    \[
        I^kN_i/I^{k+1}N_i
    \]
    is finitely presented over \(\pi_0(B)\). Therefore \(M_i\) has a finite
    filtration whose successive cofibres are the connective coherent \(B\)-modules
    \[
        H(I^kN_i/I^{k+1}N_i)[i-1].
    \]
    Each displayed quotient is concentrated in homotopy degree \(i-1\), hence
    lies in the small-extension range of
    \Ref{lem:refine-square-zero-coefficient-filtration}. Applying that lemma
    refines \(A^{(i+1)}\to A^{(i)}\) into finitely many structured square-zero
    stages with coefficients restricted from those coherent \(B\)-modules.
    \item \textbf{Propagate almost-perfectness.}
    Each refined coefficient is almost perfect over \(B\), and \(B\) is almost
    perfect over \(A\); hence \Ref{lem:almost-perfect-restriction-transitivity}
    makes each refined coefficient almost perfect over \(A\).
    Starting from the inductive hypothesis that \(A^{(i)}\) is almost perfect over
    \(A\), the square-zero fibre sequences through the finite refinement show that
    \(A^{(i+1)}\) is almost perfect over \(A\). Since \(A\) is perfect over
    itself, its cofibre \(C_{i+1}\) is also almost perfect over \(A\). This closes the induction using only the coherence input established for
    \(\pi_0(A)\).
    \item \textbf{Terminate the tower.}
    For \(i>t\), one has \(C_i\simeq0\), so \(A\simeq A^{(i)}\). Concatenating
    the finitely many refined Postnikov-killing stages with the degree-zero tower
    from \(A^{(1)}\) to \(B\) gives a finite factorization of \(A\to B\) by
    structured square-zero extensions whose coefficients are restrictions of
    connective coherent \(B\)-modules.
    \item \textbf{Recover spectral coherence.}
    It remains to prove spectral coherence of \(A\) from the constructed finite
    structured square-zero chain. Start at the coherent target \(A_0=B\) and move backward through
    the finite coherent structured square-zero chain just constructed. Suppose
    \(A_{j-1}\) is coherent and \(B\) is almost perfect over \(A_{j-1}\). The
    coefficient of \(A_j\to A_{j-1}\) is the restriction of a connective coherent
    \(B\)-module, hence is almost perfect over \(A_{j-1}\). Applying
    \Ref{lem:square-zero-almost-finite} with centre \(A_{j-1}\) proves that
    \(A_j\) is coherent and that \(A_{j-1}\) is almost perfect over \(A_j\).
    \Ref{lem:almost-perfect-restriction-transitivity} then makes \(B\) almost
    perfect over \(A_j\). Induction reaches
    \(A_m=A\) and proves that \(A\) is coherent.
    \item \textbf{Derived completeness.}
    Finally the nilpotence of \(I\) makes \(A\) derived \(I\)-complete, while
    \(B\) is already almost perfect over \(A\). Hence
    \Ref{prop:complete-aft-intrinsic-characterization} gives
    \[
        A\to B\in\CAlg(\Sp)^{\mathrm{cn},\mathrm{afp},\wedge}_{R//B}.
    \]
\end{enumerate}
\end{proof}

Finite square-zero generation gives bounded coefficients. For Schlessinger excision we
need the stronger relative statement that each stage can be refined by coefficients
restricted from connective coherent \(B\)-modules, uniformly at the chosen centre.

\subsection{Coherent refinement and relative factorization}

\begin{lemma}[Coherent refinement of an Artinian structured square-zero stage]
\label{lem:artinian-square-zero-coherent-refinement}
Let
\[
    C\longrightarrow D
\]
be a structured square-zero extension in
\(\CAlg(\Sp)^{\mathrm{cn}}_{R//B}\), with coefficient
\(M\in(\Mod_D)_{\geq0}\), and suppose that the augmented arrows
\(C\to B\) and \(D\to B\) are spectral Artinian. Then \(M\) admits a
finite filtration by \(D\)-submodules whose successive cofibres are
restrictions of connective coherent \(B\)-modules. Consequently \(C\to D\) admits a finite refinement by structured square-zero extensions with coefficients
restricted from connective coherent \(B\)-modules.
\end{lemma}

\begin{proof}
\leavevmode
\begin{enumerate}[label=\textup{(\arabic*)}]
    \item \textbf{Boundedness of the coefficient.}
    Because \(C\to B\) and \(D\to B\) are spectral Artinian, their relative
    fibres are eventually coconnective. Stable fibre calculus for the compatible
    maps to \(B\) identifies
    \[
        M\simeq\operatorname{fib}(C\to D)
    \]
    with the fibre of the induced map between those two relative fibres. Hence
    \(M\) is eventually coconnective; it is connective by the definition of the
    structured square-zero stage, so it is bounded.
    \item \textbf{Almost-perfectness.}
    By \Ref{prop:artinian-finite-square-zero-generation}, choose a finite coherent structured square-zero factorization
    \[
        D=D_r\longrightarrow D_{r-1}\longrightarrow\cdots\longrightarrow D_0=B.
    \]
    The Artinian augmentation \(C\to B\) satisfies
    \(B\in\operatorname{APerf}(C)\). Every stage coefficient in the displayed
    factorization is almost perfect over \(B\), hence is almost perfect over
    \(C\) by \Ref{lem:almost-perfect-restriction-transitivity}. Starting from
    \(D_0=B\), the square-zero fibre sequences therefore show inductively that
    \[
        D\in\operatorname{APerf}(C).
    \]
    The fibre sequence
    \[
        M\longrightarrow C\longrightarrow D
    \]
    then gives
    \[
        M\in\operatorname{APerf}(C).
    \]
    Both \(C\) and \(D\) are coherent by
    \Ref{prop:artinian-finite-square-zero-generation}. Hence
    \Ref{prop:coherent-almost-perfect-homotopy} makes every
    \(\pi_j(M)\) finitely presented over \(\pi_0(C)\).
    \item \textbf{Coherence over the intermediate centre.}
    Put
    \[
        J:=\ker\bigl(\pi_0(C)\to\pi_0(D)\bigr).
    \]
    The \(C\)-action on the square-zero coefficient factors through \(D\), so
    \(J\pi_j(M)=0\). If a module annihilated by \(J\) is finitely presented over
    \(\pi_0(C)\), tensoring a finite presentation with
    \(\pi_0(D)=\pi_0(C)/J\) makes it finitely presented over \(\pi_0(D)\).
    Thus every \(\pi_j(M)\) is finitely presented over \(\pi_0(D)\), and
    \Ref{prop:coherent-almost-perfect-homotopy} gives
    \[
        M\in\operatorname{Coh}(D)\cap(\Mod_D)_{\geq0}.
    \]
    \item \textbf{Define the centre ideal.}
    Now put
    \[
        I:=\ker\bigl(\pi_0(D)\to\pi_0(B)\bigr).
    \]
    \item \textbf{Establish nilpotence and finite presentation.}
    This ideal is nilpotent. Since \(D\to B\) is Artinian,
    \Ref{lem:augmentation-kernel-finitely-presented} makes \(I\) finitely
    presented over \(\pi_0(D)\).
    \item \textbf{Control the \(I\)-adic quotients.}
    Apply
    \Ref{lem:coherence-nilpotent-extension}. For each finitely presented module
    \(N_j:=\pi_j(M)\), the finite \(I\)-adic filtration
    \[
        N_j\supseteq IN_j\supseteq\cdots\supseteq I^sN_j=0
    \]
    has successive quotients
    \[
        I^kN_j/I^{k+1}N_j
    \]
    finitely presented over \(\pi_0(B)\).
    \item \textbf{Bound the Postnikov layers.}
    Because \(M\) is connective and eventually coconnective, its canonical
    Postnikov tower contains only finitely many nonzero layers
    \[
        HN_j[j].
    \]
    \item \textbf{Refine each Postnikov layer.}
    For each such layer, pull the finite \(I\)-adic filtration of \(N_j\) back
    successively along the canonical Postnikov projection
    \(\tau_{\leq j}M\to HN_j[j]\).
    \item \textbf{Assemble the coefficient filtration.}
    Splicing these finite refinements over the
    finitely many Postnikov degrees gives a finite filtration of \(M\) by
    \(D\)-submodules whose successive cofibres are restrictions of the modules
    \[
        H(I^kN_j/I^{k+1}N_j)[j].
    \]
    \item \textbf{Identify coherent centre coefficients.}
    Each of these is a connective coherent \(B\)-module by
    \Ref{prop:coherent-almost-perfect-homotopy}.
    \item \textbf{Verify the small-extension range.}
    Moreover it is concentrated in a
    single homotopy degree and therefore lies in the small-extension range of
    \Ref{lem:refine-square-zero-coefficient-filtration}.
    \item \textbf{Refine the structured extension.}
    Applying that lemma
    refines the original structured square-zero extension along this finite
    coefficient filtration.
\end{enumerate}
\end{proof}

\begin{proposition}[Relative factorization for surjections on \(\pi_0\) with nilpotent kernel]
\label{prop:artinian-relative-pi0-surjective-factorization}
Let
\[
    A'\longrightarrow A
\]
be a morphism in
\(\CAlg(\Sp)^{\mathrm{cn}}_{R//B}\), and suppose that
\(A\to B\) is spectral Artinian. The following conditions are equivalent:
\begin{enumerate}[label=(\roman*)]
    \item \(A'\to B\) is spectral Artinian and
    \[
        \pi_0(A')\longrightarrow\pi_0(A)
    \]
    is surjective with nilpotent kernel;

    \item there exists a finite factorization
    \[
        A'=A_m
        \longrightarrow
        A_{m-1}
        \longrightarrow
        \cdots
        \longrightarrow
        A_0=A
    \]
    in \(\CAlg(\Sp)^{\mathrm{cn}}_{R//B}\) such that every
    \(A_i\to A_{i-1}\) is a structured square-zero extension whose coefficient
    is the restriction of a connective coherent \(B\)-module.
\end{enumerate}
Under these conditions every intermediate augmentation
\[
    A_i\longrightarrow B
\]
is spectral Artinian.

Moreover, the full subcategory
\[
    \operatorname{Art}^{\mathrm{sp}}_{B/R}
\]
is stable under pullback along morphisms that are surjective on \(\pi_0\) with nilpotent kernel: if
\[
\begin{tikzcd}[column sep=large,row sep=large]
    C \arrow[r] \arrow[d]
        & A_0 \arrow[d]
    \\
    A_1 \arrow[r]
        & A_{01}
\end{tikzcd}
\]
is a pullback in connective commutative \(R\)-algebras over \(B\), all three
augmentations
\[
    A_0\to B,\qquad A_1\to B,\qquad A_{01}\to B
\]
are spectral Artinian, and \(A_0\to A_{01}\) is surjective on \(\pi_0\) with nilpotent kernel, then
\(C\to B\) is spectral Artinian.
\end{proposition}

\begin{proof}
\leavevmode
\begin{enumerate}[label=\textup{(\arabic*)}]
    \item \textbf{From Artinianity to relative module finiteness.}
    Because \(A'\to B\) is spectral Artinian,
    \Ref{thm:spectral-almost-finite-criterion} gives
    \[
        B\in\operatorname{APerf}(A').
    \]
    Choose a coherent structured square-zero factorization of the
    Artinian augmentation \(A\to B\) supplied by
    \Ref{prop:artinian-finite-square-zero-generation}:
    \[
        A=C_r\longrightarrow C_{r-1}\longrightarrow\cdots
        \longrightarrow C_0=B.
    \]
    Every stage coefficient is the restriction of a connective coherent
    \(B\)-module, hence is almost perfect over \(B\) and therefore almost perfect
    over \(A'\) by \Ref{lem:almost-perfect-restriction-transitivity}. Starting
    from \(C_0=B\in\operatorname{APerf}(A')\), the square-zero fibre sequences
    show inductively that
    \[
        C_j\in\operatorname{APerf}(A')
        \qquad(0\leq j\leq r).
    \]
    In particular,
    \[
        A\in\operatorname{APerf}(A').
    \]
    \item \textbf{Artinianity over the intermediate centre.}
    The map \(A'\to A\) is surjective on \(\pi_0\), so
    \Ref{thm:spectral-almost-finite-criterion} now makes it almost finitely
    presented. Stable fibre calculus for the composable maps
    \[
        A'\longrightarrow A\longrightarrow B
    \]
    gives a fibre sequence
    \[
        \operatorname{fib}(A'\to A)
        \longrightarrow
        \operatorname{fib}(A'\to B)
        \longrightarrow
        \operatorname{fib}(A\to B).
    \]
    The last two terms are eventually coconnective because both augmentations to
    \(B\) are Artinian. Hence \(\operatorname{fib}(A'\to A)\) is eventually
    coconnective. Together with the assumed nilpotence of the degree-zero kernel,
    this shows that
    \[
        A'\longrightarrow A
    \]
    is a spectral Artinian extension with centre \(A\). The centre \(A\) is
    coherent by \Ref{prop:artinian-finite-square-zero-generation} applied to
    \(A\to B\).
    \item \textbf{Factor over the intermediate centre.}
    Apply \Ref{prop:artinian-finite-square-zero-generation} again, now to the
    Artinian augmentation \(A'\to A\). We obtain a finite
    factorization
    \[
        A'=D_s\longrightarrow D_{s-1}\longrightarrow\cdots
        \longrightarrow D_0=A
    \]
    by structured square-zero extensions whose coefficients are restrictions of
    connective coherent \(A\)-modules. The forward implication in
    \Ref{prop:artinian-finite-square-zero-generation}, applied stage by stage,
    shows that every \(D_j\to A\) is spectral Artinian. Since \(A\to B\) is
    Artinian, each composite
    \[
        D_j\longrightarrow A\longrightarrow B
    \]
    is Artinian as well: almost finite presentation is stable under composition,
    nilpotent degree-zero kernels are stable under composition, and the relative
    fibre of a composite is an extension of the two relative fibres.
    \item \textbf{Refine coefficients over \(B\).}
    For every structured square-zero stage
    \[
        D_j\longrightarrow D_{j-1},
    \]
    both endpoint augmentations to \(B\) are therefore Artinian. Apply
    \Ref{lem:artinian-square-zero-coherent-refinement}. It refines the stage,
    into finitely many structured square-zero extensions
    whose coefficients are restrictions of connective coherent \(B\)-modules.
    Concatenating these finite refinements gives the factorization required in
    \textup{(ii)}.
    \item \textbf{Reverse induction.}
    Start with the Artinian augmentation \(A_0=A\to B\) and argue inductively.
    Suppose \(A_{i-1}\to B\) is Artinian and
    \[
        A_i\longrightarrow A_{i-1}
    \]
    is a structured square-zero extension with coefficient the restriction of a
    connective coherent \(B\)-module \(M_i\). The spectral criterion for almost finite presentation
    gives
    \[
        B\in\operatorname{APerf}(A_{i-1}),
    \]
    and restriction transitivity makes \(M_i\) almost perfect over \(A_{i-1}\).
    Then \Ref{lem:square-zero-almost-finite} makes the structured square-zero stage almost
    finitely presented. Hence the composite \(A_i\to B\) is almost finitely
    presented. Its degree-zero kernel is an extension of the nilpotent kernel of
    \(\pi_0(A_{i-1})\to\pi_0(B)\) by a square-zero kernel and is therefore
    nilpotent. Fibre transitivity expresses \(\operatorname{fib}(A_i\to B)\) as
    a finite extension of the eventually coconnective fibre of
    \(A_{i-1}\to B\) by the eventually coconnective coefficient \(M_i\).
    Therefore \(A_i\to B\) is spectral Artinian.
    \item \textbf{Complete the reverse implication.}
    Induction proves that every intermediate augmentation, in particular
    \(A'=A_m\to B\), is spectral Artinian. Since each map
    \(\pi_0(A_i)\to\pi_0(A_{i-1})\) is surjective with square-zero kernel, their
    finite composite
    \[
        \pi_0(A')\longrightarrow\pi_0(A)
    \]
    is surjective with nilpotent kernel. This proves \textup{(i)}.
    \item \textbf{Choose a factorization of the nilpotent-surjective leg.}
    Consider the Artinian morphism
    \[
        A_0\longrightarrow A_{01},
    \]
    which is surjective on \(\pi_0\) with nilpotent kernel. By the equivalence
    just proved, choose a finite factorization:
    \[
        A_0=D_m
        \longrightarrow
        D_{m-1}
        \longrightarrow
        \cdots
        \longrightarrow
        D_0=A_{01}.
    \]
    \item \textbf{Form the unrestricted pullbacks.}
    Form first, in unrestricted commutative ring spectra,
    \[
        E_k^u:=D_k\times_{A_{01}}A_1.
    \]
    \item \textbf{Prove the pullbacks are connective.}
    The composite
    \[
        D_k\longrightarrow A_{01}
    \]
    is surjective on \(\pi_0\), since it is a composite of square-zero
    surjections. Hence
    \Ref{lem:connective-nilpotent-pullback} makes every \(E_k^u\)
    connective. It therefore computes the same pullback in connective
    commutative \(R\)-algebras over \(B\); write
    \[
        E_k:=E_k^u.
    \]
    \item \textbf{Identify the endpoint pullbacks.}
    Then
    \[
        E_0\simeq A_1,
        \qquad
        E_m\simeq C.
    \]
    \item \textbf{Pull back the structured square-zero stages.}
    Algebraic pullback of structured square-zero extensions is structured by
    \Ref{lem:structured-square-zero-cartesian-pullback}; the coefficient of each
    \[
        E_k\longrightarrow E_{k-1}
    \]
    is the restriction of the same connective coherent \(B\)-module occurring in
    \(D_k\to D_{k-1}\).
    \item \textbf{Propagate Artinianity.}
    Starting from the Artinian target \(E_0=A_1\), the
    one-stage argument above shows inductively that every \(E_k\to B\) is
    Artinian. In particular \(C=E_m\to B\) is spectral Artinian.
\end{enumerate}
\end{proof}


The coherent factorization theorem has two semantic consequences. Artinian augmentations
lie inside the complete reconstruction domain, and affine realization identifies them
with the bounded coherent fixed-centre part of the finite infinitesimal geometry.

\subsection{Artinian tests as finite infinitesimal contexts}

\begin{corollary}[Artinian extensions are complete almost finitely presented]
\label{cor:artinian-complete-aft}
There is a fully faithful inclusion
\[
    \operatorname{Art}^{\mathrm{sp}}_{B/R}
    \hookrightarrow
    \CAlg(\Sp)^{\mathrm{cn},\mathrm{afp},\wedge}_{R//B}.
\]
\end{corollary}

\begin{proof}
\leavevmode
\begin{enumerate}[label=\textup{(\arabic*)}]
    \item \textbf{Invoke the Artinian generation theorem.}
    This is the final assertion of
    \Ref{prop:artinian-finite-square-zero-generation}.
\end{enumerate}
\end{proof}

\begin{proposition}[Artinian augmentations as coherent finite infinitesimal tests]
\label{prop:artinian-as-finite-infinitesimal-tests}
Let
\[
    U=\Spec(B)\longrightarrow V=\Spec(R).
\]
Affine realization defines a fully faithful functor
\[
    \bigl(\operatorname{Art}^{\mathrm{sp}}_{B/R}\bigr)^{\op}
    \longrightarrow
    \operatorname{Inf}^{\mathrm{fin}}(U/V)
\]
which sends an augmentation \(A\to B\) to the relative infinitesimal test
\[
\begin{tikzcd}[column sep=large,row sep=large]
    U \arrow[r,equal] \arrow[d,hook]
        & U \arrow[d]
    \\
    \Spec(A) \arrow[r]
        & V.
\end{tikzcd}
\]
Its essential image is the full subcategory of relative finite infinitesimal
tests whose open part is all of \(U\) and whose affine coordinate augmentation
is spectral Artinian. In particular, spectral Artinian tests form the
bounded coherent fixed-centre fragment of the finite infinitesimal geometry
constructed in Chapter~\Ref{chap:differentiation-spectral-doctrine}.
\end{proposition}

\begin{proof}
\leavevmode
\begin{enumerate}[label=\textup{(\arabic*)}]
    \item \textbf{Affine realization.}
    By \Ref{prop:artinian-finite-square-zero-generation}, an Artinian augmentation
    admits a finite factorization
    \[
        A=A_m\longrightarrow A_{m-1}\longrightarrow\cdots\longrightarrow A_0=B
    \]
    by structured square-zero extensions. Passing to affine geometry
    gives a finite composite
    \[
        \Spec(B)\longrightarrow\Spec(A_1)\longrightarrow\cdots
        \longrightarrow\Spec(A)
    \]
    of square-zero infinitesimal substitutions. Hence
    \(\Spec(B)\to\Spec(A)\) is a finite spectral infinitesimal substitution by
    \Ref{def:finite-spectral-infinitesimal-substitution}. The maps from \(R\) and
    to \(B\) give the displayed incidence diagram over \(U\to V\).
    \item \textbf{Full faithfulness.}
        Morphisms in \(\operatorname{Art}^{\mathrm{sp}}_{B/R}\) are precisely
        commutative morphisms of augmented \(R\)-algebras over \(B\). Under affine
    realization these are precisely the morphisms between the displayed relative
        infinitesimal tests whose open part is the identity of \(U\). This proves full
        faithfulness. The description of the essential image is immediate from the
        construction and the intrinsic definition of spectral Artinian extension.
        The functor depends only on the augmented arrow.
\end{enumerate}
\end{proof}




One final presentation is needed only for the Schlessinger argument. When the
coefficient of a structured square-zero stage comes from the centre, the stage can be
rewritten as a pullback against the split square-zero algebra at that centre.

\subsection{Centre presentations for Schlessinger excision}

\begin{lemma}[Centre presentation of a relative structured square-zero stage]
\label{lem:square-zero-centre-presentation}
Let
\[
    C\longrightarrow D
\]
be a structured square-zero extension in
\(\CAlg(\Sp)^{\mathrm{cn}}_{R//B}\), and suppose its coefficient \(M\) is the
restriction of a connective \(B\)-module. Then there is a canonical morphism
over \(B\)
\[
    \sigma:D\longrightarrow B\oplus M[1]
\]
such that
\[
    C\simeq D\times_{B\oplus M[1]}B,
\]
where \(B\to B\oplus M[1]\) is the zero section. The comparison is natural
in the structured relative square-zero stage.
\end{lemma}

\begin{proof}
\leavevmode
\begin{enumerate}[label=\textup{(\arabic*)}]
    \item \textbf{Choose the classifying derivation section.}
    Let
    \[
        d_\delta:D\longrightarrow D\oplus M[1]
    \]
    be the derivation section classifying \(C\to D\).
    \item \textbf{Use split square-zero base change.}
    Split square-zero formation
    for a coefficient restricted from \(B\) gives a Cartesian square
    \[
    \begin{tikzcd}[column sep=large,row sep=large]
        D\oplus M[1] \arrow[r] \arrow[d]
            & B\oplus M[1] \arrow[d]
        \\
        D \arrow[r]
            & B.
    \end{tikzcd}
    \]
    \item \textbf{Define the centre map.}
    Compose \(d_\delta\) with the upper horizontal morphism to obtain
    \(\sigma:D\to B\oplus M[1]\).
    \item \textbf{Apply pullback pasting.}
    The defining square-zero pullback and the
    preceding Cartesian square then give, by pullback pasting,
    \[
    \begin{aligned}
        C
        &=D\times_{D\oplus M[1]}D\\
        &\simeq D\times_{B\oplus M[1]}B.
    \end{aligned}
    \]
    \item \textbf{Record functoriality.}
    All comparisons are functorial consequences of split square-zero formation and
    pullback pasting.
\end{enumerate}
\end{proof}


We may now retain a compact description of the Artinian category: it is simultaneously a
full subcategory of complete fixed-centre reconstruction and a bounded coherent
subcategory of finite infinitesimal geometry. Its coherent square-zero factorizations
reduce nilpotent pullbacks to elementary centre-based stages, which is the finite test
calculus needed next.

\section{Artinian formal-moduli problems and Schlessinger localization}
\label{sec:schlessinger-formal-moduli}

The Artinian category is now a finite test theory at the fixed centre \(B\). We seek
space-valued functors that behave as nonlinear formal neighbourhoods. Rather than
choosing additional coherence data, we impose two equations by localization: the centre
is sent to a point, and every elementary square-zero pullback is sent to a homotopy
pullback.

For a generator
\[
    A_\sigma=A\times_{B\oplus M[1]}B,
\]
the algebraic pullback and the corresponding semantic condition are displayed together:
\[
\begin{tikzcd}[column sep=large,row sep=large]
    A_\sigma \arrow[r] \arrow[d]
        & A \arrow[d,"\sigma"]
    &&
    F(A_\sigma) \arrow[r] \arrow[d]
        & F(A) \arrow[d]
    \\
    B \arrow[r,"d_0"']
        & B\oplus M[1]
    &&
    F(B) \arrow[r]
        & F(B\oplus M[1]).
\end{tikzcd}
\]
The left square constructs the elementary Artinian extension, and locality requires the
right square to be Cartesian. Finite coherent factorization then promotes these
elementary equations to the usual Schlessinger excision condition for arbitrary Artinian
nilpotent pullbacks.

\begin{definition}[Artinian functor category]
\label{def:artinian-semantics-category}
Put
\[
    \mathcal P^{\mathrm{Art},\mathrm{sp}}_{B/R}
    :=
    \Fun\bigl(
        \operatorname{Art}^{\mathrm{sp}}_{B/R},
        \mathcal S
    \bigr).
\]
For \(A\to B\) in the Artinian category, write
\[
    y^A(C):=
    \Map_{\operatorname{Art}^{\mathrm{sp}}_{B/R}}(A,C)
\]
for the covariant Yoneda object. Thus
\[
    \Map(y^A,F)\simeq F(A).
\]
\end{definition}

\begin{lemma}[Schlessinger pullbacks remain Artinian]
\label{lem:schlessinger-pullback-artinian}
Let \(A\to B\) be spectral Artinian, let \(M\) be a connective coherent
\(B\)-module, and let
\[
    \sigma:A\longrightarrow B\oplus M[1]
\]
be a morphism over \(B\). Put
\[
    A_\sigma:=A\times_{B\oplus M[1]}B,
\]
where the second map is the zero section. Then \(A_\sigma\to A\) is a
structured square-zero extension with coefficient the restriction of \(M\),
and \(A_\sigma\to B\) is spectral Artinian.
\end{lemma}

\begin{proof}
\leavevmode
\begin{enumerate}[label=\textup{(\arabic*)}]
    \item \textbf{Form the unrestricted pullback.}
    Form \(A_\sigma\) first in unrestricted commutative ring spectra. The zero
    section
    \[
        B\longrightarrow B\oplus M[1]
    \]
    is an isomorphism on \(\pi_0\), so
    \Ref{lem:connective-nilpotent-pullback} makes \(A_\sigma\) connective.
    Hence this unrestricted pullback is also the pullback in connective
    commutative \(R\)-algebras over \(B\).
    \item \textbf{Identify the base-changed split square-zero square.}
    Regard \(M\) as an \(A\)-module through \(A\to B\). Split square-zero
    formation gives a Cartesian square
    \[
    \begin{tikzcd}[column sep=large, row sep=large]
        A\oplus M[1] \arrow[r] \arrow[d]
            & B\oplus M[1] \arrow[d]
        \\
        A \arrow[r]
            & B.
    \end{tikzcd}
    \]
    \item \textbf{Lift the centre map.}
    Because \(\sigma\) lies over \(A\to B\), it has a contractibly unique lift
    \[
        \widetilde\sigma:A\longrightarrow A\oplus M[1]
    \]
    over \(A\).
    \item \textbf{Identify the derivation section.}
    This is the derivation section classified by the corresponding
    degree-one derivation of \(A\) with coefficient \(M\).
    \item \textbf{Apply pullback pasting.}
    Pullback pasting gives
    \[
        A_\sigma
        \simeq
        A\times_{A\oplus M[1]}A.
    \]
    \item \textbf{Identify the structured extension.}
    Thus \(A_\sigma\to A\) is the structured square-zero extension classified by
    \(\widetilde\sigma\). Its coefficient is the restriction of the connective
    coherent \(B\)-module \(M\).
    \item \textbf{Artinianity.}
    Choose a finite coherent structured square-zero factorization of
    \(A\to B\) using
    \Ref{prop:artinian-finite-square-zero-generation}. Prepending the constructed
    stage \(A_\sigma\to A\) gives such a factorization of \(A_\sigma\to B\). The
    same proposition therefore makes \(A_\sigma\to B\) spectral Artinian.
\end{enumerate}
\end{proof}

\begin{definition}[Artinian Schlessinger generators]
\label{def:artinian-schlessinger-generators}
Let
\[
    W^{\mathrm{Sch},\mathrm{Art}}_{B/R}
\]
be the following set of morphisms in
\(\mathcal P^{\mathrm{Art},\mathrm{sp}}_{B/R}\):
\begin{enumerate}[label=(\roman*)]
    \item the centre generator
    \[
        \varnothing\longrightarrow y^B;
    \]
    \item for every spectral Artinian \(A\to B\), every connective coherent
    \(B\)-module \(M\), and every morphism over \(B\)
    \[
        \sigma:A\longrightarrow B\oplus M[1],
    \]
    the Yoneda morphism
    \[
        y^A
        \amalg_{y^{B\oplus M[1]}}
        y^B
        \longrightarrow
        y^{A_\sigma},
        \qquad
        A_\sigma=A\times_{B\oplus M[1]}B.
    \]

    The corresponding algebra square is Cartesian:
    \[
    \begin{tikzcd}[column sep=large,row sep=large]
        A_\sigma \arrow[r] \arrow[d]
            & A \arrow[d,"\sigma"]
        \\
        B \arrow[r,"d_0"']
            & B\oplus M[1].
    \end{tikzcd}
    \]
\end{enumerate}
The split square-zero augmentation \(B\oplus M[1]\to B\) is Artinian: its map on
\(\pi_0\) is an isomorphism, its fibre \(M[1]\) is eventually
coconnective because \(M\) is coherent, and its augmentation is almost
finitely presented by \Ref{lem:square-zero-almost-finite}. The object
\(A_\sigma\to B\) is Artinian by
\Ref{lem:schlessinger-pullback-artinian}.
\end{definition}

\begin{definition}[Spectral formal-moduli problem]
\label{def:spectral-formal-moduli-problem}
Define
\[
    \operatorname{FMP}^{\mathrm{sp}}_{B/R}
    \subseteq
    \mathcal P^{\mathrm{Art},\mathrm{sp}}_{B/R}
\]
to be the full subcategory of
\(W^{\mathrm{Sch},\mathrm{Art}}_{B/R}\)-local objects. Its objects are the
\textbf{spectral formal-moduli problems} at the centre \(B\) over \(R\), with \(B\) coherent.
\end{definition}

\begin{theorem}[Artinian Schlessinger reflector]
\label{thm:artinian-schlessinger-reflector}
The inclusion
\[
    \operatorname{FMP}^{\mathrm{sp}}_{B/R}
    \hookrightarrow
    \mathcal P^{\mathrm{Art},\mathrm{sp}}_{B/R}
\]
admits an accessible left adjoint
\[
\begin{tikzcd}[column sep=huge]
    \mathcal P^{\mathrm{Art},\mathrm{sp}}_{B/R}
        \arrow[r,shift left=1.2ex,"L^{\mathrm{Sch},\mathrm{Art}}_{B/R}"]
        &
    \operatorname{FMP}^{\mathrm{sp}}_{B/R}
        \arrow[l,shift left=1.2ex,hook]
\end{tikzcd}
\]
A functor \(F\) is local precisely when
\[
    F(B)\simeq *
\]
and, for every generator square,
\[
    F(A_\sigma)
    \xrightarrow{\simeq}
    F(A)\times_{F(B\oplus M[1])}F(B).
\]
\end{theorem}

\begin{proof}
\leavevmode
\begin{enumerate}[label=\textup{(\arabic*)}]
    \item \textbf{Choose a set of generators.}
    The Artinian category is essentially small in the fixed universe, so the
    Yoneda generators form a set.
    \item \textbf{Use presentability of the functor category.}
    Its space-valued functor category is
    presentable.
    \item \textbf{Construct the accessible localization.}
    Accessible localization at the strongly saturated class generated by that set
    therefore exists by
    \cite[Proposition~5.5.4.15]{LurieHigherToposTheory}.
    \item \textbf{Apply Yoneda.}
    Yoneda gives
    \[
        \Map(y^A,F)\simeq F(A).
    \]
    \item \textbf{Identify the centre condition.}
    Locality for \(\varnothing\to y^B\) is exactly \(F(B)\simeq *\).
    \item \textbf{Identify the square-zero condition.}
    Mapping the
    second generator into \(F\) turns its pushout into the displayed pullback of
    values.
\end{enumerate}
\end{proof}


\begin{proposition}[Schlessinger excision follows from the generators]
\label{prop:formal-moduli-finite-excision}
Let
\[
    F\in\operatorname{FMP}^{\mathrm{sp}}_{B/R}.
\]
Consider a pullback square in
\(\operatorname{Art}^{\mathrm{sp}}_{B/R}\)
\[
\begin{tikzcd}[column sep=large,row sep=large]
    A \arrow[r] \arrow[d]
        & A_0 \arrow[d]
    \\
    A_1 \arrow[r]
        & A_{01},
\end{tikzcd}
\qquad
A\simeq A_0\times_{A_{01}}A_1,
\]
and suppose that
\[
    A_0\longrightarrow A_{01}
\]
is surjective on \(\pi_0\) with nilpotent kernel. Then the canonical map
\[
    F(A)
    \longrightarrow
    F(A_0)
    \times_{F(A_{01})}
    F(A_1)
\]
is an equivalence.

Consequently, the generator formulation of
\Ref{def:spectral-formal-moduli-problem} is equivalent to the ordinary
Schlessinger requirement that Artinian pullback squares with one leg surjective on \(\pi_0\) with nilpotent kernel
be carried to pullback squares of spaces. The conclusion is independent of
every finite structured square-zero factorization, coherent refinement, and centre
presentation used inside the proof.
\end{proposition}

\begin{proof}
\leavevmode
\begin{enumerate}[label=\textup{(\arabic*)}]
    \item \textbf{Factor the leg surjective on \(\pi_0\) with nilpotent kernel.}
    By
    \Ref{prop:artinian-relative-pi0-surjective-factorization}, the Artinian morphism
    \[
        A_0\longrightarrow A_{01},
    \]
    which is surjective on \(\pi_0\) with nilpotent kernel, admits a finite
    factorization
    \[
        A_0=D_m
        \longrightarrow
        D_{m-1}
        \longrightarrow
        \cdots
        \longrightarrow
        D_0=A_{01},
    \]
    where every
    \[
        D_k\longrightarrow D_{k-1}
    \]
    is a structured square-zero extension whose coefficient is the restriction of
    a connective coherent \(B\)-module \(M_k\). Every \(D_k\to B\) is
    Artinian.
    \item \textbf{Pull back the factorization.}
    Pull this factorization back along
    \[
        A_1\longrightarrow A_{01}
    \]
    by forming, first in unrestricted commutative ring spectra,
    \[
        E_k^u:=D_k\times_{A_{01}}A_1.
    \]
    Every composite \(D_k\to A_{01}\) is surjective on \(\pi_0\). Therefore
    \Ref{lem:connective-nilpotent-pullback} makes each \(E_k^u\) connective,
    so it also computes the pullback in the connective algebra category. Put
    \[
        E_k:=E_k^u.
    \]
    Then
    \[
        E_0\simeq A_1,
        \qquad
        E_m\simeq A.
    \]
    By algebraic pullback stability of structured square-zero extensions,
    \Ref{lem:structured-square-zero-cartesian-pullback}, every
    \[
        E_k\longrightarrow E_{k-1}
    \]
    is a structured square-zero extension with coefficient the restriction of
    \(M_k\). The pullback clause of
    \Ref{prop:artinian-relative-pi0-surjective-factorization} makes every
    \(E_k\to B\) spectral Artinian.
    \item \textbf{Present each stage at the centre.}
    For each stage
    \[
        D_k\longrightarrow D_{k-1},
    \]
    the centre-presentation lemma
    \Ref{lem:square-zero-centre-presentation} constructs a canonical map
    \[
        \sigma_k:
        D_{k-1}
        \longrightarrow
        B\oplus M_k[1]
    \]
    and a Cartesian identification
    \[
        D_k
        \simeq
        D_{k-1}
        \times_{B\oplus M_k[1]}
        B.
    \]
    After pullback to \(E_{k-1}\), the same construction gives
    \[
        E_k
        \simeq
        E_{k-1}
        \times_{B\oplus M_k[1]}
        B.
    \]
    \item \textbf{Apply generator locality.}
    Since \(F\) is local for the Schlessinger generators,
    \[
        F(D_k)
        \xrightarrow{\simeq}
        F(D_{k-1})
        \times_{F(B\oplus M_k[1])}
        F(B)
    \]
    and
    \[
        F(E_k)
        \xrightarrow{\simeq}
        F(E_{k-1})
        \times_{F(B\oplus M_k[1])}
        F(B).
    \]
    \item \textbf{Combine the two pullback descriptions.}
    Combining these two pullback descriptions gives a canonical equivalence
    \[
        F(E_k)
        \xrightarrow{\simeq}
        F(D_k)
        \times_{F(D_{k-1})}
        F(E_{k-1}).
    \]
    \item \textbf{Identify the elementary pullback square.}
    Thus \(F\) carries every elementary pullback square
    \[
    \begin{tikzcd}[column sep=large,row sep=large]
        E_k \arrow[r] \arrow[d]
            & D_k \arrow[d]
        \\
        E_{k-1} \arrow[r]
            & D_{k-1}
    \end{tikzcd}
    \]
    to a pullback square of spaces.
    \item \textbf{Paste the elementary squares.}
    Finite pasting now gives
    \[
        F(E_m)
        \simeq
        F(D_m)
        \times_{F(D_0)}
        F(E_0).
    \]
    Using
    \[
        E_m\simeq A,
        \qquad
        D_m=A_0,
        \qquad
        D_0=A_{01},
        \qquad
        E_0=A_1,
    \]
    we obtain
    \[
        F(A)
        \simeq
        F(A_0)
        \times_{F(A_{01})}
        F(A_1).
    \]
    \item \textbf{Converse implication.}
    Conversely, every Schlessinger generator square of
    \Ref{def:artinian-schlessinger-generators} is itself an Artinian pullback
    square with one leg surjective on \(\pi_0\) with nilpotent kernel: the zero section
    \[
        B\longrightarrow B\oplus M[1]
    \]
    is an isomorphism on \(\pi_0\). Hence preservation of all such Artinian
    pullbacks implies locality for the generator squares. Together with
    the centre condition \(F(B)\simeq *\), this proves the asserted equivalence of
    the two formulations.
\end{enumerate}
\end{proof}



Thus spectral formal-moduli problems are exactly the Artinian functors satisfying the
centre condition and elementary square-zero excision; finite coherent factorization
upgrades those generators to full nilpotent excision. The centre can now be removed by
retaining only the two reconstruction equations that make sense globally.

\section{Centre-free reconstruction by continuity and nilpotent excision}
\label{sec:reconstruction-localization}

The Artinian theory is centred and finitary. Over an arbitrary spectral Zariski base
\(S\), two of its reconstruction principles survive without choosing a centre, and we
impose exactly those two principles in the slice over \(S\).

The first is Postnikov continuity: values on a connective affine test are recovered from
its finite Postnikov truncations. The second is nilpotent excision: algebraic pullbacks
with one nilpotent-surjective leg are sent to homotopy pullbacks of mapping spaces.
Accessible localization imposes these equations; left exactness is then the extra
statement that makes the localized semantics compatible with finite-limit geometry.
Their roles are summarized by
\[
\begin{tikzcd}[column sep=tiny,row sep=large]
    \text{Postnikov-continuity generators}
        \arrow[dr]
    &&
    \text{nilpotent-excision generators}
        \arrow[dl]
    \\
    & W^{\mathrm{Recon}}_S
        \arrow[d,"\text{accessible localization}"]
    &
    \\
    & \bigl((\mathcal H_{\Sp})_{/S}\bigr)^{\mathrm{Recon}}
        \arrow[d,"\text{left exactness}"]
    &
    \\
    & \text{reconstruction-local }\infty\text{-topos over }S
    &
\end{tikzcd}
\]

The reflector therefore gives only the candidate semantics. The substantive step is to
show that its generating local equivalences are stable under pullback. That result
implies base-change stability, finite-limit preservation, and the existence of centred
leaves inside the localized semantics.

\subsection{Reconstruction generators}

Fix a base
\[
    S\in\mathcal H_{\Sp}
\]
and work in the slice
\[
    (\mathcal H_{\Sp})_{/S}.
\]
An \textbf{affine spectral test over \(S\)} is a morphism
\[
    h_A\longrightarrow S
\]
with \(A\) connective; these are the tests on which the reconstruction equations are
imposed.

\begin{definition}[Representability by spectral schemes]
\label{def:representable-by-spectral-schemes}
A morphism
\[
    Y\longrightarrow S
\]
in \(\mathcal H_{\Sp}\) is
\textbf{representable by spectral schemes} if, for every affine spectral test
\(T\to S\), the fibre product
\[
    Y\times_ST
\]
is a spectral scheme. This condition is a property of the morphism. In particular, every morphism of
spectral schemes is representable by spectral schemes, by finite-limit closure
of spectral schemes from Chapter~\Ref{chap:spectral-algebraic-geometry}.
\end{definition}

\begin{definition}[Global spectral reconstruction generators]
\label{def:global-spectral-reconstruction-generators}
Let \(W^{\mathrm{Recon}}_S\) be the following set of morphisms in
\((\mathcal H_{\Sp})_{/S}\).
\begin{enumerate}[label=(\roman*)]
    \item For every affine spectral test \(h_A\to S\), the
    \textbf{Postnikov-continuity generator}
    \[
        c_A:
        \operatorname*{colim}_{n\geq0}h_{\tau_{\leq n}^{\CAlg}A}
        \longrightarrow
        h_A,
    \]
    where every \(h_{\tau_{\leq n}^{\CAlg}A}\) maps to \(S\) through
    \(h_A\to S\).

    \item Let
    \[
    \begin{tikzcd}[column sep=large, row sep=large]
        C \arrow[r] \arrow[d]
            & A_0 \arrow[d]
        \\
        A_1 \arrow[r]
            & A_{01}
    \end{tikzcd}
    \qquad
    C\simeq A_0\times_{A_{01}}A_1
    \]
    be a pullback of connective commutative ring spectra in which one
    \(A_i\to A_{01}\) is surjective on \(\pi_0\) with nilpotent kernel.
    For every affine spectral test \(h_C\to S\), regard
    \(h_{A_0},h_{A_1},h_{A_{01}}\) as objects over \(S\) through
    \(h_C\to S\). The associated \textbf{nilpotent-excision generator} is
    \[
        w_\square:
        h_{A_0}\amalg_{h_{A_{01}}}h_{A_1}
        \longrightarrow
        h_C.
    \]
\end{enumerate}
The pullback \(C\) in \textup{(ii)} is connective by
\Ref{lem:connective-nilpotent-pullback}.
\end{definition}

\begin{lemma}[Representable nilpotent excision]
\label{lem:represented-nilpotent-excision}
Let
\[
    C\simeq A_0\times_{A_{01}}A_1
\]
be a generator square of
\Ref{def:global-spectral-reconstruction-generators}, equipped with its affine
test \(\Spec(C)\to S\). If \(Y\to S\) is representable by spectral
schemes, then the canonical map
\[
    \Map_S(\Spec C,Y)
    \longrightarrow
    \Map_S(\Spec A_0,Y)
    \times_{\Map_S(\Spec A_{01},Y)}
    \Map_S(\Spec A_1,Y)
\]
is an equivalence.
\end{lemma}

\begin{proof}
\leavevmode
\begin{enumerate}[label=\textup{(\arabic*)}]
    \item \textbf{Base change to the affine source.}
    Put
    \[
        Y_C:=Y\times_S\Spec(C).
    \]
    By representability, \(Y_C\) is a spectral scheme. Since every object in the
    generator square maps to \(S\) through \(\Spec(C)\), the mapping spaces in the
    statement identify with the corresponding mapping spaces over \(\Spec(C)\)
    into \(Y_C\).
    \item \textbf{Nilpotent geometry.}
    After interchanging the two indices if necessary, assume that
    \[
        \pi_0(A_0)\longrightarrow\pi_0(A_{01})
    \]
    is surjective with nilpotent kernel \(J\). Then
    \[
        \Spec(A_{01})\longrightarrow\Spec(A_0)
    \]
    is a nilpotent closed immersion and a homeomorphism on underlying
    topological spaces.

    The pullback algebra \(C=A_0\times_{A_{01}}A_1\) is connective by
    \Ref{lem:connective-nilpotent-pullback}, and its degree-zero homotopy ring is
    \[
        \pi_0(C)
        \simeq
        \pi_0(A_0)\times_{\pi_0(A_{01})}\pi_0(A_1).
    \]
    The projection \(\pi_0(C)\to\pi_0(A_1)\) is surjective, and its kernel
    identifies with \(J\) through the first projection. In particular that
    kernel is nilpotent. Consequently
    \[
        \Spec(A_1)\longrightarrow\Spec(C)
    \]
    is also a nilpotent closed immersion and a homeomorphism on underlying
    topological spaces. This conclusion is obtained from the algebraic
    pullback; no geometric Cartesian identification of the displayed affine
    square is being asserted.
    \item \textbf{Affine target.}
    If \(Y_C=\Spec(D)\) is affine, the result is the algebraic pullback universal
    property:
    \[
    \begin{aligned}
        \Map_C(\Spec C,\Spec D)
        &\simeq
        \Map_{\CAlg(\Sp)_{C/}}(D,C)
        \\
        &\simeq
        \Map_{\CAlg(\Sp)_{C/}}(D,A_0)
        \times_{\Map_{\CAlg(\Sp)_{C/}}(D,A_{01})}
        \Map_{\CAlg(\Sp)_{C/}}(D,A_1).
    \end{aligned}
    \]
    \item \textbf{Choose affine local data.}
    For general \(Y_C\), let a point of the right-hand side be represented by
    compatible maps
    \[
        f_0:\Spec(A_0)\longrightarrow Y_C,
        \qquad
        f_1:\Spec(A_1)\longrightarrow Y_C.
    \]
    Choose a finite principal affine cover of \(\Spec(C)\), equivalently of
    \(\Spec(A_1)\), such that each restriction of \(f_1\) lands in an affine open
    of \(Y_C\). Because
    \(\Spec(A_{01})\to\Spec(A_0)\) is a nilpotent homeomorphism on support and
    \(f_0\) agrees with \(f_1\) after restriction to \(\Spec(A_{01})\), the
    corresponding restrictions of \(f_0\) land in the same affine opens.
    \item \textbf{Čech refinement.}
    Refine each intersection of affine target opens by affine opens, and refine
    the corresponding principal source intersection by a finite principal cover
    subordinate to them, following the affine-open mapping-descent argument of
    Chapter~\Ref{chap:spectral-algebraic-geometry}.
    Repeat on higher intersections. On every affine member of the resulting
    Čech refinement, the preceding affine pullback calculation applies.
    Compatibility on further affine refinements is canonical because all local
    maps are restrictions of the original compatible pair.
    \item \textbf{Mapping descent.}
    Mapping descent along the resulting effective epimorphic affine cover and its
    iterated intersections,
    \Ref{lem:mapping-descent-effective-epimorphism}, therefore reconstructs a
    unique map
    \[
        \Spec(C)\longrightarrow Y_C
    \]
    together with all higher homotopies, and identifies its restrictions with
    \(f_0\) and \(f_1\). This proves the asserted equivalence of mapping spaces.
\end{enumerate}
\end{proof}

\begin{theorem}[Global spectral reconstruction reflector]
\label{thm:global-spectral-reconstruction-reflector}
Accessible localization at \(W^{\mathrm{Recon}}_S\) exists and gives a
reflective adjunction
\[
\begin{tikzcd}[column sep=huge]
    (\mathcal H_{\Sp})_{/S}
        \arrow[r,shift left=1.2ex,"L^{\mathrm{Recon}}_S"]
        &
    \bigl((\mathcal H_{\Sp})_{/S}\bigr)^{\mathrm{Recon}}
        \arrow[l,shift left=1.2ex,hook,"\iota_S"]
\end{tikzcd}
\]
An object \(Y\to S\) is local if and only if:
\begin{enumerate}[label=(\roman*)]
    \item for every affine spectral test \(h_A\to S\),
    \[
        \Map_S(h_A,Y)
        \xrightarrow{\simeq}
        \varprojlim\nolimits_n
        \Map_S(h_{\tau_{\leq n}^{\CAlg}A},Y);
    \]
    \item for every nilpotent-excision generator,
    \[
        \Map_S(h_C,Y)
        \xrightarrow{\simeq}
        \Map_S(h_{A_0},Y)
        \times_{\Map_S(h_{A_{01}},Y)}
        \Map_S(h_{A_1},Y).
    \]
\end{enumerate}
Every morphism \(Y\to S\) representable by spectral schemes is
\(L^{\mathrm{Recon}}_S\)-local.
\end{theorem}

\begin{proof}
\leavevmode
\begin{enumerate}[label=\textup{(\arabic*)}]
    \item \textbf{Verify the localization domain.}
    The spectral Zariski slice is presentable, and after fixing universes the
    family \(W^{\mathrm{Recon}}_S\) is a set.
    \item \textbf{Construct the accessible localization.}
    Accessible localization therefore
    exists by \cite[Proposition~5.5.4.15]{LurieHigherToposTheory}.
    \item \textbf{Identify the locality equations.}
    Yoneda identifies locality for the two classes of generators with the
    two displayed equations.
    \item \textbf{Affine Postnikov continuity.}
    Let \(Y\to S\) be representable by spectral schemes. Nilpotent excision is
    \Ref{lem:represented-nilpotent-excision}. For Postnikov continuity, fix an
    affine spectral test
    \[
        T=\Spec(A)\longrightarrow S
    \]
    and put
    \[
        Y_A:=Y\times_ST.
    \]
    This is a spectral scheme. The required comparison is
    \[
        \Map_A(\Spec A,Y_A)
        \longrightarrow
        \varprojlim\nolimits_n
        \Map_A(\Spec\tau_{\leq n}^{\CAlg}A,Y_A).
    \]
    If \(Y_A=\Spec(D)\) is affine, operadic Postnikov completeness gives
    \[
        A\simeq\varprojlim\nolimits_n\tau_{\leq n}^{\CAlg}A,
    \]
    and mapping out of the relative \(A\)-algebra \(D\) gives the desired
    equivalence.
    \item \textbf{Affine-open reduction.}
    For a general spectral scheme \(Y_A\), start with a compatible point of the
    right-hand limit. Its degree-zero member determines a finite principal cover
    of \(\Spec\pi_0(A)\) subordinate to affine opens of \(Y_A\). The same
    elements of \(\pi_0(A)\) determine compatible principal covers of every
    \(\Spec(\tau_{\leq n}^{\CAlg}A)\) and of \(\Spec(A)\). Compatibility of the
    Postnikov family forces the restriction to each principal member to land in
    the same affine target.
    \item \textbf{Čech descent.}
    Choose affine open refinements of the target intersections and finite
    principal refinements of the corresponding source intersections, and continue
    this construction on the higher Čech intersections. On every affine member of the resulting refinement the
    affine Postnikov calculation gives a unique compatible map from the
    corresponding principal open of \(\Spec(A)\). Mapping descent along the
    effective epimorphic affine refinement and its iterated intersections glues
    the maps and all higher homotopies. Thus the relative mapping-space
    comparison is an equivalence for every affine test, and \(Y\to S\) is local.
\end{enumerate}
\end{proof}

\begin{definition}[Reconstruction-local objects]
\label{def:reconstruction-local-objects}
An object of
\[
    \bigl((\mathcal H_{\Sp})_{/S}\bigr)^{\mathrm{Recon}}
\]
is \textbf{reconstruction-local over \(S\)} if it is local for the reflector
\(L^{\mathrm{Recon}}_S\).
\end{definition}

\begin{remark}[Recognition by nilcompleteness and infinitesimal cohesiveness]
\label{rem:reconstruction-local-recognition-dagxiv}
In spectral DAG terminology, the Postnikov-continuity
condition above is Lurie's \textbf{nilcompleteness}. Represented spectral
Deligne--Mumford functors are cohesive and nilcomplete by
\cite[Proposition~2.1.7]{LurieDAGXIV}. Moreover, for a nilcomplete functor
admitting a cotangent complex, the pullback condition with one leg surjective on \(\pi_0\) with nilpotent kernel is
equivalent to infinitesimal cohesiveness by
\cite[Proposition~2.1.13]{LurieDAGXIV}.

These citations are recognition results only. The category
\(\bigl((\mathcal H_{\Sp})_{/S}\bigr)^{\mathrm{Recon}}\) is defined by the explicit localization
\(L^{\mathrm{Recon}}_S\) on arbitrary spectral Zariski stacks; existence of a
cotangent complex, a cohesive structure, or infinitesimal cohesiveness is not
part of the definition; the definition remains the explicit localization above.
\end{remark}

\begin{proposition}[Base change of reconstruction-local objects]
\label{prop:reconstruction-local-base-change}
Let
\[
    p:S'\longrightarrow S
\]
be a morphism of spectral Zariski stacks. Pullback carries
reconstruction-local objects over \(S\) to reconstruction-local objects over
\(S'\):
\[
    p^*:
    \bigl((\mathcal H_{\Sp})_{/S}\bigr)^{\mathrm{Recon}}
    \longrightarrow
    \bigl((\mathcal H_{\Sp})_{/S'}\bigr)^{\mathrm{Recon}}.
\]
This proposition records the stability of the localization equations under
base change.
\end{proposition}

\begin{proof}
\leavevmode
\begin{enumerate}[label=\textup{(\arabic*)}]
    \item \textbf{Slice mapping spaces.}
    Let \(Y\to S\) be reconstruction-local and put
    \[
        Y':=Y\times_SS'.
    \]
    For every affine spectral test \(T\to S'\), slice base change gives a natural
    equivalence
    \[
        \Map_{S'}(T,Y')
        \simeq
        \Map_S(T,Y),
    \]
    where the right-hand side regards \(T\) as an object over \(S\) through
    \(S'\to S\).
    \item \textbf{Postnikov continuity after base change.}
    Apply this equivalence to \(T=\Spec(A)\) and to all of its Postnikov
    truncations. The Postnikov-continuity comparison for \(Y'\) is identified
    with the corresponding comparison for \(Y\), and is therefore an
    equivalence.
    \item \textbf{Nilpotent excision after base change.}
    Likewise, a nilpotent-excision pullback diagram of connective affine tests
    over \(S'\), after composition with \(S'\to S\), is the same algebraic
    pullback diagram regarded over \(S\). The mapping-space comparison for
    \(Y'\) identifies with the corresponding comparison for \(Y\), hence is an
    equivalence. Therefore \(Y'\) satisfies both defining locality equations and
    is reconstruction-local over \(S'\).
\end{enumerate}
\end{proof}

\begin{remark}[Reconstruction locality versus formal-moduli problems]
\label{rem:reconstruction-local-versus-fmp}
The adjective \textbf{reconstruction-local} refers to the centre-free nonlinear
localization of \Ref{thm:global-spectral-reconstruction-reflector}. The centred
Artinian notion of \Ref{def:spectral-formal-moduli-problem} is recovered from
it by the centred Artinian leaf construction below.
\end{remark}

\subsection{Left exactness of the reconstruction localization}

\begin{theorem}[Left exactness of spectral reconstruction localization]
\label{thm:reconstruction-localization-left-exact}
The accessible localization
\[
    L^{\mathrm{Recon}}_S
\]
is left exact. Consequently
\[
    \bigl((\mathcal H_{\Sp})_{/S}\bigr)^{\mathrm{Recon}}
\]
is an \(\infty\)-topos, its fully faithful inclusion into
\((\mathcal H_{\Sp})_{/S}\) preserves all limits, and
\(L^{\mathrm{Recon}}_S\) preserves finite limits.
\end{theorem}

\begin{proof}
\leavevmode
\begin{enumerate}[label=\textup{(\arabic*)}]
    \item \textbf{Localization criterion.}
    Let \(W\) be the strongly saturated class generated by
    \(W^{\mathrm{Recon}}_S\). The left-exactness criterion for a localization of an \(\infty\)-topos,
    \cite[Proposition~6.2.1.1]{LurieHigherToposTheory}, reduces the theorem to
    showing that \(W\) is stable under pullback. It is enough to prove that every generating morphism
    is universally a local equivalence.
    \item \textbf{Pull back a nilpotent-excision generator.}
    First consider a nilpotent-excision generator
    \[
        w_\square:
        h_{A_0}\amalg_{h_{A_{01}}}h_{A_1}
        \longrightarrow
        h_C,
        \qquad
        C\simeq A_0\times_{A_{01}}A_1,
    \]
    and pull it back along an affine morphism \(h_E\to h_C\). Put
    \[
        E_i:=E\otimes_CA_i,
        \qquad
        E_{01}:=E\otimes_CA_{01}.
    \]
    \item \textbf{Record the underlying fibre sequence.}
    The defining pullback of algebras has an underlying fibre sequence
    \[
        C\longrightarrow A_0\oplus A_1\longrightarrow A_{01}.
    \]
    \item \textbf{Extend scalars to the affine pullback.}
    Extension of scalars from \(C\) to \(E\) is exact on module spectra, so it
    produces a fibre sequence
    \[
        E\longrightarrow E_0\oplus E_1\longrightarrow E_{01}.
    \]
    \item \textbf{Recover the algebra pullback.}
    Since limits of unrestricted commutative ring spectra are created on underlying
    spectra, the induced morphism
    \[
        E\longrightarrow E_0\times_{E_{01}}E_1
    \]
    is an equivalence of commutative ring spectra.
    \item \textbf{Base change the nilpotent-surjective degree-zero map.}
    Suppose, for definiteness,
    that \(\pi_0(A_0)\to\pi_0(A_{01})\) is surjective with nilpotent kernel. Connectivity
    of the pushouts gives
    \[
        \pi_0(E_i)
        \simeq
        \pi_0(E)\otimes_{\pi_0(C)}\pi_0(A_i).
    \]
    Thus \(\pi_0(E_0)\to\pi_0(E_{01})\) is again surjective and its kernel is a
    quotient of the scalar extension of the original nilpotent ideal; in
    particular it is nilpotent.
    \item \textbf{Apply connective pullback recognition.}
    By
    \Ref{lem:connective-nilpotent-pullback}, the displayed unrestricted pullback is
    connective.
    \item \textbf{Conclude stability of the generator.}
    Therefore the affine pullback of \(w_\square\) is again a
    nilpotent-excision generator.
    \item \textbf{Pull back the Postnikov-continuity generator.}
    Now consider a Postnikov-continuity generator
    \[
        c_A:
        \operatorname*{colim}_nh_{\tau_{\leq n}^{\CAlg}A}
        \longrightarrow
        h_A
    \]
    and an affine morphism \(h_E\to h_A\), represented by a connective algebra map
    \(A\to E\). Pullback in the spectral Zariski topos preserves colimits, so the
    pullback of \(c_A\) is
    \[
        \operatorname*{colim}_nh_{E_n}
        \longrightarrow
        h_E,
        \qquad
        E_n:=E\otimes_A\tau_{\leq n}^{\CAlg}A.
    \]
    \item \textbf{Define the Postnikov fibre.}
    Let
    \[
        K_n:=\operatorname{fib}(A\to\tau_{\leq n}^{\CAlg}A).
    \]
    \item \textbf{Record its connectivity.}
    Then \(K_n\) is \((n+1)\)-connective.
    \item \textbf{Base change the fibre.}
    Exactness of extension of scalars gives
    \[
        \operatorname{fib}(E\to E_n)
        \simeq
        E\otimes_AK_n,
    \]
    which is again \((n+1)\)-connective because \(E\) is connective.
    \item \textbf{Identify stable truncations.}
    Hence for
    every fixed \(d\),
    \[
        \tau_{\leq d}^{\CAlg}E_n
        \xrightarrow{\simeq}
        \tau_{\leq d}^{\CAlg}E
        \qquad(n\geq d).
    \]
    \item \textbf{Evaluate a reconstruction-local object.}
    If \(Z\) is reconstruction-local, then
    \[
    \begin{aligned}
        \varprojlim\nolimits_nZ(E_n)
        &\simeq
        \varprojlim\nolimits_n\varprojlim\nolimits_dZ(\tau_{\leq d}^{\CAlg}E_n)
        \\
        &\simeq
        \varprojlim\nolimits_d\varprojlim\nolimits_nZ(\tau_{\leq d}^{\CAlg}E_n)
        \\
        &\simeq
        \varprojlim\nolimits_dZ(\tau_{\leq d}^{\CAlg}E)
        \\
        &\simeq
        Z(E).
    \end{aligned}
    \]
    \item \textbf{Conclude local equivalence.}
    Thus every affine pullback of \(c_A\) is an
    \(L^{\mathrm{Recon}}_S\)-equivalence.
    \item \textbf{Pass to arbitrary pullbacks.}
    Finally let \(P\) be an arbitrary object over the target of a generator.
    Affine representables generate the corresponding slice of the spectral Zariski
    \(\infty\)-topos under colimits. Write \(P\) as such a colimit. Pullback in
    an \(\infty\)-topos preserves colimits, so the pullback of the generator along
    \(P\) is the colimit in the arrow category of its affine pullbacks. The
    localization is a left adjoint and therefore sends this colimit to a colimit of
    equivalences. Hence every generator is universally local. The class of
    morphisms all of whose pullbacks lie in \(W\) is strongly saturated and
    contains the generators, so it contains \(W\). Therefore the local-equivalence
    class is pullback-stable, which proves left exactness.
\end{enumerate}
\end{proof}

\subsection{The reconstruction localization and differential cohesion}

Let
\[
    \iota_S:
    \bigl((\mathcal H_{\Sp})_{/S}\bigr)^{\mathrm{Recon}}
    \hookrightarrow
    (\mathcal H_{\Sp})_{/S}
\]
denote the fully faithful inclusion.

\begin{definition}[Reconstruction localization endofunctor]
\label{def:reconstruction-localization-endofunctor}
Define
\[
    \operatorname{Recon}_S
    :=
    \iota_S L^{\mathrm{Recon}}_S
    :
    (\mathcal H_{\Sp})_{/S}
    \longrightarrow
    (\mathcal H_{\Sp})_{/S}.
\]
We call \(\operatorname{Recon}_S\) the \textbf{reconstruction localization endofunctor} over \(S\).
\end{definition}

\begin{proposition}[Idempotence of the reconstruction localization]
\label{prop:reconstruction-localization-idempotence}
The endofunctor \(\operatorname{Recon}_S\) is an idempotent left-exact localization endofunctor.
Its fixed objects are precisely the reconstruction-local objects over \(S\). The unit fits into the idempotence diagram
\[
\begin{tikzcd}[column sep=huge]
    Y \arrow[r,"\eta_Y"]
        &
    \operatorname{Recon}_S Y
        \arrow[r,"\operatorname{Recon}_S(\eta_Y)"]
        &
    \operatorname{Recon}_S^2Y ,
\end{tikzcd}
\]
where the second arrow is an equivalence. The first arrow is universal from
\(Y\) to an object satisfying Postnikov continuity and nilpotent excision.
\end{proposition}

\begin{proof}
\leavevmode
\begin{enumerate}[label=\textup{(\arabic*)}]
    \item \textbf{Use the reflective-localization identities.}
    Idempotence and the fixed-point characterization are the reflective-localization identities
    for the adjunction
    \[
        L^{\mathrm{Recon}}_S\dashv\iota_S.
    \]
    \item \textbf{Apply left exactness.}
    Left exactness is \Ref{thm:reconstruction-localization-left-exact}.
    \item \textbf{Apply the localization universal property.}
    The universal
    property is the defining universal property of localization at
    \(W^{\mathrm{Recon}}_S\), whose local objects are characterized by
    \Ref{thm:global-spectral-reconstruction-reflector}.
\end{enumerate}
\end{proof}

\begin{remark}[Two left-exact localization constructions]
\label{rem:infinitesimal-shape-versus-formal-reconstruction}
The infinitesimal-shape modality
\[
    \Im=i_*i^*
\]
of Chapter~\Ref{chap:differentiation-spectral-doctrine} acts on the
infinitesimal Zariski topos and records the domain endpoint of an
infinitesimal arrow. The reconstruction localization
\[
    \operatorname{Recon}_S=\iota_SL^{\mathrm{Recon}}_S
\]
acts on the ordinary spectral Zariski slice over \(S\) and enforces
Postnikov continuity together with nilpotent excision. The centred Artinian
leaf construction below provides the comparison relevant to this chapter.
\end{remark}


The centre-free nonlinear semantics is now available: reconstruction-local objects are
the fixed points of a left-exact localization imposing Postnikov continuity and
nilpotent excision. Because locality is stable under base change, we can choose an
affine centre inside such an object and restrict to its bounded Artinian thickenings.
This returns the global theory to the infinitesimal endpoint comparison of the preceding
chapter.

\section{Centred Artinian leaves and tangent recovery}
\label{sec:centres-artinian-leaves}


Choose an affine centre
\[
    U=\Spec(B)\longrightarrow V=\Spec(R)\longrightarrow S
\]
and a \(U\)-point \(u:U\to Y\) of a reconstruction-local object. Restricting \(Y\) to
coherent Artinian thickenings produces a centred formal-moduli problem; on each test,
its value is the fibre at \(u\) of the endpoint restriction underlying differential
cohesion.

The passage from nonlinear reconstruction to tangent semantics is summarized by
\[
\begin{tikzcd}[column sep=huge,row sep=large]
    \text{Artinian thickening }j_A:U\to T_A
        \arrow[r,"\text{endpoint evaluation}"]
        \arrow[d,"\text{split square-zero restriction}"']
    &
    \Map_V(T_A,Y_V)\to\Map_V(U,Y_V)
        \arrow[d,"\operatorname{fib}_u"]
    \\
    \text{connective coefficient }M[n]
        \arrow[r,"\text{centred lifting fibre}"']
    &
    \mathsf T^n_{Y/V,u}(M).
\end{tikzcd}
\]
The upper row retains the nonlinear finite infinitesimal test, while the lower row
extracts the lifting fibre in a coefficient direction. Nilpotent excision turns these
fibres into an \(\Omega\)-spectrum and hence into a stable tangent functor.

\subsection{Centred Artinian leaves}

\begin{definition}[Centred Artinian leaf]
\label{def:centred-artinian-leaf}
Let
\[
    V=\Spec(R)\longrightarrow S,
    \qquad
    U=\Spec(B)\longrightarrow V
\]
be affine spectral schemes with \(B\) coherent. Let \(Y\to S\) be a spectral
Zariski stack, and let
\[
    u:U\longrightarrow Y
\]
be a chosen \(U\)-point over \(S\). For
\(A\to B\in\operatorname{Art}^{\mathrm{sp}}_{B/R}\), define
\[
    \operatorname{Leaf}^{\mathrm{Art}}_{U/V,u}(Y)(A\to B)
    :=
    \operatorname{fib}_u
    \left(
        \Map_S(\Spec A,Y)
        \longrightarrow
        \Map_S(U,Y)
    \right).
\]

Equivalently, it is the upper-left corner of the Cartesian square
\[
\begin{tikzcd}[column sep=large,row sep=large]
    \operatorname{Leaf}^{\mathrm{Art}}_{U/V,u}(Y)(A\to B)
        \arrow[r] \arrow[d]
        &
    \Map_S(\Spec A,Y) \arrow[d]
    \\
    * \arrow[r,"u"']
        & \Map_S(U,Y).
\end{tikzcd}
\]
\end{definition}

\begin{proposition}[Centred Artinian leaves as differential-cohesive endpoint fibres]
\label{prop:artinian-leaf-endpoint-fibre}
Let \(A\to B\) be spectral Artinian and let
\[
    j_A:
    U=\Spec(B)\longrightarrow T_A=\Spec(A)
\]
be the associated finite infinitesimal test of
\Ref{prop:artinian-as-finite-infinitesimal-tests}. Put
\[
    Y_V:=Y\times_SV.
\]
Under the endpoint formulas of
Chapter~\Ref{chap:differentiation-spectral-doctrine}, the restriction map
\[
    \Map_V(T_A,Y_V)
    \longrightarrow
    \Map_V(U,Y_V)
\]
is the comparison obtained by evaluating
\[
    i_!Y_V\longrightarrow i_*Y_V
\]
at the infinitesimal arrow \(j_A\) and then restricting to the fibre over the
fixed structural maps to \(V\). Consequently
\[
    \operatorname{Leaf}^{\mathrm{Art}}_{U/V,u}(Y)(A\to B)
\]
is the homotopy fibre at the chosen point \(u\) of this differential-cohesive endpoint
restriction.
\end{proposition}

\begin{proof}
\leavevmode
\begin{enumerate}[label=\textup{(\arabic*)}]
    \item \textbf{Evaluate the endpoint Kan extensions.}
    For an infinitesimal arrow \(j:U'\to T'\), the endpoint Kan-extension formulas
    of \Ref{prop:endpoint-kan-extensions} give
    \[
        (i_!F)(j)\simeq F(T'),
        \qquad
        (i_*F)(j)\simeq F(U').
    \]
    \item \textbf{Identify the endpoint comparison.}
    The natural comparison \(i_!F\to i_*F\) is induced pointwise by restriction
    along \(U'\to T'\).
    \item \textbf{Specialize to the Artinian arrow.}
    Apply this with \(F=Y_V\) and \(j=j_A\).
    \item \textbf{Pass to relative mapping spaces.}
    Taking the
    homotopy fibres of the two evaluation spaces over their prescribed maps to
    \(V\) gives the relative restriction
    \[
        \Map_V(T_A,Y_V)\longrightarrow\Map_V(U,Y_V).
    \]
    \item \textbf{Apply slice base change.}
    The slice base-change adjunction identifies these with the corresponding
    mapping spaces over \(S\) used in
    \Ref{def:centred-artinian-leaf}.
    \item \textbf{Take the centred fibre.}
    Taking the further homotopy fibre at the
    chosen point \(u\) gives the asserted centred Artinian leaf.
\end{enumerate}
\end{proof}

\begin{proposition}[Centred Artinian leaves are formal-moduli problems]
\label{prop:reconstruction-local-artinian-leaf-fmp}
If \(Y\to S\) is reconstruction-local, then
\[
    \operatorname{Leaf}^{\mathrm{Art}}_{U/V,u}(Y)
    \in
    \operatorname{FMP}^{\mathrm{sp}}_{B/R}.
\]
\end{proposition}

\begin{proof}
\leavevmode
\begin{enumerate}[label=\textup{(\arabic*)}]
    \item \textbf{Centre condition.}
    At the identity augmentation \(B\to B\), the defining restriction map is the
    identity of \(\Map_S(U,Y)\), so its homotopy fibre over \(u\) is contractible.
    \item \textbf{Structured square-zero excision.}
    Let
    \[
        A_\sigma
        =
        A\times_{B\oplus M[1]}B
    \]
    be a Schlessinger square with \(A\to B\) Artinian and \(M\) connective
    coherent. The zero section
    \[
        B\longrightarrow B\oplus M[1]
    \]
    is an isomorphism on \(\pi_0\), hence this is a nilpotent-excision pullback.
    Reconstruction locality therefore gives
    \[
        Y(A_\sigma)
        \simeq
        Y(A)
        \times_{Y(B\oplus M[1])}
        Y(B).
    \]
    Taking the homotopy fibre of this equivalence over the compatible chosen point
    \(u\in Y(B)\), and using pullback pasting, gives exactly the Schlessinger
    condition for \(\operatorname{Leaf}^{\mathrm{Art}}_{U/V,u}(Y)\). Apply
    \Ref{def:spectral-formal-moduli-problem}.
\end{enumerate}
\end{proof}



\subsection{Tangent recovery from split square-zero lifting}
\label{sec:reconstruction-tangent-semantics}


The centred Artinian leaf already identifies nonlinear reconstruction with actual
infinitesimal lifting at a fixed point. We now restrict those tests to split square-zero
augmentations and stabilize in the coefficient variable.

For connective coefficients, \(\mathsf T^n_{Y/V,u}(M)\) is literally the centred lifting
space across \(B\oplus M[n]\to B\). Nilpotent excision supplies both the loop
equivalences and right exactness. The Spanier--Whitehead universal property then extends
the resulting spectrum-valued functor from connective to eventually connective
coefficients. No extension to all unbounded modules is claimed without an additional
continuity theorem.

\subsubsection{Actual split-square-zero lifting on connective coefficients}

\begin{definition}[Actual split-square-zero tangent levels]
\label{def:split-square-zero-tangent-levels}
Let \(Y\to S\) be reconstruction-local and let
\[
    V=\Spec(R)\longrightarrow S,
    \qquad
    U=\Spec(B)\longrightarrow V
\]
be a connective affine centre. Let
\[
    u:U\longrightarrow Y
\]
be a chosen \(U\)-point over \(S\). For a connective \(B\)-module \(M\) and
\(n\geq0\), define
\[
    \mathsf T^n_{Y/V,u}(M)
    :=
    \operatorname{fib}_u
    \left(
        \Map_V
        \bigl(\Spec(B\oplus M[n]),Y\times_SV\bigr)
        \longrightarrow
        \Map_V(U,Y\times_SV)
    \right).
\]
The map is induced by the split square-zero augmentation
\(B\oplus M[n]\to B\). The split algebra section
\[
    B\longrightarrow B\oplus M[n]
\]
induces a retraction
\[
    \Spec(B\oplus M[n])\longrightarrow U,
\]
and postcomposition of this retraction with the chosen point \(u\) defines a
canonical basepoint of the displayed homotopy fibre. Thus
\(\mathsf T^n_{Y/V,u}(M)\) is canonically pointed, functorially in \(M\).
\end{definition}

\begin{proposition}[Centred Artinian leaves recover the coherent split-square-zero tangent levels]
\label{prop:artinian-leaf-tangent-levels}
Let \(Y\to S\) be reconstruction-local and let
\[
    U=\Spec(B)\longrightarrow V=\Spec(R)\longrightarrow S
\]
be an affine centre with \(B\) coherent. Let
\[
    u:U\longrightarrow Y
\]
be a chosen \(U\)-point over \(S\). For every connective coherent \(B\)-module
\(M\) and every \(n\geq0\), the split square-zero augmentation
\[
    B\oplus M[n]\longrightarrow B
\]
is a spectral Artinian extension, and there is a canonical equivalence
\[
    \operatorname{Leaf}^{\mathrm{Art}}_{U/V,u}(Y)
    \bigl(B\oplus M[n]\to B\bigr)
    \simeq
    \mathsf T^n_{Y/V,u}(M),
\]
where the right-hand side is the split-square-zero lifting space defined in
\Ref{def:split-square-zero-tangent-levels}. Under these equivalences the
loop maps induced by
\[
    B\oplus M[n]
    \simeq
    B\times_{B\oplus M[n+1]}B
\]
agree. Hence stabilizing the centred Artinian leaf on connective coherent
coefficients gives the restriction of the recovered tangent functor
constructed below.
\end{proposition}

\begin{proof}
\leavevmode
\begin{enumerate}[label=\textup{(\arabic*)}]
    \item \textbf{Artinianity of split square-zero augmentations.}
    Since \(M\) is coherent, it is connective, almost perfect, and eventually
    coconnective. The split square-zero augmentation is an isomorphism on
    \(\pi_0\) for \(n>0\), and has square-zero degree-zero kernel for \(n=0\).
    Its fibre is \(M[n]\), hence is eventually coconnective. Almost finite
    presentation follows from \Ref{lem:square-zero-almost-finite}. Thus it is
    spectral Artinian.
    \item \textbf{Identify the leaf with the lifting space.}
    The slice base-change adjunction gives natural equivalences
    \[
        \Map_S\bigl(\Spec(B\oplus M[n]),Y\bigr)
        \simeq
        \Map_V
        \bigl(\Spec(B\oplus M[n]),Y\times_SV\bigr)
    \]
    and likewise for \(U\). Taking the homotopy fibre over the chosen point \(u\)
    identifies the defining value of the centred Artinian leaf with
    \(\mathsf T^n_{Y/V,u}(M)\). Both loop maps are obtained by applying the same
    nilpotent-excision pullback identity for consecutive split square-zero
    algebras, so the equivalences are compatible with stabilization.
\end{enumerate}
\end{proof}

\begin{lemma}[Retractive self-intersections are loop objects]
\label{lem:retractive-self-intersection-loop}
Let
\[
    X\xrightarrow{s}E\xrightarrow{r}X,
    \qquad r\circ s\simeq\id_X,
\]
be a retractive diagram of spaces and let \(x\in X\). Give
\(X\times_E X\) the two maps to \(E\) induced by \(s\). Then the homotopy
fibre over \(x\) of the first projection
\[
    X\times_E X\longrightarrow X
\]
is canonically equivalent to
\[
    \Omega_{s(x)}\operatorname{fib}_x(r:E\to X).
\]
The equivalence is natural in retractive diagrams.
\end{lemma}

\begin{proof}
\leavevmode
\begin{enumerate}[label=\textup{(\arabic*)}]
    \item \textbf{Base change to the chosen point.}
    Base change the retractive diagram along the point \(x:*\to X\).
    \item \textbf{Reduce to a pointed space.}
    It reduces
    the statement to a pointed space \((F,*)\), where \(F=\operatorname{fib}_x(r)\),
    and to the homotopy fibre of the basepoint inclusion
    \[
        *\longrightarrow F
    \]
    over the basepoint.
    \item \textbf{Identify the loop space.}
    That homotopy fibre is the based loop space
    \(\Omega F\).
    \item \textbf{Record naturality.}
    Naturality is inherited from pullback.
\end{enumerate}
\end{proof}

\begin{proposition}[Square-zero stabilization and right exactness]
\label{prop:split-tangent-omega-spectrum}
For every connective \(B\)-module \(M\), the canonical split square-zero
pullback squares give natural equivalences
\[
    \mathsf T^n_{Y/V,u}(M)
    \xrightarrow{\simeq}
    \Omega\mathsf T^{n+1}_{Y/V,u}(M),
    \qquad n\geq0.
\]
These equivalences are compatible with identities, morphisms of coefficients,
and iterated suspension. Hence the pointed functors
\[
    \mathsf T^n_{Y/V,u}:
    (\Mod_B)_{\geq0}
    \longrightarrow
    \mathcal S_*
\]
form a canonical \(\Omega\)-spectrum object in
\[
    \Fun\bigl((\Mod_B)_{\geq0},\mathcal S_*\bigr).
\]
Equivalently, they determine a canonical functor
\[
    \mathsf T^{\mathrm{cn}}_{Y/V,u}:
    (\Mod_B)_{\geq0}
    \longrightarrow
    \Sp.
\]

The resulting functor is zero-preserving and right exact. Explicitly, if
\[
    M'\longrightarrow M\longrightarrow M''
\]
is a cofibre sequence of connective \(B\)-modules, then
\[
    \mathsf T^{\mathrm{cn}}_{Y/V,u}(M')
    \longrightarrow
    \mathsf T^{\mathrm{cn}}_{Y/V,u}(M)
    \longrightarrow
    \mathsf T^{\mathrm{cn}}_{Y/V,u}(M'')
\]
is a fibre and hence also a cofibre sequence of spectra.
\end{proposition}

\begin{proof}
\leavevmode
\begin{enumerate}[label=\textup{(\arabic*)}]
    \item \textbf{Base change the reconstruction-local target.}
    By \Ref{prop:reconstruction-local-base-change},
    \[
        Y_V:=Y\times_SV
    \]
    is reconstruction-local over \(V\).
    \item \textbf{Form the coefficient pullback square.}
    For every connective \(B\)-module \(M\)
    and every \(n\geq0\), exactness in \(\Mod_B\) gives the Cartesian square
    \[
    \begin{tikzcd}[column sep=large,row sep=large]
        M[n] \arrow[r] \arrow[d]
            & 0 \arrow[d]
        \\
        0 \arrow[r]
            & M[n+1],
    \end{tikzcd}
    \]
    since \(M[n]\simeq\Omega M[n+1]\).
    \item \textbf{Apply the split square-zero functor.}
    The split square-zero functor
    \[
        \Mod_B\longrightarrow\CAlg(\Sp)_{B//B},
        \qquad
        N\longmapsto B\oplus N,
    \]
    preserves limits by
    \Ref{prop:split-square-zero-functoriality}.
    \item \textbf{Obtain the algebra pullback square.}
    Applying it to the displayed
    module square gives the canonical pullback square of connective commutative
    \(B\)-algebras
    \[
    \begin{tikzcd}[column sep=large,row sep=large]
        B\oplus M[n] \arrow[r] \arrow[d]
            & B \arrow[d,"d_0"]
        \\
        B \arrow[r,"d_0"']
            & B\oplus M[n+1].
    \end{tikzcd}
    \]
    \item \textbf{Identify the square-zero self-intersection.}
    Thus the identity
    \[
        B\oplus M[n]
        \simeq
        B\times_{B\oplus M[n+1]}B
    \]
    is a functorial square-zero pullback identity.
    \item \textbf{Record compatibility under iteration.}
    The canonical comparisons
    for successive values of \(n\) are compatible under iterated pullback by
    limit functoriality and pullback pasting.
    \item \textbf{Loop comparison.}
    The lower horizontal morphism is an isomorphism on \(\pi_0\) for every
    \(n\geq0\), and hence is a nilpotent-excision leg. Since \(Y\to S\) is
    reconstruction-local, evaluation on the displayed algebraic pullback gives a
    homotopy pullback square of mapping spaces. After passing to the slice over
    \(V\) and taking the homotopy fibre over the chosen point \(u\), the
    retractive self-intersection lemma
    \Ref{lem:retractive-self-intersection-loop} gives a natural equivalence
    \[
        \mathsf T^n_{Y/V,u}(M)
        \xrightarrow{\simeq}
        \Omega\mathsf T^{n+1}_{Y/V,u}(M).
    \]
    \item \textbf{Spectrum coherence.}
    The split square-zero pullback construction is itself functorial, and the
    comparison between a two-step and an iterated three-step pullback is the
    canonical associativity comparison for finite limits. Thus the preceding
    natural equivalences satisfy the compatibility required for an
    \(\Omega\)-spectrum object in
    \[
        \Fun\bigl((\Mod_B)_{\geq0},\mathcal S_*\bigr).
    \]
    Equivalently, they determine the functor
    \[
        \mathsf T^{\mathrm{cn}}_{Y/V,u}:
        (\Mod_B)_{\geq0}\longrightarrow\Sp.
    \]
    The spectrum structure is induced by the functorial square-zero pullback
    identities.
    \item \textbf{Coefficient cofiber square.}
    We next prove right exactness. Let
    \[
        M'\longrightarrow M\longrightarrow M''
    \]
    be a cofibre sequence of connective \(B\)-modules. For every \(n\geq0\),
    exactness in \(\Mod_B\) gives
    \[
        M'[n]\simeq\operatorname{fib}\bigl(M[n]\to M''[n]\bigr).
    \]
    Apply the limit-preserving split square-zero functor of
    \Ref{prop:split-square-zero-functoriality} to the corresponding Cartesian
    square
    \[
        M'[n]\simeq M[n]\times_{M''[n]}0.
    \]
    This gives the canonical Cartesian square
    \[
    \begin{tikzcd}[column sep=large,row sep=large]
        B\oplus M'[n] \arrow[r] \arrow[d]
            & B \arrow[d,"d_0"]
        \\
        B\oplus M[n] \arrow[r]
            & B\oplus M''[n].
    \end{tikzcd}
    \]
    \item \textbf{Nilpotent-excision hypothesis.}
    The lower horizontal morphism is surjective on \(\pi_0\). For \(n>0\) this
    is immediate because all three shifted coefficient modules have trivial
    degree-zero homotopy. For \(n=0\), the long exact sequence of the connective
    cofibre sequence gives
    \[
        \pi_0(M)\twoheadrightarrow\pi_0(M'').
    \]
    The kernel of the corresponding algebra map lies in the square-zero
    coefficient ideal, and is therefore square-zero. Thus in every degree the
    displayed square is a nilpotent-excision square.
    \item \textbf{Right exactness.}
    Reconstruction locality gives a homotopy pullback
    \[
    \begin{tikzcd}[column sep=large,row sep=large]
        Y(B\oplus M'[n]) \arrow[r] \arrow[d]
            & Y(B) \arrow[d]
        \\
        Y(B\oplus M[n]) \arrow[r]
            & Y(B\oplus M''[n]),
    \end{tikzcd}
    \]
    where the relative mapping-space notation over \(V\) and \(S\) is suppressed.
    Taking the compatible homotopy fibres over \(u\) and using pullback pasting
    gives
    \[
        \mathsf T^n_{Y/V,u}(M')
        \simeq
        \operatorname{fib}
        \left(
            \mathsf T^n_{Y/V,u}(M)
            \longrightarrow
            \mathsf T^n_{Y/V,u}(M'')
        \right).
    \]
    These fibre identifications are natural in \(n\) and commute with the
    \(\Omega\)-spectrum comparisons constructed above. Since limits in spectra
    are computed levelwise, they assemble to a fibre sequence
    \[
        \mathsf T^{\mathrm{cn}}_{Y/V,u}(M')
        \longrightarrow
        \mathsf T^{\mathrm{cn}}_{Y/V,u}(M)
        \longrightarrow
        \mathsf T^{\mathrm{cn}}_{Y/V,u}(M'').
    \]
    The target category \(\Sp\) is stable, so the same sequence is a cofibre
    sequence.
    \item \textbf{Zero coefficient.}
    For the zero coefficient,
    \[
        B\oplus0\simeq B,
    \]
    and the defining restriction map is the identity. Hence
    \[
        \mathsf T^{\mathrm{cn}}_{Y/V,u}(0)\simeq0,
    \]
    A zero-preserving functor from the connective prestable category
    \((\Mod_B)_{\geq0}\) to a stable category which carries connective
    cofibre sequences to cofibre sequences preserves finite colimits. Therefore
    \(\mathsf T^{\mathrm{cn}}_{Y/V,u}\) is right exact.
\end{enumerate}
\end{proof}


\begin{proposition}[Postnikov continuity on connective coefficients]
\label{prop:split-tangent-postnikov-continuity}
For every connective \(B\)-module \(M\), there is a canonical equivalence
\[
    \mathsf T^{\mathrm{cn}}_{Y/V,u}(M)
    \xrightarrow{\simeq}
    \varprojlim\nolimits_k
    \mathsf T^{\mathrm{cn}}_{Y/V,u}(\tau_{\leq k}M).
\]
\end{proposition}

\begin{proof}
\leavevmode
\begin{enumerate}[label=\textup{(\arabic*)}]
    \item \textbf{Levelwise reduction.}
    It is enough to prove the assertion levelwise before stabilization. Fix
    \(n\geq0\). Reconstruction locality gives
    \[
        Y(B\oplus M[n])
        \simeq
        \varprojlim\nolimits_d
        Y\bigl(\tau_{\leq d}^{\CAlg}(B\oplus M[n])\bigr)
    \]
    and
    \[
        Y(B)
        \simeq
        \varprojlim\nolimits_dY(\tau_{\leq d}^{\CAlg}B).
    \]
    \item \textbf{Stabilization of truncations.}
    For fixed \(d\), use the canonical Postnikov truncation
    \[
        M\longrightarrow\tau_{\leq k}M.
    \]
    It induces a morphism of split square-zero algebras
    \[
        B\oplus M[n]
        \longrightarrow
        B\oplus\tau_{\leq k}M[n].
    \]
    Its underlying fibre is \(\tau_{\geq k+1}M[n]\), hence is
    \((k+n+1)\)-connective. Therefore, once \(k\) is sufficiently large relative
    to \(d\) and \(n\), it induces an equivalence
    \[
        \tau_{\leq d}^{\CAlg}(B\oplus M[n])
        \xrightarrow{\simeq}
        \tau_{\leq d}^{\CAlg}
        \bigl(B\oplus\tau_{\leq k}M[n]\bigr).
    \]
    Thus, for every fixed \(d\), the truncated square-zero evaluation stabilizes
    along the Postnikov tower of \(M\).
    \item \textbf{Define the truncated centred fibres.}
    For brevity let \(u_d\) denote the restriction of the chosen centre point
    \(u\) along
    \[
        \Spec(\tau_{\leq d}^{\CAlg}B)\longrightarrow U
    \]
    and put
    \[
        \Phi_d(N)
        :=
        \operatorname{fib}_{u_d}
        \left(
            Y\bigl(\tau_{\leq d}^{\CAlg}(B\oplus N[n])\bigr)
            \longrightarrow
            Y(\tau_{\leq d}^{\CAlg}B)
        \right).
    \]
    \item \textbf{Use reconstruction locality.}
    Reconstruction locality gives
    \[
        \mathsf T^n_{Y/V,u}(N)
        \simeq
        \varprojlim\nolimits_d\Phi_d(N)
    \]
    for \(N=M\) and for every \(N=\tau_{\leq k}M\).
    \item \textbf{Use eventual constancy.}
    By the preceding
    stabilization calculation, for each fixed \(d\) the tower
    \(\Phi_d(\tau_{\leq k}M)\) is eventually constant with value
    \(\Phi_d(M)\).
    \item \textbf{Interchange the inverse limits.}
    Therefore
    \[
    \begin{aligned}
        \varprojlim\nolimits_k
        \mathsf T^n_{Y/V,u}(\tau_{\leq k}M)
        &\simeq
        \varprojlim\nolimits_k\varprojlim\nolimits_d
        \Phi_d(\tau_{\leq k}M)
        \\
        &\simeq
        \varprojlim\nolimits_d\varprojlim\nolimits_k
        \Phi_d(\tau_{\leq k}M)
        \\
        &\simeq
        \varprojlim\nolimits_d\Phi_d(M)
        \\
        &\simeq
        \mathsf T^n_{Y/V,u}(M).
    \end{aligned}
    \]
    \item \textbf{Identify the stated comparison.}
    Inverse limits commute because both are limits in spaces. Reversing the
    displayed canonical equivalence gives the comparison stated in the
    proposition.
    \item \textbf{Promote the comparison to spectra.}
    The equivalences are natural in \(M\) and commute with the loop
    identifications of \Ref{prop:split-tangent-omega-spectrum}; they therefore
    assemble to the asserted equivalence of spectra.
\end{enumerate}
\end{proof}


\begin{remark}[Accessibility is a separate finiteness property]
\label{rem:tangent-accessibility-separate}
Accessibility of the tangent functor requires a separate finiteness
hypothesis on the nonlinear target. For example, an accessibility or
local almost finite presentation condition can be used to prove that
square-zero evaluation preserves the relevant filtered colimits.
\end{remark}

\subsubsection{The eventually connective stable envelope}

\begin{definition}[Eventually connective coefficient category]
\label{def:eventually-connective-coefficients}
For a connective commutative ring spectrum \(B\), define
\[
    \Mod_B^{+}
    :=
    \bigcup_{r\in\mathbb Z}(\Mod_B)_{\geq r}
    \subseteq
    \Mod_B.
\]
Thus \(\Mod_B^{+}\) is the full stable subcategory of \(B\)-modules which are
bounded below, with the lower bound allowed to depend on the object.
\end{definition}

\begin{proposition}[Spanier--Whitehead envelope of connective modules]
\label{prop:spanier-whitehead-connective-modules}
There is a canonical equivalence of stable \(\infty\)-categories
\[
    \operatorname{SW}\bigl((\Mod_B)_{\geq0}\bigr)
    \xrightarrow{\simeq}
    \Mod_B^{+}.
\]
For every stable \(\infty\)-category \(\mathcal D\), restriction along the
connective inclusion induces an equivalence
\[
    \Fun^{\mathrm{ex}}(\Mod_B^{+},\mathcal D)
    \xrightarrow{\simeq}
    \Fun^{\mathrm{rex}}
    \bigl((\Mod_B)_{\geq0},\mathcal D\bigr),
\]
where \(\mathrm{rex}\) denotes right-exact functors, equivalently
zero-preserving finite-colimit-preserving functors.
\end{proposition}

\begin{proof}
\leavevmode
\begin{enumerate}[label=\textup{(\arabic*)}]
    \item \textbf{Identify the stabilization.}
    The connective aisle \((\Mod_B)_{\geq0}\) is a prestable
    \(\infty\)-category: it is pointed, admits finite colimits, and its suspension
    functor is fully faithful. Its Spanier--Whitehead stabilization is obtained by
    formally inverting suspension. The inclusion
    \[
        (\Mod_B)_{\geq0}\hookrightarrow\Mod_B
    \]
    therefore induces an exact functor from the stabilization to \(\Mod_B\). Its
    essential image consists precisely of the finite desuspensions of connective
    modules, namely the eventually connective modules. Full faithfulness follows
    from the stabilization mapping formula and the fact that suspension is already
    invertible in \(\Mod_B\). This proves the first equivalence.
    \item \textbf{Universal property.}
    The second assertion is the universal property of the
    Spanier--Whitehead stabilization: for a pointed \(\infty\)-category with
    finite colimits and a stable target, exact functors out of its stabilization
    are equivalent to right-exact functors out of the original category. See
    \cite{LurieHigherAlgebra}.
\end{enumerate}
\end{proof}

\begin{theorem}[Eventually connective tangent functor]
\label{thm:reconstruction-tangent-spectrum}
Let \(Y\to S\) be reconstruction-local and let
\[
    U=\Spec(B)\longrightarrow V=\Spec(R)\longrightarrow S
\]
be a connective affine centre. Let
\[
    u:U\to Y
\]
be a chosen \(U\)-point over \(S\). There is a canonical exact functor
\[
    \mathsf T^{\mathrm{sp},+}_{Y/V,u}:
    \Mod_B^{+}
    \longrightarrow
    \Sp
\]
whose restriction to connective modules is
\(\mathsf T^{\mathrm{cn}}_{Y/V,u}\). It is characterized, through a contractible
space of choices, by this property.

For every connective coefficient \(M\), its value is therefore the actual
stabilized split-square-zero lifting spectrum of
\Ref{def:split-square-zero-tangent-levels}, and that connective restriction is
Postnikov-continuous by
\Ref{prop:split-tangent-postnikov-continuity}.
\end{theorem}

\begin{proof}
\leavevmode
\begin{enumerate}[label=\textup{(\arabic*)}]
    \item \textbf{Use right exactness on connective modules.}
    By \Ref{prop:split-tangent-omega-spectrum},
    \[
        \mathsf T^{\mathrm{cn}}_{Y/V,u}:
        (\Mod_B)_{\geq0}\to\Sp
    \]
    is right exact.
    \item \textbf{Apply the Spanier--Whitehead universal property.}
    Apply the universal property of
    \Ref{prop:spanier-whitehead-connective-modules} with
    \(\mathcal D=\Sp\).
    \item \textbf{Construct the exact extension.}
    It supplies a unique exact extension to
    \(\Mod_B^+\), and the space of such extensions is contractible.
    \item \textbf{Identify the remaining assertions.}
    The final
    statements are the defining restriction and
    \Ref{prop:split-tangent-postnikov-continuity}.
\end{enumerate}
\end{proof}

\begin{proposition}[Cotangent representation of the represented tangent functor]
\label{prop:represented-tangent-cotangent}
For the represented target \(U=\Spec(B)\to V=\Spec(R)\) with chosen
\(U\)-point \(\id_U:U\to U\), there is a natural equivalence of exact functors on eventually
connective modules
\[
    \mathsf T^{\mathrm{sp},+}_{U/V}(M)
    \simeq
    \map_{\Mod_B}(L_{B/R},M),
    \qquad
    M\in\Mod_B^+.
\]
The displayed equivalence concerns \(\Mod_B^+\). The mapping-spectrum
functor on the right extends canonically to all of \(\Mod_B\).
\end{proposition}

\begin{proof}
\leavevmode
\begin{enumerate}[label=\textup{(\arabic*)}]
    \item \textbf{Identify the connective lifting spaces.}
    For connective \(M\), the \(n\)-th space of the actual lifting spectrum is the
    space of \(R\)-algebra sections of
    \[
        B\oplus M[n]\longrightarrow B
    \]
    restricting to the identity.
    \item \textbf{Apply cotangent classification.}
    Cotangent classification identifies it with
    \[
        \Omega^\infty
        \map_{\Mod_B}(L_{B/R},M[n]).
    \]
    \item \textbf{Compare loops with suspension.}
    The loop identifications used in
    \Ref{prop:split-tangent-omega-spectrum} agree with suspension in the mapping
    spectrum, so
    \[
        \mathsf T^{\mathrm{cn}}_{U/V}(M)
        \simeq
        \map_{\Mod_B}(L_{B/R},M)
    \]
    \item \textbf{Obtain the connective mapping-spectrum formula.}
    for connective \(M\).
    \item \textbf{Use exactness of the mapping-spectrum functor.}
    The mapping-spectrum functor on the right is exact on
    all \(B\)-modules.
    \item \textbf{Restrict to eventually connective modules.}
    Restrict it to \(\Mod_B^+\); both sides are then exact
    extensions of the same right-exact functor on connective modules.
    \item \textbf{Apply the stabilization universal property.}
    The
    Spanier--Whitehead universal property identifies them uniquely.
\end{enumerate}
\end{proof}

\begin{definition}[Relative tangent functor of a centred map]
\label{def:relative-tangent-semantics-centre}
Let \(Y\to S\) be reconstruction-local and let
\[
    U=\Spec(B)\longrightarrow V=\Spec(R)\longrightarrow S
\]
be a connective affine centre. Let
\[
    u:U\longrightarrow Y
\]
be a chosen \(U\)-point over \(S\). Postcomposition with \(u\) gives a natural
morphism of exact functors on \(\Mod_B^+\)
\[
    \map_{\Mod_B}(L_{B/R},-)|_{\Mod_B^+}
    \longrightarrow
    \mathsf T^{\mathrm{sp},+}_{Y/V,u}.
\]
Define the \textbf{relative tangent functor of the centre} by
\[
    \mathsf T^{\mathrm{sp},+}_{U/Y;V}
    :=
    \operatorname{fib}
    \left(
        \map_{\Mod_B}(L_{B/R},-)|_{\Mod_B^+}
        \longrightarrow
        \mathsf T^{\mathrm{sp},+}_{Y/V,u}
    \right).
\]
It is an exact functor \(\Mod_B^+\to\Sp\).
\end{definition}

\begin{remark}[Domain of the tangent functor]
\label{rem:tangent-functor-domain}
For a general reconstruction-local target, the construction above produces an
exact tangent functor on the eventually connective category
\[
    \Mod_B^+.
\]
An extension to all of \(\Mod_B\) requires an additional continuity theorem,
such as a right-left or pro-coherent continuity result. Represented targets
already carry the full exact functor
\[
    M\longmapsto\map_{\Mod_B}(L_{B/R},M)
\]
on \(\Mod_B\).
\end{remark}



The reconstruction theory has now returned to stable coefficients. At a represented
centre the tangent functor is represented by the cotangent complex; for a general
reconstruction-local target it is an exact functor on eventually connective
coefficients. With both nonlinear reconstruction and tangent recovery in hand, we can
finally compare them with the differential-linear branch and with the stable boundary of
reconstruction localization.

\section{Comparison with the differentiated spectral doctrine}
\label{sec:formal-reconstruction-comparisons}

All objects needed for the deferred comparison have now been constructed. The following
diagram can therefore be read from top to bottom without introducing future
dependencies:
\[
\begin{tikzcd}[column sep=huge,row sep=large]
    & \text{spectral doctrine}
        \arrow[d,"\text{cotangent linearization}"]
        &
    \\
    & \text{cotangent classifier}
        \arrow[dl,"\text{tangent linear fibre}"']
        \arrow[dr,"\text{square-zero classification}"]
        \arrow[ddd,"\text{represented tangent recovery}" description]
        &
    \\
    !_X,\ \mathsf d
        &
        &
    \text{finite infinitesimal geometry}
        \arrow[d,"\text{completion and excision}"]
    \\
        &
        &
    \text{formal reconstruction}
        \arrow[dl,"\text{stabilization}"']
    \\
    & \text{recovered tangent functor}
        &
\end{tikzcd}
\]
The left branch is the differential-linear structure internal to the stable coefficient
fibre. The right branch is the nonlinear infinitesimal geometry used for reconstruction.
They are not identified; their common interface is the cotangent classifier. The
represented finite-projective case is where that interface becomes an explicit universal
derivation.

There are three comparisons. First, the differential exponential and the square-zero
reconstruction comonad act on different categories. Second, stabilization of the
centre-free localization produces only the unit-boundary sector of the four-adjoint
fracture formalism. Third, on represented finite-projective models, the
differential-linear derivative, infinitesimal first-order variation, and recovered
tangent functor meet at the universal cotangent derivation.

The cotangent-mediated comparison, and the absence of a direct comparison between the
two comonads, is summarized by
\[
\begin{tikzcd}[column sep=small,row sep=large]
    !_B\ \text{on }\Mod_B
        \arrow[dr,"\text{K\"ahler comparison}"']
    &&
    \mathbb C^{\mathrm{sqz}}_{B/R}\ \text{on anchored cotangent data}
        \arrow[dl,"\text{cotangent anchor}"]
    \\
    & \text{cotangent classifier }L
        \arrow[d,"\text{represented tangent recovery}"]
    &
    \\
    & \mathsf T^{\mathrm{sp},+}
    &
\end{tikzcd}
\]
The diagonal arrows are therefore mediated by the cotangent classifier; no natural
transformation between the two comonads is asserted.

\subsection{The differential exponential and the square-zero reconstruction comonad}

Chapter~\Ref{chap:differentiation-spectral-doctrine} constructed, in every tangent fibre
\(\Mod_B\simeq\QCoh(\Spec B)\), the cofree cocommutative-coalgebra exponential
\[
    !_B:\Mod_B\longrightarrow\Mod_B
\]
with its Seely comparisons, primitive infinitesimal coalgebras, codereliction, and
deriving transformation. The present chapter has constructed a different canonical
comonad,
\[
    \mathbb C^{\mathrm{sqz}}_{B/R}
    =
    \operatorname{cot}_{B/R}\operatorname{sqz}_{B/R},
\]
on the anchored shifted-cotangent category
\[
    \operatorname{AnchCot}^{\mathrm{aperf}}_{B/R,\geq0}.
\]

\begin{remark}[The two reconstruction-related comonads]
\label{rem:two-comonads-differentiated-doctrine}
The differential exponential
\[
    !_B:\Mod_B\longrightarrow\Mod_B
\]
governs differential-linear use of coefficients in the tangent fibre. The
comonad
\[
    \mathbb C^{\mathrm{sqz}}_{B/R}
    =
    \operatorname{cot}_{B/R}\operatorname{sqz}_{B/R}
\]
acts on the anchored shifted-cotangent category
\(\operatorname{AnchCot}^{\mathrm{aperf}}_{B/R,\geq0}\) and governs comonadic
reconstruction of complete augmented algebras. Their comparison in this
chapter is mediated by the cotangent complex and, on finite-projective models,
by the universal derivation described in
\Ref{cor:seely-formal-reconstruction-cotangent-interface}.
\end{remark}



The fixed-centre comparison above is nonlinear and comonadic. The centre-free
localization has a different stable shadow: after stabilization, its unit has an exact
fibre. We record that boundary next and compare only its formal shape with the BNV
fracture theorem.

\subsection{Stable reconstruction defect and the BNV comparison}
\label{subsec:stable-reconstruction-defect-bnv}

The left-exact reconstruction localization yields a canonical stable localization
sequence. This is intentionally weaker than the BNV-type fracture theorem of
Chapter~\Ref{chap:differentiation-spectral-doctrine}: a reflective adjunction supplies
one exact unit boundary after stabilization, whereas the remaining BNV sectors require
the four-adjoint formalism constructed there.

\begin{definition}[Stable reconstruction categories]
\label{def:stable-reconstruction-categories}
Put
\[
    \mathcal D_{\Sp/S}
    :=
    \operatorname{Sp}\bigl(((\mathcal H_{\Sp})_{/S})_*\bigr),
    \qquad
    \mathcal D^{\mathrm{Recon}}_{\Sp/S}
    :=
    \operatorname{Sp}\bigl((\bigl((\mathcal H_{\Sp})_{/S}\bigr)^{\mathrm{Recon}})_*\bigr).
\]
These are the stabilizations of the pointed ordinary and reconstruction-local
spectral Zariski slices.
\end{definition}

\begin{theorem}[Stabilized reconstruction localization]
\label{thm:stable-reconstruction-localization}
The adjunction
\[
    L^{\mathrm{Recon}}_S
    \dashv
    \iota_S
\]
stabilizes canonically to an exact reflective localization
\[
\begin{tikzcd}[column sep=huge]
    \mathcal D_{\Sp/S}
        \arrow[r,shift left=1.2ex,"L^{\mathrm{Recon},\mathrm{st}}_S"]
        &
    \mathcal D^{\mathrm{Recon}}_{\Sp/S}
        \arrow[l,shift left=1.2ex,"\iota^{\mathrm{st}}_S"]
\end{tikzcd}
\]
with \(\iota^{\mathrm{st}}_S\) fully faithful. Consequently
\[
    \operatorname{Recon}^{\mathrm{st}}_S
    :=
    \iota^{\mathrm{st}}_S
    L^{\mathrm{Recon},\mathrm{st}}_S
\]
is an exact idempotent localization endofunctor of \(\mathcal D_{\Sp/S}\).
\end{theorem}

\begin{proof}
\leavevmode
\begin{enumerate}[label=\textup{(\arabic*)}]
    \item \textbf{Use left exactness before stabilization.}
    By \Ref{thm:reconstruction-localization-left-exact},
    \(L^{\mathrm{Recon}}_S\) preserves finite limits as well as all colimits, and
    \(\iota_S\) is a fully faithful right adjoint preserving limits.
    \item \textbf{Pass the adjunction to pointed objects.}
    The induced
    adjunction on pointed objects therefore preserves the finite-limit structure
    used in stabilization.
    \item \textbf{Stabilize the adjunction.}
    Functoriality of spectrum objects for finite-limit
    preserving functors,
    \cite[Section~1.4.2]{LurieHigherAlgebra}, gives an adjunction of stable
    \(\infty\)-categories
    \[
        L^{\mathrm{Recon},\mathrm{st}}_S
        \dashv
        \iota^{\mathrm{st}}_S.
    \]
    \item \textbf{Prove exactness of the stabilized functors.}
    The left adjoint is exact because it is induced by a finite-limit preserving
    functor on pointed objects, and the right adjoint is exact because it preserves
    finite limits between stable categories.
    \item \textbf{Prove full faithfulness after stabilization.}
    Full faithfulness of
    \(\iota^{\mathrm{st}}_S\) follows from full faithfulness of \(\iota_S\).
    \item \textbf{Apply the localization identities.}
    The reflective-localization identities then make
    \(\operatorname{Recon}^{\mathrm{st}}_S\) exact and idempotent.
\end{enumerate}
\end{proof}

\begin{definition}[Stable reconstruction defect]
\label{def:stable-reconstruction-defect}
For \(E\in\mathcal D_{\Sp/S}\), define
\[
    \mathsf K^{\mathrm{Recon}}_S(E)
    :=
    \operatorname{fib}
    \left(
        E
        \longrightarrow
        \operatorname{Recon}^{\mathrm{st}}_S(E)
    \right).
\]
Let
\[
    \mathcal K^{\mathrm{Recon}}_S
    :=
    \ker\bigl(L^{\mathrm{Recon},\mathrm{st}}_S\bigr)
    \subseteq
    \mathcal D_{\Sp/S}.
\]
We call \(\mathsf K^{\mathrm{Recon}}_S(E)\) the
\textbf{stable reconstruction defect} of \(E\).
\end{definition}

\begin{proposition}[Stable reconstruction defect triangle]
\label{prop:stable-reconstruction-defect-triangle}
For every \(E\in\mathcal D_{\Sp/S}\) there is a canonical exact triangle
\[
    \mathsf K^{\mathrm{Recon}}_S(E)
    \longrightarrow
    E
    \longrightarrow
    \operatorname{Recon}^{\mathrm{st}}_S(E).
\]
Moreover:
\begin{enumerate}[label=(\roman*)]
    \item
    \[
        L^{\mathrm{Recon},\mathrm{st}}_S
        \mathsf K^{\mathrm{Recon}}_S(E)
        \simeq0;
    \]
    \item the functor
    \[
        \mathsf K^{\mathrm{Recon}}_S:
        \mathcal D_{\Sp/S}\longrightarrow\mathcal K^{\mathrm{Recon}}_S
    \]
    is right adjoint to the fully faithful inclusion
    \[
        \mathcal K^{\mathrm{Recon}}_S
        \hookrightarrow
        \mathcal D_{\Sp/S};
    \]
    \item there is therefore a canonical stable localization sequence
    \[
        \mathcal K^{\mathrm{Recon}}_S
        \hookrightarrow
        \mathcal D_{\Sp/S}
        \xrightarrow{\ L^{\mathrm{Recon},\mathrm{st}}_S\ }
        \mathcal D^{\mathrm{Recon}}_{\Sp/S}.
    \]
\end{enumerate}
The fibre construction is functorial in \(E\).
\end{proposition}

\begin{proof}
\leavevmode
\begin{enumerate}[label=\textup{(\arabic*)}]
    \item \textbf{Kernel property.}
    The exact triangle is the defining fibre sequence of
    \Ref{def:stable-reconstruction-defect}. Apply the exact localization functor
    \(L^{\mathrm{Recon},\mathrm{st}}_S\). Since the localization unit becomes an
    equivalence after localization,
    \[
        L^{\mathrm{Recon},\mathrm{st}}_S(E)
        \xrightarrow{\simeq}
        L^{\mathrm{Recon},\mathrm{st}}_S
        \operatorname{Recon}^{\mathrm{st}}_S(E),
    \]
    its fibre vanishes. This proves \textup{(i)}.
    \item \textbf{Right-adjoint property.}
    Let \(K\in\mathcal K^{\mathrm{Recon}}_S\). By the stabilized adjunction,
    \[
    \begin{aligned}
        \map_{\mathcal D_{\Sp/S}}
        \bigl(
            K,
            \operatorname{Recon}^{\mathrm{st}}_S(E)
        \bigr)
        &\simeq
        \map_{\mathcal D^{\mathrm{Recon}}_{\Sp/S}}
        \bigl(
            L^{\mathrm{Recon},\mathrm{st}}_S K,
            L^{\mathrm{Recon},\mathrm{st}}_S E
        \bigr)
        \\
        &\simeq0.
    \end{aligned}
    \]
    Applying \(\map_{\mathcal D_{\Sp/S}}(K,-)\) to the defining exact triangle therefore
    gives
    \[
        \map_{\mathcal D_{\Sp/S}}
        \bigl(K,\mathsf K^{\mathrm{Recon}}_S(E)\bigr)
        \xrightarrow{\simeq}
        \map_{\mathcal D_{\Sp/S}}(K,E).
    \]
    Together with \textup{(i)}, this is precisely the right-adjoint universal
    property in \textup{(ii)}. The localization sequence in \textup{(iii)} is the
    corresponding exact reflective localization.
\end{enumerate}
\end{proof}

\begin{remark}[Comparison with the stable exact-adjoint fracture]
\label{rem:stable-reconstruction-defect-bnv}
The fibre
\[
    \mathsf K^{\mathrm{Recon}}_S(E)
    =
    \operatorname{fib}
    \bigl(
        E\to
        \iota^{\mathrm{st}}_SL^{\mathrm{Recon},\mathrm{st}}_SE
    \bigr)
\]
has the unit-boundary form
\[
    \mathsf A(E)=\operatorname{fib}(E\to jLE)
\]
from the exact-adjoint fracture theorem of
Chapter~\Ref{chap:differentiation-spectral-doctrine}. The reconstruction
localization supplies the reflective adjunction
\[
    L^{\mathrm{Recon},\mathrm{st}}_S
    \dashv
    \iota^{\mathrm{st}}_S,
\]
so the comparison identifies the unit-boundary sector of that fracture
pattern. Recovering the secondary and cycle sectors and the characteristic
morphism of the four-adjoint BNV formalism requires the additional right
adjoint constructed in the preceding chapter.
\end{remark}



The final comparison is the represented finite-projective case, where the
differential-linear, infinitesimal, and reconstructed tangent constructions all admit
the same explicit cotangent representative.

\subsection{Common cotangent interface with differential-linear differentiation}

\begin{corollary}[Differential-linear differentiation and formal reconstruction meet at the cotangent classifier]
\label{cor:seely-formal-reconstruction-cotangent-interface}
Let \(B\) be connective, let \(P\) be a finite projective \(B\)-module, and put
\[
    A:=\operatorname{Sym}_B(P^\vee).
\]
Then
\[
    L_{A/B}\simeq A\otimes_BP^\vee,
\]
and the represented tangent functor of the affine context
\[
    \Spec(A)\longrightarrow\Spec(B)
\]
is
\[
    \mathsf T^{\mathrm{sp},+}_{\Spec(A)/\Spec(B)}(N)
    \simeq
    \map_{\Mod_A}(L_{A/B},N)
\]
for every \(N\in\Mod_A^+\).

Under the differential-linear construction of
Chapter~\Ref{chap:differentiation-spectral-doctrine}, the deriving
transformation
\[
    \mathsf d_P:
    !_BP\otimes_BP
    \longrightarrow
    !_BP
\]
induces the convolution derivation
\[
    D:
    \mathcal O_!(P)
    \longrightarrow
    \mathcal O_!(P)\otimes_BP^\vee,
    \qquad
    \mathcal O_!(P)
    :=
    \underline{\operatorname{Hom}}_B(!_BP,B).
\]
The comonad counit determines the canonical algebra morphism
\[
    j_P:
    A=\operatorname{Sym}_B(P^\vee)
    \longrightarrow
    \mathcal O_!(P),
\]
and the differential-linear comparison theorem gives a canonical homotopy
\[
    D\circ j_P
    \simeq
    (j_P\otimes1)\circ d_{\mathrm{univ}},
\]

equivalently, the square
\[
\begin{tikzcd}[column sep=huge,row sep=large]
    A \arrow[r,"j_P"] \arrow[d,"d_{\mathrm{univ}}"']
        & \mathcal O_!(P) \arrow[d,"D"]
    \\
    A\otimes_BP^\vee
        \arrow[r,"j_P\otimes1"']
        & \mathcal O_!(P)\otimes_BP^\vee
\end{tikzcd}
\]
commutes up to the canonical homotopy. Here
\[
    d_{\mathrm{univ}}:
    A\longrightarrow A\otimes_BP^\vee
    \simeq L_{A/B}
\]
is the spectral universal derivation.

The canonical first infinitesimal neighbourhood of the diagonal constructed in
Chapter~\Ref{chap:differentiation-spectral-doctrine} has first-order
difference given by the same universal relative derivation. Hence, on this
common represented finite-projective domain, differential-linear
differentiation, differential-cohesive first-order variation, and tangent
recovery from formal reconstruction meet at the cotangent classifier
\(L_{A/B}\), through the canonical comparisons displayed above.
\end{corollary}

\begin{proof}
\leavevmode
\begin{enumerate}[label=\textup{(\arabic*)}]
    \item \textbf{Identify the cotangent complex.}
    The cotangent calculation
    \[
        L_{A/B}\simeq A\otimes_BP^\vee
    \]
    is \Ref{prop:cotangent-free-algebra}.
    \item \textbf{Identify the represented tangent functor.}
    The represented tangent identification is \Ref{prop:represented-tangent-cotangent}.
    \item \textbf{Convolution derivative.}
    The cofree function algebra and the canonical map
    \[
        j_P:
        \operatorname{Sym}_B(P^\vee)\longrightarrow\mathcal O_!(P)
    \]
    are \Ref{def:cofree-function-algebra}. The differential-linear deriving
    transformation induces the convolution derivation \(D\) by
    \Ref{con:spectral-convolution-derivation}, and
    \Ref{thm:dill-kahler-comparison} gives the canonical commuting comparison
    \[
        D\circ j_P
        \simeq
        (j_P\otimes1)\circ d_{\mathrm{univ}}.
    \]
    Thus the differential-linear derivative is compatible with the universal
    spectral K\"ahler derivative after passage along the canonical comparison
    \(j_P\); the comparison is therefore mediated by \(j_P\).
    \item \textbf{First infinitesimal neighbourhood.}
    Finally,
    \Ref{thm:first-neighbourhood-universal-derivation} identifies the
    first-order difference of the two retractions of the canonical first
    infinitesimal neighbourhood with the same universal relative derivation.
    Together with cotangent representation of the recovered tangent functor,
    these comparisons establish the claimed common cotangent interface.
\end{enumerate}
\end{proof}



The comparison is now precise. Formal reconstruction neither identifies the
differential-linear and nonlinear comonads nor reproduces the full BNV fracture.
Instead, the cotangent classifier is the shared interface: it generates formal order,
classifies square-zero variation, represents the recovered tangent functor on
represented targets, and mediates the finite-projective K\"ahler comparison.

\section{The formal reconstruction interface}
\label{sec:formal-reconstruction-interface}

The chapter exports five packages rather than one undifferentiated list of
constructions: approximation and completion; fixed-centre comonadic reconstruction;
Artinian and centre-free reconstruction semantics; centred tangent recovery; and
comparison with the differentiated doctrine. The theorem below records the individual
results that make up those packages so that later chapters can cite them precisely.

\begin{theorem}[Formal reconstruction of the spectral doctrine]
\label{thm:formal-reconstruction-interface}
The constructions of this chapter provide the following.
\begin{enumerate}[label=(\roman*)]
    \item \textbf{Postnikov reconstruction.}
    Every connective commutative ring spectrum is recovered as
    \[
        B\simeq\varprojlim\nolimits_n\tau_{\leq n}^{\CAlg}B,
    \]
    and each positive Postnikov transition is a structured square-zero
    extension. The finite composite
    \[
        \tau_{\leq n}^{\CAlg}B\longrightarrow H\pi_0(B)
    \]
    is therefore a finite composite of structured square-zero extensions; its
    corresponding affine morphism is a finite spectral infinitesimal substitution.

    \item \textbf{Canonical adic filtration.}
    Filtration-zero augmentation admits a left adjoint
    \[
        \operatorname{adic}\dashv\operatorname{aug},
    \]
    and filtered completion is a symmetric-monoidal localization. Together
    they assign a functorial complete filtered algebra to every augmentation.

    \item \textbf{Cotangent description of the associated graded.}
    For every augmentation \(A\to B\),
    \[
        \operatorname{Gr}
        \widehat{\operatorname{adic}}(A\to B)
        \simeq
        \operatorname{Sym}^{\mathrm{gr}}_B
        \bigl(L_{B/A}[-1]\langle1\rangle\bigr).
    \]
    For connective maps surjective on \(\pi_0\), filtration zero of the
    completion is the Amitsur totalization. Under almost finite
    presentation it is derived completion for
    \(I=\ker(\pi_0(A)\to\pi_0(B))\).

    \item \textbf{Complete almost finitely presented augmentations.}
    On the domain of connective almost finitely presented augmentations that are surjective on \(\pi_0\), the fixed
    points of canonical adic completion are exactly the derived
    \(I\)-complete augmentations.

    \item \textbf{Comonadic shifted-cotangent--structured-square-zero reconstruction.}
    For a connective centre \(B\) over \(R\), the adjunction
    \[
    \begin{tikzcd}[column sep=huge]
        \CAlg(\Sp)^{\mathrm{cn},\mathrm{afp},\wedge}_{R//B}
            \arrow[r,shift left=1.2ex,"\operatorname{cot}_{B/R}"]
            &
        \operatorname{AnchCot}^{\mathrm{aperf}}_{B/R,\geq0}
            \arrow[l,shift left=1.2ex,"\operatorname{sqz}_{B/R}"]
    \end{tikzcd}
    \]
    is comonadic. Hence
    \[
        \CAlg(\Sp)^{\mathrm{cn},\mathrm{afp},\wedge}_{R//B}
        \simeq
        \operatorname{Coalg}_{\mathbb C^{\mathrm{sqz}}_{B/R}}
        \bigl(\operatorname{AnchCot}^{\mathrm{aperf}}_{B/R,\geq0}\bigr).
    \]
    The associated cobar resolution reconstructs each complete
    almost finitely presented augmentation from its comonadic shifted-cotangent data.

    \item \textbf{Artinian infinitesimal tests.}
    At a centre \(B\) over \(R\), with \(B\) coherent, spectral Artinian extensions are generated
    by finitely many structured square-zero stages with coefficients in connective coherent \(B\)-modules. Affine realization identifies them with the
    bounded coherent fixed-centre subcategory of the finite infinitesimal
    geometry of Chapter~\Ref{chap:differentiation-spectral-doctrine}.

    \item \textbf{Schlessinger formal-moduli problems.}
    The explicit localization of
    \[
        \Fun(\operatorname{Art}^{\mathrm{sp}}_{B/R},\mathcal S)
    \]
    at the centre and square-zero excision maps exists and has as local
    objects the spectral formal-moduli problems
    \(\operatorname{FMP}^{\mathrm{sp}}_{B/R}\).

    \item \textbf{Reconstruction localization over a base.}
    For every spectral Zariski stack \(S\), the localization
    \(L^{\mathrm{Recon}}_S\) generated by Postnikov continuity and nilpotent
    excision is left exact. Its essential image
    \(\bigl((\mathcal H_{\Sp})_{/S}\bigr)^{\mathrm{Recon}}\) is therefore an \(\infty\)-topos, and
    \[
        \operatorname{Recon}_S
        =
        \iota_S L^{\mathrm{Recon}}_S
    \]
    is an idempotent left-exact localization endofunctor of
    \((\mathcal H_{\Sp})_{/S}\).

    \item \textbf{Stable localization boundary.}
    Stabilization gives an exact reflective localization
    \[
    \begin{tikzcd}[column sep=huge]
        \mathcal D_{\Sp/S}
            \arrow[r,shift left=1.2ex,"L^{\mathrm{Recon},\mathrm{st}}_S"]
            &
        \mathcal D^{\mathrm{Recon}}_{\Sp/S}
            \arrow[l,shift left=1.2ex,"\iota^{\mathrm{st}}_S"]
    \end{tikzcd}
    \]
    and a functorial exact triangle
    \[
        \mathsf K^{\mathrm{Recon}}_S(E)
        \longrightarrow
        E
        \longrightarrow
        \operatorname{Recon}^{\mathrm{st}}_S(E).
    \]
    The first term belongs to
    \(\ker(L^{\mathrm{Recon},\mathrm{st}}_S)\) and is right adjoint to the
    inclusion of that kernel.

    \item \textbf{Centred Artinian leaves.}
    At an affine centre with \(B\) coherent,
    \[
        U=\Spec(B)\longrightarrow V=\Spec(R)\longrightarrow S,
    \]
    and a chosen \(U\)-point \(u:U\to Y\) of a reconstruction-local object
    \(Y\to S\), the centred Artinian leaf is a spectral formal-moduli problem. Objectwise it is the centred fibre of the
    endpoint restriction associated with the corresponding finite
    infinitesimal arrow.

    \item \textbf{Eventually connective tangent functor.}
    Lifting along split square-zero augmentations on connective \(B\)-modules forms an
    \(\Omega\)-spectrum and is right exact. Its Spanier--Whitehead extension
    is an exact functor
    \[
        \mathsf T^{\mathrm{sp},+}_{Y/V,u}:
        \Mod_B^+\longrightarrow\Sp.
    \]
    For represented targets,
    \[
        \mathsf T^{\mathrm{sp},+}_{\Spec(B)/\Spec(R)}(M)
        \simeq
        \map_{\Mod_B}(L_{B/R},M).
    \]

    \item \textbf{Comparison with differential-linear differentiation.}
    For \(P\) finite projective and
    \(A=\operatorname{Sym}_B(P^\vee)\), the represented tangent functor, the
    first infinitesimal neighbourhood of the diagonal, and the
    differential-linear convolution derivative all meet at the universal
    derivation
    \[
        d_{\mathrm{univ}}:
        A\longrightarrow L_{A/B}
        \simeq A\otimes_BP^\vee.
    \]
    The comparison with the convolution derivative is mediated by the
    canonical map
    \[
        j_P:
        \operatorname{Sym}_B(P^\vee)\longrightarrow\mathcal O_!(P).
    \]
\end{enumerate}
\end{theorem}

\begin{proof}
\leavevmode
\begin{enumerate}[label=\textup{(\arabic*)}]
    \item \textbf{Cite Postnikov reconstruction.}
    Item \textup{(i)} is
    \Ref{thm:postnikov-square-zero-stages} together with
    \Ref{cor:infinitesimal-postnikov-reconstruction}.
    \item \textbf{Cite filtered reconstruction.}
    Items
    \textup{(ii)}--\textup{(iv)} follow from
    \Ref{prop:filtered-completion-reflector},
    \Ref{thm:adic-associated-graded},
    \Ref{thm:adic-adams-completion},
    \Ref{thm:adic-adams-derived-completion}, and
    \Ref{prop:complete-aft-intrinsic-characterization}.
    \item \textbf{Cite comonadic reconstruction.}
    Item \textup{(v)} is
    \Ref{thm:cotangent-square-zero-comonadicity} and its two corollaries
    \Ref{cor:cotangent-square-zero-coalgebra-classification} and
    \Ref{cor:canonical-square-zero-cobar-resolution}.
    \item \textbf{Cite Artinian reconstruction.}
    Item \textup{(vi)} follows
    from \Ref{prop:artinian-finite-square-zero-generation} and
    \Ref{prop:artinian-as-finite-infinitesimal-tests}.
    \item \textbf{Cite Schlessinger reconstruction.}
    Item \textup{(vii)} is
    \Ref{thm:artinian-schlessinger-reflector} together with
    \Ref{prop:formal-moduli-finite-excision}.
    \item \textbf{Cite reconstruction localization.}
    Item \textup{(viii)} follows from
    \Ref{thm:global-spectral-reconstruction-reflector},
    \Ref{thm:reconstruction-localization-left-exact}, and
    \Ref{prop:reconstruction-localization-idempotence}.
    \item \textbf{Cite the stable boundary.}
    Item \textup{(ix)} is
    \Ref{thm:stable-reconstruction-localization} and
    \Ref{prop:stable-reconstruction-defect-triangle}.
    \item \textbf{Cite the centred leaf comparison.}
    Item \textup{(x)} is
    \Ref{prop:artinian-leaf-endpoint-fibre} and
    \Ref{prop:reconstruction-local-artinian-leaf-fmp}.
    \item \textbf{Cite tangent recovery.}
    Item \textup{(xi)} is
    \Ref{thm:reconstruction-tangent-spectrum} together with
    \Ref{prop:represented-tangent-cotangent}.
    \item \textbf{Cite the differential-linear comparison.}
    Item \textup{(xii)} is
    \Ref{cor:seely-formal-reconstruction-cotangent-interface}.
\end{enumerate}
\end{proof}

\section{Chapter summary}
\label{sec:formal-reconstruction-summary}

Two approximation coordinates now organize reconstruction. Postnikov height resolves a
connective spectral algebra into canonical finite square-zero stages over its discrete
centre and recovers it by inverse limit. Formal order resolves an augmentation through
its canonical adic filtration, whose complete associated graded is the free graded
commutative algebra on the shifted relative cotangent complex. These coordinates are
different, but the cotangent classifier links formal order to the infinitesimal geometry
of the preceding chapter.

At a fixed centre, that link becomes the principal reconstruction theorem. Complete
almost finitely presented augmentations are comonadically recovered from the anchored
shifted cotangent object
\[
    L_{B/R}[-1]\longrightarrow L_{B/A}[-1]
\]
with coalgebra coherence generated by structured square-zero formation. Its bounded
coherent fragment is the Artinian test category and yields Schlessinger formal-moduli
semantics. Removing the centre leaves Postnikov continuity and nilpotent excision as the
global reconstruction equations, whose left-exact localization defines the centre-free
semantics.

Returning to a centre closes the loop. Split square-zero lifting stabilizes to an exact
tangent functor on eventually connective coefficients, and represented targets recover
the cotangent complex as its representing object. The differential-linear comonad, the
square-zero reconstruction comonad, and the stable localization boundary remain distinct
constructions; what unifies them is the cotangent classifier, which measures first-order
variation, generates formal order, classifies square-zero reconstruction, and represents
tangent behaviour.
